\section{Half-line smoothing, characteristic functions, and resource enclosures}
\label{ms:fourier:appendix}

\subsection{Exact half-line smoothing}
\label{ms:fourier:kernel}

The main conserved identity reduces the problem to bounding explicit
Fourier integrals and moment expressions. This appendix develops the
half-line smoothing kernel, characteristic-function enclosures, complete
low- and high-frequency bounds, and the resource-to-confinement maps
used in those estimates. Section~\ref{ms:certificates} records their
certificate conditions and verification status.

Write $\phi$ and $\Phi$ for the standard normal density and distribution
function. All terminating decimals are exact rational numbers. The
imaginary unit is $i$. Symbols introduced within a subsection, including
$D$, $R$, $L$, $s$ and $K$, are local to that subsection unless explicitly
carried forward. In particular the moment defect $D$, an auxiliary distance
bound, and the Bernoulli inward displacement $d$ are different quantities.

The one-sided kernel is
\[
 K(t)=\frac{1-t}{2}+\frac{i}{2}
       \left((1-t)\cot(\pi t)+\frac1\pi\right),\qquad 0<t<1.
\]
For $\kappa\in\{K,\overline K\}$, $0<b\le1$ and
$g(t)=g(0)+O(t)$ with real $g(0)$, the convention throughout is
\begin{equation}\label{ms:fourier:kernel-convention}
 \mathfrak I_{\kappa,b}(g)
 =2\lim_{\epsilon\downarrow0}\operatorname{Re}
       \int_\epsilon^b\kappa(t)g(t)\,dt.
\end{equation}
Only the real-part improper limit is intended; the unmodified complex
integral need not converge. The proof below establishes directly that
\[
 M(x)=\tfrac12+\mathfrak I_{K,1}(e^{2\pi itx})
       \ge\mathbf1_{(-\infty,0]}(x),\qquad
 1-M(-x)\le\mathbf1_{(-\infty,0)}(x),
\]
including atoms at the boundary, and justifies the expectation identity
for an integrable random variable.

\subsubsection{A real-variable Beurling construction}

Set

\[
 W(x)=\left(\frac{\sin\pi x}{\pi x}\right)^2\quad(x\ne0),
 \qquad W(0)=1,
\]

and define the continuous odd function

\begin{equation}
 H_B(x)=\sum_{m=1}^{\infty}\{W(x-m)-W(x+m)\}+2xW(x).
 \label{ms:fourier:kernel-1}
\end{equation}

The series converges absolutely and locally uniformly on the real line:
for \(|x|\le R\) and \(m>2R+1\), each shifted term is bounded by a
constant times \(m^{-2}\). Thus the function and all the integer values in
what follows are unambiguous. In particular,

\begin{equation}
 H_B(0)=0,\qquad H_B(j)=\operatorname{sgn}(j)
 \quad(j\in\mathbb Z\setminus\{0\}).                 \label{ms:fourier:kernel-2}
\end{equation}

Here \(\operatorname{sgn}(0)=0\). The corresponding classical Beurling
majorant is \(B=H_B+W\). For noninteger \(x\), its series is

\begin{equation}
 B(x)=\left(\frac{\sin\pi x}{\pi}\right)^2
 \left\{\frac2x+\sum_{m=0}^{\infty}\frac1{(x-m)^2}
                    -\sum_{m=1}^{\infty}\frac1{(x+m)^2}\right\}.
 \label{ms:fourier:kernel-3}
\end{equation}

We prove the needed inequality directly, without importing an extremality
or Fourier-transform theorem about \(B\).

The elementary partial-fraction identity

\begin{equation}
 \sum_{m\in\mathbb Z}\frac1{(y-m)^2}
       =\frac{\pi^2}{\sin^2\pi y},\qquad y\notin\mathbb Z,
 \label{ms:fourier:kernel-4}
\end{equation}

has the following short proof. Integrate
\(\pi\cot(\pi z)/(z-y)^2\) counterclockwise over the square with vertices
\(\pm(N+\tfrac12)\pm i(N+\tfrac12)\). The cotangent is uniformly bounded
on these edges: on the vertical edges it is a multiple of a hyperbolic
tangent, and on the horizontal edges its modulus is bounded by
\(\coth(\pi(N+\tfrac12))\). The denominator is of order \(N^2\) and the
perimeter of order \(N\), so the integral tends to zero. Its residues are
\((m-y)^{-2}\) at \(m=-N,\ldots,N\), and
\(-\pi^2\csc^2(\pi y)\) at \(y\). The residue theorem and absolute
convergence give \eqref{ms:fourier:kernel-4}.

For \(y>0\), \(y\notin\mathbb Z\), put

\[
 A(y)=\left(\frac{\sin\pi y}{\pi}\right)^2,
 \qquad S(y)=\sum_{m=1}^{\infty}(y+m)^{-2}.
\]

Subtracting \eqref{ms:fourier:kernel-4} from \eqref{ms:fourier:kernel-1} gives exactly

\begin{equation}
 H_B(y)-1=A(y)\left\{\frac2y-\frac1{y^2}-2S(y)\right\}.
 \label{ms:fourier:kernel-5}
\end{equation}

The decreasing positive function \(s\mapsto(y+s)^{-2}\) satisfies

\begin{equation}
 S(y)\le\int_0^\infty(y+s)^{-2}\,ds=\frac1y
       \le\frac1{y^2}+S(y).
 \label{ms:fourier:kernel-6}
\end{equation}

Consequently the expression in braces in \eqref{ms:fourier:kernel-5} lies between
\(-y^{-2}\) and \(y^{-2}\). Hence

\[
 |H_B(y)-1|\le W(y)\quad(y>0).
\]

Positive integer points follow directly from \eqref{ms:fourier:kernel-2}. Oddness of \(H_B\) and
evenness of \(W\), together with their values at zero, now prove

\begin{equation}
 |H_B(x)-\operatorname{sgn}(x)|\le W(x)
 \qquad(x\in\mathbb R).                              \label{ms:fourier:kernel-7}
\end{equation}

This is the required sign-function majorization, with its normalization
proved rather than assumed.

\subsubsection{Recovering the exact Fourier kernel}

Elementary integration gives, for every real \(x\),

\begin{equation}
 W(x)=2\int_0^1(1-t)\cos(2\pi tx)\,dt.              \label{ms:fourier:kernel-8}
\end{equation}

For \(0<r<1\), insert \eqref{ms:fourier:kernel-8} into the absolutely convergent Abel-weighted
series. Termwise integration is legitimate because
\(\sum_{m\ge1}r^m<\infty\). The cosine subtraction formula gives

\begin{equation}
 \begin{aligned}
 \sum_{m\ge1}r^m\{W(x-m)-W(x+m)\}
   &=4\int_0^1(1-t)\sin(2\pi tx)
               \sum_{m\ge1}r^m\sin(2\pi mt)\,dt\\
   &=4\int_0^1(1-t)\sin(2\pi tx)
        \frac{r\sin(2\pi t)}{1-2r\cos(2\pi t)+r^2}\,dt.
 \end{aligned}                                                    \label{ms:fourier:kernel-9}
\end{equation}

For \(0<t<1\),

\begin{equation}
 \left|\frac{r\sin(2\pi t)}{1-2r\cos(2\pi t)+r^2}\right|
 \le\frac12|\cot(\pi t)|,
 \label{ms:fourier:kernel-10}
\end{equation}

since the denominator is
\((1-r)^2+4r\sin^2(\pi t)\). The right side of \eqref{ms:fourier:kernel-9}, including its
remaining factors, is therefore dominated by

\[
 2(1-t)|\sin(2\pi tx)|\,|\cot(\pi t)|.
\]

This is integrable at both endpoints: the sine is \(O_x(t)\) at zero,
and \((1-t)\cot(\pi t)\) is bounded at one. Dominated convergence on
the integral side and absolute convergence on the series side allow
\(r\uparrow1\). The fraction in \eqref{ms:fourier:kernel-9} tends to
\(\tfrac12\cot(\pi t)\). Also, direct integration gives

\begin{equation}
 2xW(x)=\frac2\pi\int_0^1\sin(2\pi tx)\,dt,          \label{ms:fourier:kernel-11}
\end{equation}

including \(x=0\). Combining \eqref{ms:fourier:kernel-1} and \eqref{ms:fourier:kernel-9}--\eqref{ms:fourier:kernel-11},

\begin{equation}
 H_B(x)=2\int_0^1 R(t)\sin(2\pi tx)\,dt,
 \qquad R(t)=(1-t)\cot(\pi t)+\frac1\pi.
 \label{ms:fourier:kernel-12}
\end{equation}

Every integral in \eqref{ms:fourier:kernel-12} is a convergent real integral. No principal-value
Fourier manipulation has been used.

Now take exactly the manuscript's kernel

\begin{equation}
 K(t)=\frac{1-t}{2}+\frac i2 R(t),\qquad0<t<1.
 \label{ms:fourier:kernel-13}
\end{equation}

Equations \eqref{ms:fourier:kernel-8} and \eqref{ms:fourier:kernel-12} establish the exact identity

\begin{equation}
 \boxed{\displaystyle
 M(x):=\frac12+2\lim_{\epsilon\downarrow0}\operatorname{Re}
       \int_\epsilon^1K(t)e^{2\pi itx}\,dt
       =\frac{1-H_B(x)+W(x)}2=\frac{1+B(-x)}2.}
 \label{ms:fourier:kernel-14}
\end{equation}

In particular, the exponent has the \emph{positive} sign and the imaginary
part of \(K\) has the \emph{positive} sign. The minus sign on \(H_B\) comes
from taking the real part of \(iR(t)e^{2\pi itx}\).

\subsubsection{Closed boundary, cardinal values and the minorant}

For \(x<0\), \eqref{ms:fourier:kernel-7} gives \(M(x)\ge1\); for \(x>0\), it gives
\(M(x)\ge0\). At the boundary, \eqref{ms:fourier:kernel-2} and \(W(0)=1\) give

\begin{equation}
 M(0)=1.                                             \label{ms:fourier:kernel-15}
\end{equation}

Thus this majorizes the \emph{closed} half-line, including a point mass at
zero, rather than merely the half-height convention for a jump:

\begin{equation}
 M(x)\ge1_{(-\infty,0]}(x),\qquad
 1-M(-x)\le1_{(-\infty,0)}(x).
 \label{ms:fourier:kernel-16}
\end{equation}

The second inequality follows by subtracting the first inequality at
\(-x\) from one; its endpoint is open. Direct cardinal values are

\begin{equation}
 M(j)=1_{\{j\le0\}},\qquad1-M(-j)=1_{\{j<0\}}
 \quad(j\in\mathbb Z).                              \label{ms:fourier:kernel-17}
\end{equation}

The proof also gives the useful global error bound

\begin{equation}
 0\le M(x)-1_{(-\infty,0]}(x)\le W(x)
 \quad(x\in\mathbb R).                              \label{ms:fourier:kernel-18}
\end{equation}

At zero the left-hand error is exactly zero; it is not obtained by
assigning half of an atom to each side. In particular \(0\le M\le2\).

\subsubsection{Expectations and the compensated identity}

Suppose \(U\) is real and \(\E|U|<\infty\), and use the convention
\(\psi_U(t)=\E e^{2\pi itU}\). Since
\(2\int_0^1\operatorname{Re}K(t)\,dt=1/2\), \eqref{ms:fourier:kernel-14} can be written

\begin{equation}
 M(u)=1+2\operatorname{Re}\int_0^1
             K(t)\{e^{2\pi itu}-1\}\,dt.             \label{ms:fourier:kernel-19}
\end{equation}

The complex integral in \eqref{ms:fourier:kernel-19} is absolutely convergent. Indeed,
\(K(t)=i/(2\pi t)+1/2+O(t)\) at zero, \(K\) is bounded near one,
and

\[
 \E|e^{2\pi itU}-1|\le2\pi t\,\E|U|.
\]

Consequently
\(\int_0^1|K(t)|\E|e^{2\pi itU}-1|\,dt<\infty\), so ordinary
absolute Fubini applies to \eqref{ms:fourier:kernel-19}. Restoring the constant real part gives

\begin{equation}
 \E M(U)=\frac12+
     2\lim_{\epsilon\downarrow0}\operatorname{Re}
        \int_\epsilon^1K(t)\psi_U(t)\,dt.            \label{ms:fourier:kernel-20}
\end{equation}

Since \(\psi_{-U}=\overline{\psi_U}\), conjugation similarly yields

\begin{equation}
 \E\{1-M(-U)\}=\frac12-
     2\lim_{\epsilon\downarrow0}\operatorname{Re}
        \int_\epsilon^1\overline{K(t)}\psi_U(t)\,dt.
 \label{ms:fourier:kernel-21}
\end{equation}

These are the two expectation identities under the real-part
improper convention \eqref{ms:fourier:kernel-convention}. All expectations are finite; \eqref{ms:fourier:kernel-18} even gives
boundedness of the majorant without a moment assumption. The first
absolute moment is a sufficient hypothesis for the stated compensated
Fourier/Fubini argument. The uncompensated complex integral itself need
not converge.

\subsubsection{Four-summand smoothing normalization}

Let \(Y\) be integrable, let \(h\in\mathbb R\), and let \(T>0\). For
\(\chi_Y(s)=\E e^{isY}\), apply \eqref{ms:fourier:kernel-16} and \eqref{ms:fourier:kernel-20} to
\(U=T(Y-h)/(2\pi)\). The resulting exact normalization is

\begin{equation}
 \Pp(Y\le h)\le\frac12+
  2\lim_{\epsilon\downarrow0}\operatorname{Re}
       \int_\epsilon^1K(t)e^{-iTth}\chi_Y(Tt)\,dt.
 \label{ms:fourier:kernel-22}
\end{equation}

For the four-summand application,
\(Z=(X_1+\cdots+X_4)/2\), \(\varphi_P(s)=\E e^{isX_1}\),
and iid summands, choose \(T=6\). Independence gives
\(\chi_Z(6t)=\varphi_P(3t)^4\). Thus \eqref{ms:fourier:kernel-22} is exactly

\begin{equation}
 F_Z(x)\le\frac12+\int_0^1\operatorname{Re}
   \{\widehat K(t)e^{-6itx}\varphi_P(3t)^4\}\,dt,
 \qquad \widehat K=2K,
 \label{ms:fourier:kernel-23}
\end{equation}

the physical four-summand smoothing normalization. Writing \(D=x-Z\) and
\(S(t)=tR(t)\), its real integrand is exactly

\begin{equation}
 (1-t)\E\cos(6Dt)+S(t)\E[6D\operatorname{sinc}(6Dt)],
 \qquad \operatorname{sinc}(v)=\sin(v)/v,
 \label{ms:fourier:kernel-24}
\end{equation}

with the continuous value at zero. This verifies both the factor six
and the sign of the sine term. No non-atomicity, finite-support, density,
or continuity-at-the-threshold assumption is needed. For \eqref{ms:fourier:kernel-23}, iid
independence and integrability are sufficient; the manuscript's bounded
support assumptions imply integrability.

\subsection{Characteristic functions and conserved resources}
\label{ms:fourier:resources}

The following arguments retain the actual marginal law and raw threshold.
Attainment and contact are inputs where explicitly stated; the numerical
support, moments and probability panels are acquired hypotheses. The
finite induction at the end justifies a declared schedule of valid
updates, not the acceptance of any particular schedule. The classical
four-term Prawitz smoothing inequality is explicitly cited at its use.
The preceding pointwise construction alone is not substituted for that
external inequality without deriving its four-term normalization.

\subsubsection{Elementary characteristic-function bounds}
\label{ms:fourier:afcore:sec:fourier}

Throughout this section, $X$ is real, $\E X=0$, $\E X^2=1$, and
$\rho=\E|X|^3<\infty$. Write
\[
 f(u)=\E e^{iuX},\qquad g(u)=e^{-u^2/2},\qquad \gamma=\E X^3.
\]
No support, extremality, or predecessor assumption is made here.

\paragraph{Real and imaginary coordinates}

\begin{lemma}[Anchored coordinates]\label{ms:fourier:afcore:lem:anchor}
For $0\le u\le\pi$,
\begin{align}
 \cos u\le\Re f(u)&\le\cos u+\frac{\rho-1}{6}u^3,
 \label{ms:fourier:afcore:eq:real-anchor}\\
 |\Im f(u)|^2&\le\frac{(\rho-1)(\rho+5/3)}{36}u^6.
 \label{ms:fourier:afcore:eq:imag-anchor}
\end{align}
\end{lemma}
\begin{proof}
Put $\operatorname{sinc}v=\int_0^1\cos(tv)\,dt$, including $v=0$.
The elementary bounds $|\sin v|\le|v|$ give
\[
 |\operatorname{sinc}a-\operatorname{sinc}b|
 \le\frac{|a-b|}{2},\qquad
 |\operatorname{sinc}a-\operatorname{sinc}b|
 \le\frac{|a^2-b^2|}{6}.
\]
For the second inequality, use
$|\cos\sqrt A-\cos\sqrt B|\le|A-B|/2$ for $A,B\ge0$,
and integrate in $t$. Evenness handles either sign of $a,b$.

For $r\ge0$, integration from $1$ to $r$ of the first sinc bound gives
\[
 \left|\cos(ur)-\cos u+
       \frac{u\sin u}{2}(r^2-1)\right|
 \le\frac{u^3}{12}(r-1)^2(2r+1).
\]
Indeed the derivative of the expression inside the absolute value is
$u^2r\{\operatorname{sinc}u-\operatorname{sinc}(ur)\}$, and
$\int_{\min(1,r)}^{\max(1,r)}t|t-1|\,dt
=(r-1)^2(2r+1)/6$.
For $R=|X|$, the weight $W(R)=(R-1)^2(2R+1)$ satisfies
$\E W(R)=2(\rho-1)$. Averaging proves the upper real bound, and in fact
the corresponding two-sided anchor error.

For the sharper lower bound, take $0<u\le\pi$ and
$a=\sin u/(2u)\ge0$. The function $v\mapsto\cos v+av^2$ has its
minimum on $[0,\pi]$ at $u$, since sinc decreases there. For $v\ge\pi$,
\[
 \cos v+av^2\ge-1+a\pi^2\ge\cos u+au^2.
\]
Thus this minimum is global on $[0,\infty)$. Substitute $v=uR$ and
average; $\E R^2=1$ cancels the quadratic term. The case $u=0$ is immediate.

Centering and the second sinc bound give
\[
 |\Im f(u)|
 \le\frac{u^3}{6}\E[R|R^2-1|].
\]
Cauchy--Schwarz, with $W$ as above, yields
\[
 \{\E[R|R^2-1|]\}^2
 \le\E W(R)\,
       \E\frac{R^2(R+1)^2}{2R+1}.
\]
There is no singularity at $R=1$: the factorization used here is
$R|R^2-1|=\sqrt{W(R)}R(R+1)/\sqrt{2R+1}$.
The identity
\[
 r^3+\frac{29r^2+1}{18}
 -\frac{2r^2(r+1)^2}{2r+1}
 =\frac{(r-1)^2(4r+1)}{18(2r+1)}\ge0
\]
shows that the second expectation is at most $(\rho+5/3)/2$.
This proves \eqref{ms:fourier:afcore:eq:imag-anchor}. Every expectation used is
finite under the third-moment assumption.
\end{proof}

\begin{lemma}[The sine circle]\label{ms:fourier:afcore:lem:sine-circle}
Let
\[
 r=\E X_+=\E X_->0,\quad A=\E X_+^2,\quad B=\E X_-^2,
 \qquad \zeta=\frac{2AB}{r}.
\]
Then $A+B=1$ and
\begin{equation}\label{ms:fourier:afcore:eq:sine-circle}
 \zeta^2+\frac{36|\Im f(u)|^2}{u^6}\le\rho^2
 \qquad(u\ne0),\qquad
 \zeta\ge b(\rho):=\frac{\sqrt{\rho^2+8}-\rho}{2}.
\end{equation}
In particular, for $0\le u\le\pi$ the coefficient in the imaginary
bound may be replaced by
\[
 |\Im f(u)|^2\le\frac{u^6}{36}
 \min\{(\rho-1)(\rho+5/3),\,\rho^2-b(\rho)^2\}.
\]
Both bounds also hold for $\gamma^2$ after deleting $u^6/36$.
\end{lemma}
\begin{proof}
Let $\mu_+,\mu_-$ be the restrictions of the law of $|X|$ to the
strictly positive and strictly negative classes, respectively. For
$a,b>0$, let $P_{a,b}$ put masses $b/(a+b),a/(a+b)$ at $a,-b$.
The measure
\[
 d\nu(a,b)=\frac{a+b}{r}\,d\mu_+(a)d\mu_-(b)
\]
has mass $\Pp(X\ne0)$. Mixing $P_{a,b}$ against $\nu$, and retaining
the original atom at zero, recovers the law of $X$.
For a component, write $f_{a,b}(u)=\E_{P_{a,b}}e^{iuX}$ and
\[
 e_{a,b}=\frac{2a^2b^2}{a+b},\qquad
 \rho_{a,b}=\frac{ab(a^2+b^2)}{a+b}.
\]
The squared-argument sinc bound gives
\[
 \left|\frac{6\Im f_{a,b}(u)}{u^3}\right|
 \le\frac{ab|a^2-b^2|}{a+b}.
\]
Consequently the Euclidean norm of
$(e_{a,b},6\Im f_{a,b}(u)/u^3)$ is at most $\rho_{a,b}$.
Integrate this vector inequality. The two integrals on the right and
in the first coordinate are $\rho$ and $2AB/r$, respectively.
The triangle inequality proves the circle. This argument also covers
an atom at zero, whose vector contribution is zero.

Signwise Cauchy inequalities give
\[
 A^2\le r\E X_+^3,\qquad B^2\le r\E X_-^3,\qquad
 \frac{r^2}{A}+\frac{r^2}{B}\le\Pp(X\ne0)\le1.
\]
Hence $\rho+\zeta\ge1/r$, $AB\ge r^2$ and $\zeta\ge2r$.
All quantities are positive, so
$\zeta(\rho+\zeta)\ge2$. Solving this quadratic gives the stated
lower bound. Finally
$6\Im f(u)/u^3\longrightarrow-\gamma$ as $u\to0$ by dominated
convergence, using $|\sin v-v|\le|v|^3/6$. This also proves the assertions
about $\gamma$.
\end{proof}

On a closed moment band $\rho\in[\rho_-,\rho_+]$, any rational $Z>0$
satisfying $Z^2+\rho_+Z\le2$ is a valid uniform lower bound for
$b(\rho)$. Thus no numerical square root is needed to certify this
part of the imaginary enclosure.

\paragraph{An explicit scalar certificate and its convex profile}

\begin{lemma}[Scalar cosine bound]\label{ms:fourier:afcore:lem:cosine}
Assume the finite Bernstein positivity condition specified in the proof
below. Then, for every real $x$,
\begin{equation}\label{ms:fourier:afcore:eq:cosine}
 \cos x\le1-\frac{x^2}{2}+\kappa|x|^3,
 \qquad \kappa=\frac{62}{625}.
\end{equation}
\end{lemma}
\begin{proof}
By evenness take $x\ge0$. For $x\ge T:=625/124$ the right side is
at least one. Let $P_{16}(x)=\sum_{j=0}^8(-1)^jx^{2j}/(2j)!$.
The sixteenth derivative of $\cos x-P_{16}(x)$ is $\cos x-1\le0$,
and its first sixteen initial data vanish. Therefore
$\cos x\le P_{16}(x)$ on the whole nonnegative half-line.
It remains to prove positivity on $[0,T]$ of
\[
 Q(x)=\frac{1-x^2/2+\kappa x^3-P_{16}(x)}{x^3}
      =\sum_{j=0}^{13}q_jx^j,
\]
where the removable value at zero is used and the exact coefficient
vector is
\begin{align*}
 (q_0,\ldots,q_6)&=
 \left(\frac{62}{625},-\frac1{24},0,\frac1{720},0,
 -\frac1{40320},0\right),\\
 (q_7,\ldots,q_{13})&=
 \left(\frac1{3628800},0,-\frac1{479001600},0,
 \frac1{87178291200},0,-\frac1{20922789888000}\right).
\end{align*}
Here is a fully specified finite rational certificate. For $0\le r<80$
set $(a_r,b_r)=(r/16,(r+1)/16)$, and set $(a_{80},b_{80})=(5,625/124)$.
For $0\le j,k\le13$ define
\begin{align}
 c_{r,j}&=\sum_{\ell=j}^{13}q_\ell\binom\ell j
               a_r^{\ell-j}(b_r-a_r)^j,\nonumber\\
 B_{r,k}&=\sum_{j=0}^k c_{r,j}\frac{\binom kj}{\binom{13}j}.
 \label{ms:fourier:afcore:eq:bernstein-data}
\end{align}
All $81\cdot14$ numbers are thereby specified without external data.
The finite arithmetic condition required here is
\begin{equation}\label{ms:fourier:afcore:eq:bernstein-min}
 \min_{0\le r\le80,\ 0\le k\le13}B_{r,k}
 =B_{63,12}=\frac{2479189}{65585520000}>0.
\end{equation}
The identity
\[
 Q(a_r+(b_r-a_r)t)
 =\sum_{k=0}^{13}B_{r,k}\binom{13}k t^k(1-t)^{13-k}
\]
follows by expanding either side. Its basis is nonnegative and sums
to one for $0\le t\le1$. Equations
\eqref{ms:fourier:afcore:eq:bernstein-data}--\eqref{ms:fourier:afcore:eq:bernstein-min}
therefore imply positivity on every entire closed cell when that finite
condition holds. The cells cover $[0,T]$, completing the conditional proof.
\end{proof}

\begin{lemma}[Compatible profile]\label{ms:fourier:afcore:lem:profile}
Put $c_{\cos}=1-\cos4$, $a=129/1250$ and $A=16a/c_{\cos}$. The function
\begin{equation}\label{ms:fourier:afcore:eq:profile}
 q(v)=\begin{cases}
  1/2-\kappa v,&0\le v\le4,\\
  A(1-\cos v)/v^2,&4\le v\le2\pi,\\
  0,&v\ge2\pi
 \end{cases}
\end{equation}
is continuous, nonnegative, convex and nonincreasing on $[0,\infty)$,
and $\cos v\le1-v^2q(v)$ there.
\end{lemma}
\begin{proof}
The scalar bound at four gives $c_{\cos}\ge16a>0$, hence $0<A\le1$.
The pieces agree at both endpoints. On the curved piece,
\[
 q'(v)=A\frac{v\sin v-2(1-\cos v)}{v^3}\le0,
 \qquad
 q''(v)=\frac{A}{v^4}
       \{6+(v^2-6)\cos v-4v\sin v\}.
\]
Write the brace as $B(v)$. On $[4,3\pi/2]$,
$B'(v)=-2v\cos v+(2-v^2)\sin v\ge0$.
The elementary bounds $25/8<\pi<22/7$ imply
$7\pi/6<4<3\pi/2$, so $B(4)\ge6-10+8=4$.
On $[3\pi/2,2\pi]$ each remaining term in $B$ is nonnegative,
and $B\ge6$. Thus the curved branch is convex, with derivative zero
at $2\pi$. The upward derivative join at four is equivalent to
\[
 258\sin4+119(1-\cos4)\ge0.
\]
Let $S_{15}(x)=\sum_{j=0}^7(-1)^jx^{2j+1}/(2j+1)!$.
The fifteenth derivative of $\sin x-S_{15}(x)$ is $1-\cos x\ge0$,
with vanishing initial data; hence $\sin x\ge S_{15}(x)$ for $x\ge0$.
The exact rational lower bound at the join is
\[
 258S_{15}(4)+119\{1-P_{16}(4)\}
 =\frac{967616288}{638512875}>1.
\]
This proves the required join. The profile majorizes the cosine defect
by the scalar inequality on $[0,4]$, by $A\le1$ on $[4,2\pi]$,
and by $\cos v\le1$ thereafter.

The bounds $25/8<\pi<22/7$ follow from the rational comparisons in
Lemma~\ref{ms:gaussian-constants}.
\end{proof}

\begin{lemma}[Whole-law modulus]\label{ms:fourier:afcore:lem:modulus}
Let $X'$ be an independent copy of $X$, and let
$D\ge\E|X-X'|^3/2$. For $u\ge0$,
\begin{equation}\label{ms:fourier:afcore:eq:modulus}
 |f(u)|^2\le1-2u^2q(uD).
\end{equation}
The universal choice $D=\rho+1$ is valid. In particular the two
nontrivial branches are
\[
 |f(u)|^2\le
 \begin{cases}
  1-u^2+2\kappa Du^3,&uD\le4,\\
  1-2A(1-\cos(uD))/D^2,&4\le uD\le2\pi.
 \end{cases}
\]
\end{lemma}
\begin{proof}
Since $\E(X-X')^2=2$, the measure
$(X-X')^2d\Pp/2$ is a probability measure. Its mean of $|X-X'|$ is
$\E|X-X'|^3/2$. The cosine defect bound and Jensen's inequality give
\[
 1-|f(u)|^2
 \ge\E[(u(X-X'))^2q(u|X-X'|)]
 \ge2u^2q\bigl(u\E|X-X'|^3/2\bigr)
 \ge2u^2q(uD).
\]
Finally $|x-y|^3\le(x-y)^2(|x|+|y|)$ and centering imply
$\E|X-X'|^3\le2\rho+2\E|X|\le2(\rho+1)$.
No stronger universal independent-copy bound is being asserted.
\end{proof}

\paragraph{A pointwise complex-power enclosure}

\begin{lemma}[Rectangle gate]\label{ms:fourier:afcore:lem:rectangle}
Let $n\ge2$ be an integer, $a\ge0$, $0<x_0\le x_1$, and $Y\ge0$.
If
\[
 (n-2)Y\le\frac32x_0,\qquad
 (2n-1)x_0^n\ge(n-1)a,
\]
then throughout $x\in[x_0,x_1]$, $|y|\le Y$,
\[
 |(x+iy)^n-a|
 \le\max\{|(x_0+iY)^n-a|,|(x_1+iY)^n-a|\}.
\]
\end{lemma}
\begin{proof}
Put $z=x+iy$, $r=|z|$, and $F(x,y)=|z^n-a|^2$.
Since $|\arg z|\le|y|/x$, the first gate gives
$\Re z^{n-2}\ge0$. Direct differentiation yields
\begin{align*}
 F_{yy}&=2nr^{2n-4}\{r^2+2(n-1)y^2\}
             +2an(n-1)\Re z^{n-2}\ge0,\\
 F_{xx}&=2nr^{2n-4}\{r^2+2(n-1)x^2\}
             -2an(n-1)\Re z^{n-2}.
\end{align*}
Divide the second expression by $2nr^{n-2}>0$. Its positive term is
$r^{n-2}\{(2n-1)x^2+y^2\}\ge(2n-1)x_0^n$, and its negative term
has magnitude at most $(n-1)a$. Thus $F$ is separately convex.
Its maximum on the rectangle is bounded by its four corners; conjugate
corners have equal values. Equality in either gate, $n=2$, and $Y=0$
are included.
\end{proof}

For a frequency $0\le u<\pi/2$ one may take
\[
 x_0=\cos u,\quad x_1=\cos u+(\rho_+-1)u^3/6,\quad
 a=e^{-nu^2/2},
\]
and any valid uniform imaginary upper $Y$. The gates must hold on the
whole integration cell before this enclosure is used there.
Without those gates, an enclosing disk still gives a complete bound.
Indeed, if $M\ge|f(u)|$ and $E\ge|f(u)-g(u)|$, then for every $n\ge1$
\begin{equation}\label{ms:fourier:afcore:eq:telescoper}
 |f(u)^n-g(u)^n|
 \le\min\left\{M^n+g(u)^n,
 E\sum_{j=0}^{n-1}M^jg(u)^{n-1-j}\right\}.
\end{equation}
This is the finite factorization of $z^n-w^n$ and remains valid at equal
or zero bases. For a signed real interval, both endpoint squares are
retained when constructing either disk.

\paragraph{The smoothing interface}

We use Prawitz's smoothing inequality in the four-term normalization
of \cite[Lemma~4.1, p.~263]{Shevtsova2012}. This is a cited external inequality,
not a consequence of the preceding elementary lemmas. Put
\[
 \mathcal K(t)=(1-t)+i\{(1-t)\cot(\pi t)+1/\pi\},\qquad0<t<1.
\]
For $c>0$, $0<\tau<1$, and every integer $n\ge1$, it gives
\begin{align}
 \Delta_n\le{}&
 \int_0^\tau |\mathcal K(t)|\,
             |f(ct)^n-g(ct)^n|\,dt
 +\int_\tau^1|\mathcal K(t)|\,|f(ct)|^n\,dt\nonumber\\
 &+\int_0^\tau\left|\mathcal K(t)-\frac{i}{\pi t}\right|
                    g(ct)^n\,dt
 +\frac1\pi\int_\tau^\infty\frac{g(ct)^n}{t}\,dt.
 \label{ms:fourier:afcore:eq:prawitz}
\end{align}
Here $\Delta_n$ is the Kolmogorov distance for $S_n/\sqrt n$.
All four terms are nonnegative. The last is
$E_1(nc^2\tau^2/2)/(2\pi)$, where
$E_1(x)=\int_x^\infty e^{-v}\,dv/v\le e^{-x}/x$ for $x>0$.
Normal continuity equates the suprema defined with strict and closed
CDFs, so atoms introduce no additional restriction in this application.

On a moment band, multiply a complete upper bound for the right side
by $\sqrt n/\rho_-$, not by $\sqrt n/\rho_+$.
A finite partition yields an integral upper bound by summing width
times a valid whole-cell majorant. If a fixed $C^2$ majorant $F$ is
used on a cell of width $w$ and midpoint $m$, Taylor integration gives
\[
 \int F\le wF(m)+\frac{w^3}{24}\max\{0,\sup F''\}.
\]
Taking a minimum with a separately valid Darboux upper is legitimate.
Selecting a favorable branch at the midpoint without proving it on
the entire cell is not. Endpoint pieces and the infinite normal tail
remain separate obligations in every finite certificate.

\subsubsection{Complete Fourier remainders and finite resource enclosure}
\label{ms:fourier:rescore:sec:setup}
Use the contact normalization and the gap quantities of
Section~\ref{ms:fourier:afcore:sec:conserved}. Thus $C$ is the actual positive
closed-CDF score of one unrestricted attaining law, $H$ is its raw
sum threshold, $\rho=\E|X|^3$, $\gamma=\E X^3$,
$\beta=\E[X|X|]$. The quantities $D,J$ are defined in
\eqref{ms:physical-defect-definition}, and $J_1,J_2,L$ in
\eqref{ms:fourier:afcore:eq:gap-definitions}. Let
$\mu_j=\E X^j$ for even $j$, $\eta_5=\E|X|^5$, and
$\varphi=\phi(H/\sqrt n)$. Support $[-R_-,R_+]$ is acquired before
these higher moments are invoked; write $R=\max(R_-,R_+)$.
In applications below $1\le R\le6$.
The exact raw remainder $\mathcal R$ is defined in
\eqref{ms:fourier:rescore:re-raw}; its contact derivation, compensated
inversion and phase price are proved in Section~\ref{ms:conserved}.
Here we bound its Fourier and moment terms on explicit domains.
The elementary constant enclosures used here are
\[
 997/2500<\phi(0)<2/5,\qquad
 0.4095<\CE<0.40974,\qquad 333/106<\pi<22/7.
\]
The Gaussian and Esseen-constant bounds follow from
Lemma~\ref{ms:gaussian-constants}. For the slightly sharper lower bound
on $\pi$ used here, the same alternating sums give
\[
 \frac{333}{106}
 <16\sum_{j=0}^{5}\frac{(-1)^j}{(2j+1)5^{2j+1}}-\frac4{239}
 <\pi.
\]
The first comparison is rational; the second uses exactly the lower
arctangent truncations in that lemma.
The fractional-part estimate below requires finer directed bounds.
With
\[
 A_m(t)=\sum_{j=0}^{m-1}\frac{(-1)^jt^{2j+1}}{2j+1},\qquad
 \epsilon_5=\frac{5^{-65}}{65},\quad
 \epsilon_{239}=\frac{239^{-21}}{21},
\]
the alternating remainders and the same Machin identity give
\begin{align*}
 3.141592653589793
 &<16A_{32}(1/5)-4\{A_{10}(1/239)+\epsilon_{239}\}<\pi,\\
 \pi&<16\{A_{32}(1/5)+\epsilon_5\}-4A_{10}(1/239)
 <3.141592653589794.
\end{align*}
The two outer comparisons involve only the displayed rational numbers.

A finite-index profile means
\begin{equation}\label{ms:fourier:rescore:u-profile}
\begin{gathered}
 N\le n\le L,\qquad n(C-\CE)\ge K\ge0,\qquad C\in[c_0,c_1],\\
 n(C_{n-1}-C)_+\le q_*,\qquad c_0=0.4095+K/L.
\end{gathered}
\end{equation}
A band-uniform bound is used only under the explicit monotonicity
gates proved below. Otherwise every formula is evaluated at the
actual integer $n$. In particular, no monotonicity of $C_n$ is assumed.
The quantitative support implication is a separate contact argument,
not an output of the Fourier integration.

\paragraph{One complete low-frequency transfer}
Put $f(t)=\E e^{itX}$, $h(t)=e^{-t^2/2-i\gamma t^3/6}$,
$e_n(t)=e^{-nt^2/2}(1-in\gamma t^3/6)$, and
\[
 I_0(H)=\int_0^{1/2}|\Im(e^{-itH}(f^n-e_n))|\frac{dt}{t^3},\qquad
 I_\infty=\int_{1/2}^\infty|f^n-e_n|\frac{dt}{t^3}.
\]
The inversion bound~\eqref{ms:fourier:afcore:eq:error-ledger} charges
$2\sqrt n(I_0+I_\infty)/\pi$; both integrals below are complete.

\begin{lemma}[Signed Gaussian transfer]\label{ms:fourier:rescore:fourier-early-low}
\label{ms:fourier:rescore:fourier-paired-low}\label{ms:fourier:rescore:entryleap-fourier-transfer}
Suppose $n\ge N>2$, $\mu_4\le M\le11.1$, $|\gamma|\le g\le1.85$,
$\mu_6\le M_6$, and the actual coefficients
\[
 a=(3-\mu_4)/24,\quad z=(10\gamma-\E X^5)/120,\quad c=a\gamma/6
\]
satisfy $|a|\le a_*$, $|z|\le Z$. Put
\begin{equation}\label{ms:fourier:rescore:fourier-taylor}
 E_6=M_6/720+g^2/72+(1/2+g^2/144)^3/5,\quad
 B_6=a_*+Z/2+E_6/4.
\end{equation}
Let $a_0=(N-1)/(2N)$,
$0<b_0\le(45-M)(1-2/N)/96$, and
$\lambda_N=\sqrt{\pi/2}(N/(N-1))^{3/2}$.
Then, for every real $H$,
\begin{equation}\label{ms:fourier:rescore:fourier-low-form}
 I_0(H)\le(L_0+L_H|H|)/\sqrt n+L_2/n,
 \qquad L_2=E_6/(2a_0^2)+B_6^2/(2b_0^3).
\end{equation}
Two admissible choices, each used as a complete triple, are as follows.
The triangle choice is
\begin{equation}\label{ms:fourier:rescore:fourier-early}
 L_0=c_0Z+c_1a_*g+251g^3/17280,\quad
 L_H=c_0a_*+251g^2/4800,
\end{equation}
where $c_0\ge\lambda_N$, $c_1\ge\lambda_N/2$.
For the cancellation choice, take any rectangle containing the actual
$(z,c)$, $I_u\ge\sqrt{\pi/2}$ and $F_N\ge(N/(N-1))^{3/2}$. Define
\begin{gather*}
 U(z,c)=\begin{cases}
 I_u(|z|+3|c|),&zc\le0,\\
 I_u\bigl||z|-3|c|\bigr|+6\sqrt{|zc|}e^{-|z|/(2|c|)},&zc>0,
 \end{cases}\\
 J_u=\max_{\rm corners}U(z,c).
\end{gather*}
Then
\begin{align}\label{ms:fourier:rescore:fourier-paired}
 L_0&=F_NJ_u+\frac{F_NI_u}{N-1}
 (35a_*g^3/48+35Zg^2/24)+251g^3/17280,\notag\\
 L_H&=F_NI_u\{a_*+5Zg/[2(N-1)]\}+251g^2/4800.
\end{align}
Neither the transfer nor its cancellation choice assumes signs for
$a,z,c$ or $\gamma$.
\end{lemma}
\begin{proof}
Taylor with integral remainder gives
$f-h=-at^4-izt^5+r_6$, $|r_6|\le E_6t^6$ for $0\le t\le1/2$.
Indeed for $w=-t^2/2-i\gamma t^3/6$,
$|w|\le t^2(1/2+g^2/144)<1/6$; the terms through degree five in
$1+w+w^2/2$ are exact, its degree-six term costs $g^2t^6/72$, and
the remaining exponential series costs $|w|^3e^{|w|}/6\le|w|^3/5$.
Add the true sixth-moment remainder.

Four integrations of $1-\cos u\ge0$ give
$\cos u\le1-u^2/2+u^4/24$. Symmetrization, using
$\E(X-X')^4=2\mu_4+6$, therefore gives
$|f(t)|,|h(t)|\le e^{-(45-M)t^2/96}$ on $[0,1/2]$.
Every point on their joining segment has the same bound, so polynomial
Taylor expansion yields
\[
 f^n-h^n=n(f-h)h^{n-1}+R_2,\qquad
 |R_2|\le\tfrac12n^2B_6^2t^8e^{-b_0nt^2}.
\]
The $r_6$ and $R_2$ integrals are at most
\begin{equation}\label{ms:fourier:rescore:fourier-common-error}
 \frac1n\{E_6/(2a_0^2)+B_6^2/(2b_0^3)\}=L_2/n.
\end{equation}
For the remaining comparator put $u=n\gamma t^3/6$. Its imaginary part
is bounded by $e^{-nt^2/2}(|u|^3/6+|H|tu^2/2)$, so its divided integral
is at most
\begin{equation}\label{ms:fourier:rescore:fourier-comparator-low}
 n^{-1/2}(251g^3/17280+251g^2|H|/4800).
\end{equation}
Here the Gaussian moments follow from integration by parts,
$\int_0^\infty x^{2j}e^{-x^2/2}dx=(2j-1)!!\sqrt{\pi/2}$,
and $\sqrt{2\pi}<251/100$.

It remains to price the leading quartic and fifth terms. Set $b=n-1$,
$v=b\gamma t^3/6$. The triangle estimate uses
$|\sin(Ht+v)|\le|H|t+bg t^3/6$, $|\cos|\le1$; the Gaussian
moments of orders two and four give respectively $\lambda_N$ and
$\lambda_N/2$. This proves~\eqref{ms:fourier:rescore:fourier-early}.
For cancellation, after division by $t^3$, rotation by $Ht$ bounds the
leading imaginary part by
\[
 |zt^2-atv|+a_*t|\sin v-v|+Zt^2|\cos v-1|
       +|H|t\{a_*t+Zt^2|\sin v|\}.
\]
Multiply by $ne^{-bt^2/2}$ and use the global sine/cosine remainders.
The first integral equals $nJ(z,c)/b^{3/2}$, where
$J(z,c)=\int_0^\infty e^{-x^2/2}|zx^2-cx^4|dx$.
This norm is convex and $J(-z,-c)=J(z,c)$, so its rectangle maximum
is at a corner. For $zc\le0$ the integral is
$\sqrt{\pi/2}(|z|+3|c|)$. For $z,c>0$, split at $r=\sqrt{z/c}$;
integration by parts gives
\[
 J=\sqrt{\pi/2}(z-3c)+6cr e^{-r^2/2}
       +2(3c-z)\int_r^\infty e^{-x^2/2}dx.
\]
Discard the last term when nonpositive, and otherwise bound its integral
by $\sqrt{\pi/2}$. This proves $J\le U$; simultaneous sign reflection
covers the remaining quadrant. The other four phase terms use moments
of orders two, six, eight and ten and give the coefficients
$1,5/2,35/24,35/48$ in~\eqref{ms:fourier:rescore:fourier-paired}. Finally $n/(n-1)$,
$1/(n-1)$, $a_0^{-1}$ and the allowable damping inverses decrease with
$n$, proving the all-index bounds at $N$.
\end{proof}

\begin{lemma}[Retaining the signed Gaussian tail]\label{ms:fourier:rescore:entry189-mills}
For $r>0$, put
\[
 \mathcal M(r)=e^{r^2/2}\int_r^\infty e^{-t^2/2}\,dt,
 \qquad
 \mathcal M_k(r)=\cfrac1{r+\cfrac1{r+\cfrac2{\ddots+\cfrac{k}{r}}}}
 \quad(k\ge1),\qquad \mathcal M_0(r)=1/r.
\]
Odd truncations are lower bounds for $\mathcal M$; even truncations
are upper bounds. For $z,c>0$, set $x=z/c$ and
\[
 d_{k+1}=1,\qquad d_j=1+\frac j{xd_{j+1}}\quad(j=k,\ldots,1),
 \qquad m_k(x)=d_1^{-1}.
\]
Then the cancellation norm in Lemma~\ref{ms:fourier:rescore:entryleap-fourier-transfer}
satisfies
\begin{equation}\label{ms:fourier:rescore:entry189-mills-bound}
 J(z,c)\le \sqrt{\pi/2}(z-3c)+\sqrt{zc}\,e^{-x/2}
 \left(6+\frac{2(3-x)}x m_k(x)\right),
\end{equation}
where $k=11$ if $x\ge3$ and $k=12$ if $x<3$.
The factor in parentheses is positive. One may take the minimum with
the previous bound $U(z,c)$ at each corner of the coefficient rectangle.
\end{lemma}
\begin{proof}
The positive integrals
$I_j(r)=\int_0^\infty t^j e^{-rt-t^2/2}\,dt$ satisfy
$I_1=1-rI_0$ and $I_{j+1}=jI_{j-1}-rI_j$ by integration by parts.
Thus, writing $q_j=I_j/I_{j-1}>0$,
\[
 I_0=\frac1{r+q_1},\qquad q_j=\frac j{r+q_{j+1}}.
\]
Replacing the final positive $q_{k+1}$ by zero reverses the inequality
at each substitution. This proves the alternating finite bounds, since
translation gives $I_0=\mathcal M(r)$. No convergence of an infinite
continued fraction is required. The displayed rational recursion gives
$m_k(x)=r\mathcal M_k(r)$ when $r=\sqrt x$.

Splitting the norm at $r=\sqrt{z/c}$, as in the preceding proof, gives
the exact identity
\[
 J(z,c)=\sqrt{\pi/2}(z-3c)+\sqrt{zc}\,e^{-x/2}
 \left(6+\frac{2(3-x)}x r\mathcal M(r)\right).
\]
The multiplier of $\mathcal M$ is nonpositive when $x\ge3$ and
positive when $x<3$. Substitute the corresponding finite lower or
upper bound. Positivity of the remaining factor is immediate for
$x<3$; for $x\ge3$, use $0<m_{11}(x)<1$ to bound it below by
$6-2(x-3)/x>4$. Finally the exact norm $J$, rather than either
upper-bound formula, is convex, so its containing rectangle is still
exhausted by its four corners.
\end{proof}

At every positive coefficient corner use the displayed finite Mills bound.
A negative Gaussian coefficient uses a lower enclosure for
$\sqrt{\pi/2}$; a nonnegative coefficient uses an upper one. The remaining
factor is positive, so outward upper square-root and exponential enclosures
are valid. A simultaneous reflection of both odd coefficients leaves the
norm unchanged.

\paragraph{The two coefficient choices in the certificates.}
All coefficients come from the actual marginal. For $Q=X^2-\gamma X-1$,
multiplication gives
\begin{align}\label{ms:fourier:rescore:fourier-pearson}
 \E Q^2&=\mu_4-1-\gamma^2,&
 \E X^5&=\gamma^3+2\gamma+\E[Q^2(X+2\gamma)],\notag\\
 \mu_6&=1+3\gamma^2+\gamma^4+
       \E[Q^2\{(X+\gamma)^2+2\gamma^2+2\}].
\end{align}
For any radius $R_0\ge|X|$ with $R_0\ge2g$, the two signs of
$\E X^5-10\gamma$ are bounded, at $t=|\gamma|$, by
\begin{align*}
 P_{R_0}(t)&=R_0(M-1)+(10-2M)t-R_0t^2+t^3,\\
 Q_{R_0}(t)&=R_0(M-1)+(2M-10)t-R_0t^2-t^3.
\end{align*}
The gate $R_0-2g\ge0$ permits substituting $M$ for $\mu_4$ in both
signs. The triangle choice uses the valid enclosing radius $R_0=6$, $a_* =\max(2,M-3)/24$,
$M_6=RE_5$, $b_0=(45-M)(1-2/N)/96$ and
\begin{equation}\label{ms:fourier:rescore:fourier-z0}
\begin{gathered}
 120Z=\min\{E_5+10g,Z_0\},\\
 Z_0=\begin{cases}
 P_6(g),&M\le5, P'_6(g)\ge0,\\
 6(M-1)+(10-2M)^2/[4(6-g)],&M\le5, P'_6(g)<0,\\
 6(M-1)+(2M-10)^2/24,&M>5.
 \end{cases}
\end{gathered}
\end{equation}
If $M\le5$, $P_6\ge Q_6$ and $P_6''<0$ on $[0,g]$; use the
increasing endpoint or $t^3\le gt^2$ and complete the square.
If $M>5$, $P_6\le6(M-1)$ and the upper quadratic for $Q_6$ gives
the third row. One may also use the actual radius by setting
$r_0=\max(R,2g)>g$ and
\[
\begin{gathered}
 t_*=\begin{cases}
 \min\{g,(10-2M)/[2(r_0-g)]\},&M\le5,\\
 \min\{g,(2M-10)/(2r_0)\},&M>5,
 \end{cases}\\
 Z_r=\begin{cases}
 r_0(M-1)+(10-2M)t_*-(r_0-g)t_*^2,&M\le5,\\
 r_0(M-1)+(2M-10)t_*-r_0t_*^2,&M>5.
 \end{cases}
\end{gathered}
\]
Indeed for $M\le5$, $P_{r_0}\ge Q_{r_0}$ and $t^3\le gt^2$;
for $M>5$, $P_{r_0}\le r_0(M-1)$ and dropping $-t^3$ bounds
$Q_{r_0}$. Each displayed quadratic is maximized at its clipped vertex.
Thus a permitted coefficient choice is
$120Z=\min(E_5+10g,Z_0,Z_r)$.
It uses $c_0\ge\sqrt{\pi/(16a_0^3)}=\lambda_N$ and $c_1=c_0/2$.
The weaker choice $c_1\ge\sqrt{\pi/(256a_0^5)}
=\lambda_N N/[2(N-1)]$ is also permitted by the same transfer lemma.

For the paired choice the already proved cell has
$0\le\gamma_l\le|\gamma|\le g$, $M<3$, $g^2<2$, and $R>2g$. After reflecting only this auxiliary Fourier integral if necessary, set
\begin{gather*}
 a_*=(2-\gamma_l^2)/24,\qquad b_0=((45-M)/96)(1-2/N),\\
 M_6=\min\{RE_5,1+3g^2+g^4+(R^2+2Rg+3g^2+2)(M-1-\gamma_l^2)\},\\
 t_*=\min\{g,(10-2M)/[2(R-g)]\},\\
 120Z=\min\{E_5+10g,R(M-1)+(10-2M)t_*-(R-g)t_*^2\},\\
 z_l=\max\{-Z,[8\gamma_l-\gamma_l^3-(R+2\gamma_l)(M-1-\gamma_l^2)]/120\},
 \quad z_u=Z,\\
 c_l=(3-M)\gamma_l/144,\qquad c_u=(2-\gamma_l^2)g/144.
\end{gather*}
Check $M-1-\gamma_l^2\ge0$ and that the rectangle is ordered.
The extension from the formerly sufficient $M\le2$, $g\le1$ domain is
legitimate: Pearson gives $\mu_4\ge1+\gamma^2$, hence
$0\le a\le(2-\gamma_l^2)/24$; $M<3$ keeps $c_l\ge0$ and $g^2<2$
keeps the upper quartic coefficient nonnegative. The general transfer
hypotheses remain true since $M<3<11.1$ and $g<\sqrt2<1.85$.
The damping must be $(45-M)(1-2/N)/96$; the older value
$43(1-2/N)/96$ is not licensed when $M>2$. Here
$P_R\ge Q_R$, and $P_R(t)\le R(M-1)+(10-2M)t-(R-g)t^2$,
whose maximum is at $t_*$. The numerator of the lower $z$ bound increases
in $\gamma\ge0$ because its derivative is
$10-2M+2R\gamma+3\gamma^2>0$. Pearson's first and third identities
give $a_*$ and $M_6$. If $3g^2\le2$, the increasing function
$2\gamma-\gamma^3$ further licenses $c_u\le(2g-g^3)/144$.
Thus this rectangle is a specialization of the signed transfer on the
whole enlarged domain. Its actual moment inputs, including any shared
lower bound $r_*\le\rho$, are acquired before it is used. The signed
phase calculation keeps the original $\gamma$; reflection here acts only
on the auxiliary Fourier norm and both its odd coefficients.

\paragraph{A low-frequency bound on a finite interval}
\label{ms:fourier:rescore:sec:finite-low}

The preceding low-frequency bounds extend positive Gaussian majorants to infinity.
At small indices this can be wasteful. The following alternative keeps
the original integration interval; it uses no physical-gap premise.

\begin{lemma}[Finite-window transfer]\label{ms:fourier:rescore:lem:finite-low}
Suppose $n>2$ is an integer, $1\le R\le6$, $|X|\le R$,
$\mu_4\le M\le111/10$, $|\gamma|\le g\le37/20$,
and $\eta_5\le E_5$. Let $0\le g_l\le|\gamma|$ with $g_l^2\le M-1$.
Set
\begin{gather*}
 M_6=\min\{RE_5,
 1+3g^2+g^4+(R^2+2Rg+3g^2+2)(M-1-g_l^2)\},\\
 A_*\ge |(3-\mu_4)/24|,\quad
 Z\ge|(10\gamma-\E X^5)/120|,\\
 E=M_6/720+g^2/72+(1/2+g^2/144)^3/5,\qquad
 B=A_*+Z/2+E/4,\\
 a_n=(n-1)/2,\qquad b_n=(n-2)(45-M)/96,\qquad
 I_j(A)=\int_0^{1/2}t^je^{-At^2}\,dt.
\end{gather*}
Then $I_0(H)$ is bounded by $(L_0+L_H|H|)/\sqrt n+L_2/n$, with
the following complete triple:
\begin{align}
 L_0={}&\sqrt n\{nZI_2(a_n)
       +n(n-1)A_*gI_4(a_n)/6+n^3g^3I_6(n/2)/1296\},\nonumber\\
 L_H={}&\sqrt n\{nA_*I_2(a_n)+n^2g^2I_4(n/2)/72\},\nonumber\\
 L_2={}&n\{nEI_3(a_n)+n(n-1)B^2I_5(b_n)/2\}.
 \label{ms:fourier:rescore:eq:finite-low-triple}
\end{align}
The bound holds for every real raw threshold $H$.
One always may use
$A_*=\max\{|3-M|,|2-g_l^2|\}/24$ and $Z=(E_5+10g)/120$;
the sharper proved Pearson coefficient bounds above may be substituted.
\end{lemma}
\begin{proof}
Pearson's sixth-moment identity
\eqref{ms:fourier:rescore:fourier-pearson}, whose residual multiplier is positive,
gives the stated $M_6$; also $\mu_6\le R\eta_5$ pointwise.
Put $a=(3-\mu_4)/24$ and $z=(10\gamma-\E X^5)/120$.
For $h(t)=e^{-t^2/2-i\gamma t^3/6}$,
\[
 f(t)-h(t)=-at^4-izt^5+r_6(t),\qquad |r_6(t)|\le Et^6
 \quad(0\le t\le1/2).
\]
Indeed, with $w=-t^2/2-i\gamma t^3/6$,
$|w|\le t^2(1/2+g^2/144)<1/6$. Comparing with $1+w+w^2/2$
costs $g^2t^6/72+|w|^3e^{|w|}/6\le
\{g^2/72+(1/2+g^2/144)^3/5\}t^6$, besides the true sixth-moment
remainder $\mu_6t^6/720$. Thus $|f-h|\le Bt^4$ on the whole interval.

Symmetrization gives
\[
 |f(t)|^2\le1-t^2+(M+3)t^4/12
 \le1-(45-M)t^2/48.
\]
Both $|f(t)|$ and $|h(t)|$ are therefore at most
$e^{-(45-M)t^2/96}$, as is every point on their complex line segment.
Taylor's integral formula for the polynomial $z^n$ yields
\[
 f^n-h^n=n(f-h)h^{n-1}+R_2,\qquad
 |R_2|\le\frac{n(n-1)}2B^2t^8e^{-b_nt^2}.
\]
The $r_6$ and $R_2$ terms, after division by $t^3$ and integration,
cost $nEI_3(a_n)+n(n-1)B^2I_5(b_n)/2$.
The leading imaginary part is at most
\[
 ne^{-a_nt^2}\{(Z+A_*|H|)t^5+(n-1)A_*gt^7/6\}.
\]
Finally, for $u=n\gamma t^3/6$,
\[
 |\Im[e^{-itH}(h^n-e_n)]|
 \le e^{-nt^2/2}\{|u|^3/6+|H|t u^2/2\}.
\]
Integrating these terms gives every coefficient in
\eqref{ms:fourier:rescore:eq:finite-low-triple}, with no replacement of $n(n-1)$
by $n^2$ and no omitted frequency interval.
\end{proof}

For rational $A\ge0$, exact enclosures for these finite integrals are
obtained from
\begin{equation}\label{ms:fourier:rescore:eq:finite-gaussian-integral}
 P_{m,j}(A)=\sum_{k=0}^m
 \frac{(-A)^k}{k!(j+2k+1)2^{j+2k+1}},\qquad
 \max\{0,P_{21,j}(A)\}\le I_j(A)\le P_{20,j}(A).
\end{equation}
The integral Taylor remainder of $e^{-x}$ has sign $(-1)^{m+1}$
for every $x\ge0$, not merely on an alternating-series domain.
Integrate its polynomial bounds after $x=At^2$.
A finite table restricted to $A\in[0,8]$ and $0\le j\le16$
is used only when these conditions hold.
Every weight in \eqref{ms:fourier:rescore:eq:finite-low-triple} is nonnegative.
An alternative low-frequency choice is always one entire valid triple; taking
coordinatewise minima of triples is not justified.

\paragraph{The complete high-frequency integral}
For $b,\kappa>0$ and $n>0$, put
\[
 T_n=\sqrt n\int_b^\infty e^{-n\kappa t^2}t^{-3}\,dt,
 \qquad x=n\kappa b^2.
\]
\begin{lemma}[Rational Laguerre tail bound]\label{ms:fourier:laguerre}
For every $x>0$,
\[
 T_n\le\frac{\sqrt n e^{-x}}{2b^2}
                  \frac{x+2}{x^2+4x+2}.
\]
If $N\kappa b^2\ge1/4$, the right side decreases with $n\ge N$.
\end{lemma}
\begin{proof}
Substitute $t^2=b^2(1+s)$ and use positive Fubini to obtain
\[
 T_n=\frac{\sqrt n e^{-x}}{2b^2}
 \int_0^\infty\frac{e^{-xs}}{(1+s)^2}\,ds
 =\frac{\sqrt n e^{-x}}{2b^2}
 \int_0^\infty\frac{u}{x+u}e^{-u}\,du.
\]
Put $P(u)=u^2-4u+2$ and $D(x)=x^2+4x+2=P(-x)>0$.
Exact division gives
\[
 \frac1{x+u}-\frac{Q_x(u)}{D(x)^2}
       =\frac{P(u)^2}{(x+u)D(x)^2}\ge0,
\]
where
\[
 Q_x(u)=-u^3+(x+8)u^2-(x^2+8x+20)u
                  +x^3+8x^2+20x+16.
\]
Integration uses only $\int_0^\infty u^j e^{-u}du=j!$, proved by
integration by parts. It gives
$\int Q_x(u)e^{-u}du=(x+3)D(x)$ and hence
\[
 \int_0^\infty\frac{u}{x+u}e^{-u}du
 =1-x\int_0^\infty\frac{e^{-u}}{x+u}du
 \le1-\frac{x(x+3)}{D(x)}=\frac{x+2}{D(x)}.
\]
For the all-index statement, differentiate
$\sqrt{x}e^{-x}(x+2)/D(x)$. After multiplication by positive factors,
its derivative has numerator
$4-10x-22x^2-13x^3-2x^4$. This polynomial is negative at $x=1/4$
and strictly decreasing for $x>0$. Thus it is negative on the entire
required half-line.
\end{proof}


Let $\E X=0$, $\E X^2=1$, $\E X^4\le M$, and write
$\gamma=\E X^3$, $f(t)=\E e^{itX}$. For independent copies suppose an
actual event $E$ satisfies
\[
 \Pp(E)\ge a>0,\qquad |X-X'|\ge d_0>0\text{ on }E,\qquad
 \E[(X-X')^2\mid E]\le M_E.
\]
The event is measured under the product of the same marginal law.
Define the Taylor comparator
$e_n(t)=e^{-nt^2/2}(1-in\gamma t^3/6)$ and
\[
 I_\infty=\int_{1/2}^{\infty}|f(t)^n-e_n(t)|t^{-3}\,dt,
 \qquad L(x)=\frac{x+2}{x^2+4x+2}.
\]
The sign of the imaginary term is immaterial for its absolute-value bound.
Use $b=9/10$, $0<r\le1/4$,
\begin{gather*}
 h=\frac{333}{106}-\frac{(22/7)^3r^2}{6}>0,\quad
 \kappa=a(23/48)^2d_0^2,\\
 S_3=\sum_{j=1}^{20}j^{-3}+\frac1{800},\qquad
 S_5=\sum_{j=1}^{20}(j-r)^{-5}+\frac1{4(20-r)^4}.
\end{gather*}
Set
\[
 \begin{gathered}c=\frac{1-\max_{u\in\{1/4,81/100\}}
              [1-u+(M+3)u^2/12]}2,\\
 c_0=\frac18-\frac{M+3}{384},\quad
 s=\frac{1-(M+3)53/600}{2}.\end{gathered}
\]
The usual damping domain requires $c>0$.

\begin{lemma}[Six-part high-frequency bound]\label{ms:fourier:window50}
For every integer $n\ge5$, the complete upper bound is
\begin{align*}
 \sqrt n I_\infty\le{}& A_n
 +\frac{M_E}{4\pi^{3/2}h\sqrt a}
       \left(S_3+\frac{3S_5}{nah^2}\right)
 +\frac{M_E\sqrt n}{2\pi^2}e^{-nah^2r^2}\\
 &+Z_n+\frac{\sqrt n}{2b^2}e^{-na(23/48)^2}
 +\sqrt n e^{-n/8}\left(2L(n/8)+\frac{\sqrt{M-1}}3\right),
\end{align*}
where always $A_n=2\sqrt n e^{-nc}$ is permitted, and, if $s>0$,
it can be replaced by the smaller of itself and
\[
 2\sqrt n e^{-nc_0}L(ns/4).
\]
The zero-cell term $Z_n$ may be taken as the smaller of
\[
 \frac{e^{-n\kappa b^2}}{2b^4\kappa\sqrt n},\qquad
 \frac{\sqrt n e^{-n\kappa b^2}}{2b^2}L(n\kappa b^2).
\]
Each term may be evaluated at a common floor $N$ for all $n\ge N$
provided
\[
 2Nc'>1\quad(c'\in\{c,ah^2r^2,a(23/48)^2,1/8\}).
\]
For the new first term also require $2Nc_0>1$;
for the second zero-cell option require $N\kappa b^2\ge1/4$.
If either extra condition is unavailable, retain the corresponding simpler term.
Any fixed collection of permitted windows may be evaluated, and the minimum
of its complete six-term bounds remains valid.
\end{lemma}

\begin{proof}
We give the entire frequency partition to make clear what changes.
On $[1/2,b]$, symmetrization and
$\cos u\le1-u^2/2+u^4/24$ yield
\[
 |f(t)|^2\le1-t^2+(M+3)t^4/12,
 \qquad |f(t)|^n\le e^{-n c(t^2)},
 \quad c(u)=u/2-(M+3)u^2/24.
\]
Bounding $c(t^2)$ below by $c$ and using
$\int_{1/2}^{b}t^{-3}dt<2$ gives the first permitted term.
Concavity supplies the more precise chord bound
\[
 c(u)\ge c_0+s(u-1/4),\qquad 1/4\le u\le81/100,
\]
because the difference is exactly
$(M+3)(u-1/4)(81/100-u)/24$. If $s>0$, extend only the resulting
positive exponential majorant to the half-line. The preceding rational
Gaussian-tail lemma, with lower endpoint $1/2$, gives
\[
 \sqrt n\int_{1/2}^{b}e^{-nc(t^2)}t^{-3}dt
 \le e^{-n(c_0-s/4)}\sqrt n
       \int_{1/2}^{\infty}e^{-ns t^2}t^{-3}dt
 \le2\sqrt n e^{-nc_0}L(ns/4).
\]

For $t\ge b$, conditioning on the actual event and convex Jensen give
\[
 |f(t)|^n\le e^{-nE\sin^2(t(X-X')/2)}
 \le \E[e^{-na\sin^2(td/2)}\mid E],\qquad d=|X-X'|.
\]
Positive Tonelli allows a separate partition for every $d>0$.
Put $\ell=2\pi/d$. The zero cell is $[0,\ell/2]$; the other cells are
$\ell[j-1/2,j+1/2]$, $j\ge1$. In each nonzero cell take the core
$t=\ell(j+v)$, $|v|\le r$. These cores are disjoint even for $r\le1/2$.
For our smaller range $r\le1/4$, the elementary sine bound gives
$|\sin(\pi v)|\ge h|v|$ throughout a core, and
$|\sin(\pi v)|\ge hr$ throughout its wings. The latter uses that
$\sin(\pi|v|)$ increases on $[0,1/2]$.

Since $j-r>0$, Taylor's formula on the two sides gives
\[
 \frac{(j-v)^{-3}+(j+v)^{-3}}2
 \le j^{-3}+6(j-r)^{-5}v^2.
\]
Integrate this majorant against $e^{-nah^2v^2}$ on the full real line.
The result is
$\sqrt\pi[j^{-3}/\sqrt{nah^2}
  +3(j-r)^{-5}/(nah^2)^{3/2}]$.
The Jacobian is
$t^{-3}dt=d^2(4\pi^2)^{-1}(j+v)^{-3}dv$.
Summing and averaging $d^2$ gives the second term. The two series tails
are bounded by their decreasing integrals after $20$, producing exactly
$S_3,S_5$. The total wing weight is at most
$\int_{\ell/2}^{\infty}t^{-3}dt=d^2/(2\pi^2)$, giving the third term.
Extending cores or wings below $b$ only adds a positive integral.

In the zero cell write $u=td/2\le\pi/2$.
For $u\le1/2$, use $\sin u\ge(23/24)u$ and the rational
Gaussian-tail lemma with rate $\kappa$; its two choices give $Z_n$.
For $u\ge1/2$, use $\sin u\ge23/48$ and total weight
$1/(2b^2)$, giving the fifth term. Shared endpoints have zero integration
measure, including when the conditional distance distribution has atoms.

Finally keep the comparator on the entire interval $[1/2,\infty)$.
Its absolute-value bound splits its integral into
\[
 \sqrt n\int_{1/2}^{\infty}e^{-nt^2/2}t^{-3}dt
 +\frac{n^{3/2}|\gamma|}{6}
                 \int_{1/2}^{\infty}e^{-nt^2/2}dt.
\]
The first part is at most $2\sqrt n e^{-n/8}L(n/8)$.
For the second, $t\ge1/2$ implies
\[
 \int_{1/2}^{\infty}e^{-nt^2/2}dt
 \le2\int_{1/2}^{\infty}t e^{-nt^2/2}dt
 =2e^{-n/8}/n.
\]
Pearson's elementary identity
$\E[(X^2-1-\gamma X)^2]=\E X^4-1-\gamma^2\ge0$
gives $|\gamma|\le\sqrt{M-1}$ and proves the sixth term.
Any separately proved upper bound $|\gamma|\le g$ may replace
$\sqrt{M-1}$ in that term. A permitted choice is
$g=\min(1.85,\sqrt{M-1})$.

For uniformity, $\sqrt n e^{-nc'}$ decreases when $2nc'>1$,
whereas $n^{-1}$ and $n^{-1/2}e^{-nc'}$ always decrease.
The derivative of $L(x)$ is
$-(x^2+4x+6)/(x^2+4x+2)^2<0$.
The new first term is thus decreasing under its stated gate, and the
comparator decreases for every $n\ge5$. The earlier tail lemma proves
the separate $N\kappa b^2\ge1/4$ gate for its new zero-cell option.
These comparisons concern the bounding expressions, so they are uniform
even when the marginal law varies with $n$.
\end{proof}\paragraph{Events, not artificial kernels.}\label{ms:fourier:rescore:fourier-events}
The entry uses only the following actual choices, with the direct rate
from the same fourth moment in each case:
\[
\begin{array}{c|c|c|c}
 \text{event}&a&d_0&M_E\\\hline
 (X-X')^2\ge9/25 &(2-9/25)^2/(2M+6)&3/5&2/a\\
 XX'<0, D\le d\le0.8 &1/(2r_*^2)&\sqrt{2-d}&4r_*^2\\
 XX'<0, p\in[l,u] &2\min(l(1-l),u(1-u))
 &\max(A_l+B_l,\sqrt{2-d})&2/a
\end{array}
\]
The first row follows from Cauchy using
$\E[(X-X')^2;E]\ge2-d_0^2$ and
$\E(X-X')^4=2\mu_4+6$. The second is Lemma~\ref{ms:fourier:rescore:entry201-event}; all its moment and $\alpha<3$ gates are checked at use.
In the third, $D\le d<2$ ensures a genuine central gap,
$\Pp(XX'<0)=2pq$, and restriction of the squared distance gives $M_E=2/a$.
Taking a minimum of \emph{complete} high-frequency bounds is legitimate. No tail is
dropped and the third construction also allows a cell crossing $1/2$.

\paragraph{A compensated high-frequency bound with bounded support}
\label{ms:fourier:rescore:sec:compensated-high}

Here is a second complete high-frequency inequality. It is useful at fixed small
indices and requires neither a physical gap nor a predecessor estimate.
Its support and moments are hypotheses, not outputs of the integration.

\begin{lemma}[Supported compensated-power bound]
\label{ms:fourier:rescore:lem:compensated-high}
Let $n\ge2$ be an integer, $|X|\le R$, $\rho\le r$,
$\E X^4\le M_4$, $\E|X|^5\le M_5$, and $\gamma^2\le G$.
Let $1/2<T<\pi$, and let $m(u)$ be any valid upper bound for $|f(u)|$
on $[1/2,T]$. Put $g=e^{-u^2/2}$ and
\begin{gather*}
 b(u)=\sqrt G\,u^3/6,\qquad W(u)=g\sqrt{1+b(u)^2},\\
 d_0(u)=\cos u-g,\qquad d_1(u)=d_0(u)+(r-1)u^3/6,\\
 A(u)=\max\{|d_0(u)|,|d_1(u)|\},\\
 B(u)=\min\{M_4u^4/24,M_5u^5/120\}+b(u)(1-g),\qquad
 R_f(u)=\sqrt{A(u)^2+B(u)^2},\\
 \mathcal C_n(u)=R_f(u)\sum_{j=0}^{n-1}m(u)^jW(u)^{n-1-j}
 +g^n\frac{n(n-1)b(u)^2}{2}(1+b(u)^2)^{(n-2)/2},\\
 \mathcal T_n(u)=m(u)^n+g^n\sqrt{1+n^2b(u)^2}.
\end{gather*}
For every partition $1/2=t_0<\cdots<t_s=T$, let $U_j$ be a valid
whole-cell upper bound for $\min(\mathcal C_n,\mathcal T_n)$ on
$[t_{j-1},t_j]$. Then
\begin{align}
 \sqrt n I_\infty\le\sqrt n\left\{
 \sum_{j=1}^s U_j\frac{t_{j-1}^{-2}-t_j^{-2}}2
 +\frac1{2T^2}
 +\frac{e^{-nT^2/2}}{nT}
       \left(T^{-3}+\frac{n\sqrt G}{6}\right)\right\}.
 \label{ms:fourier:rescore:eq:compensated-high}
\end{align}
\end{lemma}
\begin{proof}
Write $v=-i\gamma u^3/6$ and $w=g(1+v)$.
The anchored real interval of Lemma~\ref{ms:fourier:afcore:lem:anchor} gives
$|\Re(f-w)|\le A$. Taylor's theorem for the sine gives the two
independent bounds
\[
 |\Im f(u)+\gamma u^3/6|\le M_4u^4/24,
 \qquad |\Im f(u)+\gamma u^3/6|\le M_5u^5/120.
\]
Consequently $|f-w|\le R_f$ and $|w|\le W$.
The finite telescoping factorization gives
\[
 |f^n-w^n|\le R_f\sum_{j=0}^{n-1}m^jW^{n-1-j}.
\]
No division, $W\le1$, or distinct-root hypothesis is used.
Taylor's integral formula for $s\mapsto(1+sv)^n$ on $[0,1]$ gives
\[
 |w^n-g^n(1+nv)|\le
 g^n\frac{n(n-1)b^2}{2}(1+b^2)^{(n-2)/2},
\]
since $v$ is purely imaginary. These two inequalities prove the
compensated bound $\mathcal C_n$. The independent triangle bound is
$\mathcal T_n$; both bound the same full remainder.
Integration on a closed cell uses the exact positive weight
$\int_a^b u^{-3}\,du=(a^{-2}-b^{-2})/2$.

Beyond $T$, use $|f|\le1$ and
$|1-in\gamma u^3/6|\le1+n\sqrt G u^3/6$.
The characteristic contribution is $1/(2T^2)$.
The Gaussian contribution is at most
\[
 \left(T^{-3}+n\sqrt G/6\right)
 \int_T^\infty e^{-nu^2/2}\,du
 \le\frac{e^{-nT^2/2}}{nT}
       \left(T^{-3}+n\sqrt G/6\right).
\]
This includes the entire infinite tail and proves the result.
\end{proof}

Valid choices of $m$ and $G$ are supplied by
Lemmas~\ref{ms:fourier:afcore:lem:sine-circle} and~\ref{ms:fourier:afcore:lem:modulus}
and the signed coordinate disk. The support-moment bounds below supply
$M_4,M_5$ for any actual radius $R$; a historical radius used in a
different index is not an input at the present index.

The real endpoints can be enclosed sharply without differentiating a
characteristic function: on a cell $[a,b]$, midpoint $m_c$ and radius
$h_c$, the mean value theorem gives
\[
 d_j([a,b])\subset d_j(m_c)+[-h_c,h_c]D_j,
\]
where $D_j$ encloses the explicit derivative on the whole cell and
\[
 d_0'(u)=-\sin u+ug(u),\qquad
 d_1'(u)=d_0'(u)+(r-1)u^2/2.
\]
Intersect this interval with the direct scalar enclosure. Both contain
the same function; no midpoint sampling replaces a remainder bound.
The minimum of several complete valid high-frequency bounds is allowed. Taking
selected favorable frequency pieces without proving their whole-cell
majorants is not.

\subsubsection{Turning one defect budget into simultaneous enclosures}
Lemma~\ref{ms:fourier:rescore:re-weight} gives the unweighted defect
budget after a first raw ceiling licenses its sign. Corollary~\ref{ms:conserved:phase-resource}
retains the additional phase curvature under its stated positivity gates.

The same identity also gives a useful score update. If this cell forces
$D\ge d_l\ge0$ and a preceding pass gives raw ceiling $r$, then
\[
 C\le 0.40974+
 \frac{r/N-0.3988(1-5r/(2N))d_l/6}{\rho_l}.
\]
A negative numerator excludes the cell because $e\ge0$. Otherwise its
numerator increases with $1/n$, and division by the positive lower
$\rho_l$ is outward. Thus the update is valid for every $n\ge N$ and
uses only the previous raw bound.

\paragraph{One positive-measure rule}
Suppose zero lies in the gap $(-A,B)$, $A,B>0$. Let
\[\begin{gathered}
 p=\Pp(X>0),\quad q=1-p,\quad m=\E X_+=\E X_-,\quad V=\E[X\mid\mathbf1_{X>0}],\\
 W=X-V,\quad\theta_\pm=\E[W^2;\pm X>0],\quad\theta=\theta_++\theta_-.
\end{gathered}\]
Then $m^2=(1-\theta)pq$, and the conditional means have magnitudes
$\bar A=m/q$, $\bar B=m/p$. Set
$\Delta_A=\bar A-A\ge0$, $\Delta_B=\bar B-B\ge0$.
We also use the unnormalised endpoint displacements
\[
 M_+=\E[(X-B)\mathbf1_{X>0}]=p\Delta_B,\qquad
 M_-=\E[(-X-A)\mathbf1_{X<0}]=q\Delta_A.
\]

\begin{lemma}[Positive-measure pricing]\label{ms:fourier:rescore:re-pricing}
For $x=|X|$, $a=|V|$, and a function $h$ constant on each sign class,
\begin{equation}\label{ms:fourier:rescore:re-conditional}
 \E[h(|X|^j-|V|^j)]
 =\E\left[hW^2\sum_{k=0}^{j-2}(k+1)a^kx^{j-2-k}\right],\quad j\ge2.
\end{equation}
Also, for $d\mu_0=(|X|-1)^2(|X|+1/2)\,dP$,
\begin{align*}
 \mu_0(\mathbb R)&=\rho-1,\\
 M_4-1&=\int\frac{(x+1)^2}{x+1/2}\,d\mu_0,\\
 M_5-1&=\int\frac{x^3+2x^2+3x+3/2}{x+1/2}\,d\mu_0,\\
 K_1+3\rho\beta&=\int
  \frac{\operatorname{sgn}(X)x(x+2-3\rho/2)}{x+1/2}\,d\mu_0.
\end{align*}
For any nonnegative resources $z_i$ satisfying $\sum a_i z_i\le r$,
$a_i>0$, a signed linear cost obeys
\begin{equation}\label{ms:fourier:rescore:re-simplex}
 \sum b_i z_i\le r\max\{0,b_1/a_1,\ldots,b_k/a_k\}.
\end{equation}
\end{lemma}
\begin{proof}
Divide $x^j-a^j-ja^{j-1}(x-a)$ by $(x-a)^2$; the linear term has
zero conditional mean. The radial identities follow by multiplying
out and using $\E X=0$, $\E X^2=1$. Finally each $b_i$ is at most
$a_i$ times the displayed maximum, and all $z_i$ are nonnegative.
\end{proof}

This one rule accounts for the moment, signed Taylor, endpoint movement,
outer-gap, and confinement-error estimates. In particular, no two costs
may each spend the full residual and then be added.

We use the central-gap notation and quantitative decomposition
\eqref{ms:defect-decomposition} from Section~\ref{ms:geometry-counts}.
In particular, $T_\pm$ measure the third-moment excess beyond the inner
endpoints and $J_o$ sums the outer gaps. The simultaneous inequality
\eqref{ms:simultaneous-gap-formula} supplies the common residual below.

\paragraph{Simultaneous bounds and their refinement}
For a closed probability interval $I=[l,u]$, suppose the following bounds hold:
\[
 |X|\le R,\quad C\in[c_0,c_1],\quad\mathcal R\le r,
 \quad nD\le Nd,\quad\theta\le\Theta<1,
 \quad A\ge A_l>0,\quad B\ge B_l>0,
\]
together with simultaneous signed-moment and coefficient intervals.
The distinction between $nD\le Nd$ and $D\le d$ is essential.
Let $s_+=\max_Ip(1-p)$, $s_-=\min_Ip(1-p)$ and
\begin{gather*}
 b_I=\frac{(6+2\sigma)(1-\Theta)^{3/2}}{\sqrt{s_+}}
                 \operatorname{dist}(p_*,I)^2,
 \qquad r_I=d-b_I,\\
 K_+=3(\sigma+1)B_l,\quad K_-=3(\sigma+3)A_l,
 \quad K=\min(K_+,K_-),\quad m_l=\sqrt{(1-\Theta)s_-}.
\end{gather*}
If $r_I<0$, exclude the entire closed cell. Otherwise \eqref{ms:simultaneous-gap-formula}
gives, simultaneously,
\begin{equation}\label{ms:fourier:rescore:re-residual}
 6m(A_l+B_l)(\Delta_A+\Delta_B)+K_+\theta_++K_-\theta_-\le r_I.
\end{equation}
The guaranteed inequalities for the actual law are
\begin{gather}
 \theta\le\min\{\Theta,r_I/K\},\notag\\
 A\ge\max\left\{A_l,\sqrt{(1-\Theta)l/(1-l)}-
                  \frac{r_I}{6m_l(A_l+B_l)}\right\},\notag\\
 B\ge\max\left\{B_l,\sqrt{(1-\Theta)(1-u)/u}-
                  \frac{r_I}{6m_l(A_l+B_l)}\right\}.
 \label{ms:fourier:rescore:re-geometry}
\end{gather}
The new variance ceiling is an outward upper bound for the first
right side, and each new endpoint floor is an outward lower bound for its
right side. They are not arbitrary numbers satisfying an inequality in
the opposite direction.
Each right side is evaluated from established bounds; sequential
improvements may then be used by the next expression. A nonpositive
endpoint estimate fails this update, rather than excluding a law.
Initialization uses the same formulas with the valid lower sum
$A+B\ge\sqrt{2-d}$.

In Bernoulli coordinates set $Y=B_0+U$, $\E[U\mid B_0]=0$,
$B_0\sim\mathrm{Bernoulli}(p)$, $X=\ell(Y-p)$ and
$d_Y=(\Delta_A+\Delta_B)/(\bar A+\bar B)$.
Then $\ell^2=(1-\theta)/(pq)$ and $\lambda=n\E U^2=npq\theta/(1-\theta)$.
The \emph{weighted} versions of \eqref{ms:fourier:rescore:re-residual} prove
\begin{equation}\label{ms:fourier:rescore:re-exports}
 n\theta\le\frac{Nr_I}{K},\quad
 nd_Y\le\frac{Nr_I}{6(1-\Theta)(A_l+B_l)},\quad
 \lambda\le\frac{Ns_+r_I}{(1-\Theta)K}.
\end{equation}
Indeed $nD\le Nd$ whereas $nb_I\ge Nb_I$.
These exports are derived from the weighted resource; multiplying an
unweighted variance ceiling by $N$ would not prove them.


For a joint bound on displacement and noise, apply
Theorem~\ref{ms:joint-resource} with
$d_D=d$, $v_*=s_+$, $S_*=A_l+B_l$, $K=\min(K_+,K_-)$ and $b=b_I$.
The premise follows from \eqref{ms:simultaneous-gap-formula} by dropping
$6m\Delta_A\Delta_B\ge0$; the required index-scaled input is $nD\le Nd_D$.
On a closed band $L_0\le\lambda\le U_0\le\Lambda$ this gives
$d_Y\le d_*(1-L_0/\Lambda)$.
Either positive intercept may first be enlarged, since both summands
are nonnegative. The theorem also handles zero and negative residuals.

For $j=3,4,5$ set
\[
 L_3(v)=R+2v,\quad L_4(v)=R^2+2Rv+3v^2,
 \quad L_5(v)=R^3+2R^2v+3Rv^2+4v^3.
\]
Put $v_+=\sqrt{(1-l)/l}$, $v_-=\sqrt{u/(1-u)}$ and
$k_j=[1-(1-\Theta)^{j/2}]/\Theta$, with $k_j=j/2$ at $\Theta=0$.
For $u<1/2$, valid Bernoulli moment ceilings are, with $s=l(1-l)$,
\[
 B_3^*=(1-2s)/\sqrt s,\quad B_4^*=1/s-3,
 \quad B_5^*=(1-4s+2s^2)/s^{3/2}.
\]
They decrease towards $p=1/2$. Convexity places $(1-\theta)^{j/2}$
below its endpoint chord on $[0,\Theta]$; positive-measure pricing gives
\begin{equation}\label{ms:fourier:rescore:re-moments}
 M_j\le B_j^*+r_I\max\left\{0,
 \frac{L_j(v_+)-k_jB_j^*}{K_+},
 \frac{L_j(v_-)-k_jB_j^*}{K_-}\right\}.
\end{equation}
For clarity, conditional Taylor expansion first gives
\[
 M_j\le(1-\theta)^{j/2}B_j^*
               +L_j(v_+)\theta_++L_j(v_-)\theta_-.
\]
The chord estimate subtracts $k_jB_j^*$ from each variance coefficient.
Multiplying the maximum in \eqref{ms:fourier:rescore:re-moments} by
$K_+\theta_++K_-\theta_-$, then using the single residual inequality,
proves \eqref{ms:fourier:rescore:re-moments}; a negative numerator is harmless because zero
is among the maximized values. For the outer charge, the deterministic
positive cubic term is
$m(\bar B^2-B^2)\le2m\bar B\Delta_B$.
Thus $T_+\le2mv_+\Delta_B+L_3(v_+)\theta_+$, with the reflected
formula on the negative class. Applying the same simplex rule bounds
\begin{equation}\label{ms:fourier:rescore:re-outerresource}
\begin{aligned}
 W_o&:=(\sigma-1)T_++(\sigma+1)T_-\\
 &\le r_I\max\left\{
 \frac{(\sigma-1)v_+}{3(A_l+B_l)},
 \frac{(\sigma+1)v_-}{3(A_l+B_l)},
 \frac{(\sigma-1)L_3(v_+)}{K_+},
 \frac{(\sigma+1)L_3(v_-)}{K_-}\right\}.
\end{aligned}
\end{equation}
Intersect every new bound with the previous one. The versions which
merge $\theta_+$ and $\theta_-$ are obtained by replacing both costs by
their maximum and both prices by their minimum; they need no new proof.

\paragraph{The enclosure theorem}\label{ms:fourier:rescore:re-operator}
On each cell, the formula atlas bounds the four coefficients in
\eqref{ms:fourier:rescore:re-raw}. The result, for each actual sign $s=\operatorname{sgn}H$,
is a finite family of quadratics
\begin{equation}\label{ms:fourier:rescore:re-quadratics}
 \mathcal R\le\max_j(A_j+F_{j,s}|H|-Q_j^{\rm raw}H^2),
 \qquad Q_j^{\rm raw}>0.
\end{equation}
The maxima retain all score--moment corners. In the initial signed
cover, intervals contain both signs of $\beta,\gamma$; no orientation is
discarded. In the probability cover, the central integral bound is a
maximum of affine functions of $\rho$, hence convex. The other
$\rho$-dependence is affine or the maximum of convex quadratics. For
fixed $H$, the score dependence is convex. Endpoints are therefore taken
\emph{before} completing the square.

\begin{theorem}[Finite resource enclosure]\label{ms:fourier:rescore:re-enclosure}
Consider a finite ordered schedule of the following operations:
\begin{enumerate}
\item Split a cell into an explicitly covered family of closed cells.
\item Intersect its enclosures with an inequality proved in the atlas.
\item Exclude a cell by an empty necessary interval or a strictly
negative nonnegative resource.
\item Apply a separately proved support or confinement implication,
checking all its hypotheses against this cell's already acquired data.
\item Export a raw, weighted, or signed phase bound by
\eqref{ms:fourier:rescore:re-quadratics}, \eqref{ms:fourier:rescore:re-phaseprice}, and \eqref{ms:fourier:rescore:re-exports}.
\end{enumerate}
If exact outward arithmetic checks every domain, corner and cover,
every output contains all attaining laws satisfying the input. The
conclusion holds for all integers $n\ge N$ wherever the formula's
monotonicity gates are checked. It includes zero conditional variance,
singleton intersections and arbitrarily small atoms.
\end{theorem}
\begin{proof}
Induct over the schedule. A split preserves the union, intersections use
necessary inequalities, and an excluded cell contains no feasible law.
Each support or confinement implication uses previously established hypotheses for that law.
For a scalar raw output, use
$fh-qh^2\le(f_+)^2/(4q)$ for $h\ge0$, $q>0$; using $f^2$ is a valid
weaker bound. For a resource output first add \eqref{ms:fourier:rescore:re-extra} to each
curvature. A negative maximum excludes the cell; a nonnegative one may
be divided by $\kappa N$ to give its next defect parameter.
The phase bound follows from the positive root of
$A+fh-qh^2\ge0$, replacing $A$ by $A_+$ if necessary.
A strict output is exported only when the certificate retains a strict
upper gap, such as its strict Fourier or target enclosure; closed-domain
coverage does not change strictness. All constants are frozen before evaluating
corners. Thus no iteration
step relies on an unproved fixed point. The stated index uniformity is
that of the displayed finite formulas, not a fixed-law limit.
\end{proof}

For the Fourier cancellation estimate only, reflect $(X,H)$ if needed:
$I_0(-X,-H)=I_0(X,H)$. Both odd coefficients change sign, and
$\int e^{-x^2/2}|zx^2-cx^4|\,dx$ is unchanged by $(z,c)\mapsto(-z,-c)$.
Consequently a lower endpoint used there bounds $|\gamma|$, not necessarily
the original signed $\gamma$. The actual contact, gap and phase data keep
their original orientation. A useful additional signed lower bound is
\begin{equation}\label{ms:fourier:rescore:re-gamma}
 \gamma\ge(1-\Theta)^{3/2}\frac{1-2u}{\sqrt{u(1-u)}}
       -\left(R+2\sqrt{\frac{u}{1-u}}\right)\Theta.
\end{equation}
To prove it expand $\E X^3=\E V^3+\E[W^2(X+2V)]$; the positive
residual is nonnegative and the negative residual has the displayed
lower bound. The binary term decreases in both $p<1/2$ and $\theta$.
Thus $\max(0,\text{any signed lower bound})$ is a valid lower bound for
$|\gamma|$ inside this auxiliary Fourier estimate.

\subsubsection{Coefficient atlas for the enclosure}\label{ms:fourier:rescore:re-atlas}

The atlas has four stages. First, signed first, second and third moments
bound the fourth and fifth moments and the gap integrals before a
probability panel is known. Second, an acquired central gap and its
conditional variances price the same expressions more sharply.
Third, signed phase bounds combine those moment prices with the complete
Fourier error into an affine-in-score quadratic in $|H|$. Finally,
probability and support feedback export one joint displacement/noise
budget to the count comparison.

The Fourier part is assembled before choosing these moment coordinates:
the low-frequency transfer treats the compensated Taylor terms, the
high-frequency formulas bound the entire complementary integral including
its infinite tail, and the contact identity supplies the signed
constant, linear and quadratic moment terms. The formulas below preserve
those separate contributions and state the gates for each alternative
bound. In particular, a lattice characteristic function is not assumed
to decay on the whole high-frequency half-line. The final output is the
quadratic family in \eqref{ms:fourier:rescore:re-quadratics}, with the
phase addition of \eqref{ms:fourier:rescore:re-extra} made once.

These formulas are inequalities for one actual law. They are the two
coordinate descriptions of the same enclosure: signed moments before a
probability gap is located, and conditional moments afterwards. A bound
is used only with its stated gates, and every refinement is an intersection.
\paragraph{Signed moments with a common first moment}
Put $a_\pm=(1\pm\beta)/2$, $b_\pm=(\rho\pm\gamma)/2$.
Both signs have positive probability and positive first moment, since a
centered variance-one law cannot be supported on one half-line. Thus
$|\beta|<1$ and the signwise denominators below are positive whenever
used; degenerate correction denominators are handled separately.
Cauchy gives $m^2\le a_+a_-$ and $b_\pm m\ge a_\pm^2$.
Consequently
\[
 \frac{(1+\beta)^{3/2}}{\sqrt{1-\beta}}-\rho
 \le\gamma\le
 \rho-\frac{(1-\beta)^{3/2}}{\sqrt{1+\beta}},\qquad
 m\ge\max\left\{\frac{1+\beta^2}{2\rho},
                  \frac{a_+^2}{b_+},\frac{a_-^2}{b_-}\right\}.
\]
These statements also hold with a zero atom. Projection of $X|X|$ onto
the span of $X$ and $|X|$ gives
\[
 \mu_4\ge\rho^2+\frac{(\gamma-\rho\beta)^2}{1-\beta^2}.
\]
For completeness, this follows by minimizing the nonnegative expectation
$\E(X|X|-uX-v|X|)^2$ over $u,v$; its Gram matrix has diagonal entries
one and off-diagonal entry $\beta$, and is positive definite.

Two support-localizing moment matrices give the additional upper bounds
as follows. If
$D_\pm=Ra_\pm-b_\pm>0$, then
\[
 2m+\sum_\pm\frac{(Rm-a_\pm)^2}{D_\pm}\le R,
 \qquad
 \mu_4\le R\rho-\sum_\pm\frac{D_\pm^2}{Rm-a_\pm}.
\]
Indeed the first follows by summing the determinant inequalities for
\[\E[(R-|X|)(1,|X|)^{\mathsf T}(1,|X|);\pm X>0].\]
The second follows from the corresponding matrices with an additional
factor $|X|$. One more factor $|X|$ yields, sign by sign,
\[
 \E[X_\pm^5]\le R \E[X_\pm^4]
 -\frac{(Rb_\pm-\E[X_\pm^4])^2}{D_\pm}.
\]
Zero denominators are treated by dropping their nonnegative correction;
they never imply exclusion without a separate contradictory numerator.
The common $m$ is retained throughout.

A second useful upper bound for $m$ follows by integrating
$(x-a)^2(R-x)\ge0$ for $x=|X|\in[0,R]$. For every fixed $a>0$,
\[
 2m\le\frac{R-\rho+2a+Ra^2}{a(2R+a)}.
\]
The cover uses only the eleven rational choices $a=2/10,\ldots,12/10$.

Write $\delta=\gamma-\rho\beta$. Cauchy and a square completion on
each sign give the signed constant bound
\begin{align*}
 \frac\rho2\mu_4-\frac14\eta_5
 &\le\frac{\rho^3}{4}
 -\frac{(\rho+\gamma)(\delta_+)^2}{8(1+\beta)^2}
 -\frac{(\rho-\gamma)((-\delta)_+)^2}{8(1-\beta)^2}.
\end{align*}
Indeed on either sign, write $a=\E X_s^2$, $b=\E X_s^3$,
$c=\E X_s^4$, $f=\E X_s^5$. Cauchy gives $c\ge b^2/a$,
$f\ge c^2/b$, and
$\rho c/2-f/4\le\rho^2b/4-(c-\rho b)^2/(4b)$.
Distance to the half-line $c\ge b^2/a$ proves the stated penalty
on summing both signs. Independently,
$\mu_4^2\le\rho\eta_5$ and the projection bound above give
\[
 \frac\rho2\mu_4-\frac14\eta_5
 \le\frac{\rho^3}{4}
 -\frac{(\gamma-\rho\beta)^4}{4\rho(1-\beta^2)^2}.
\]
The larger of the two nonnegative deductions is retained.
There is a third independent deduction which is often stronger.
For $Y=|X|$ and $\delta=\gamma-\rho\beta$, put
\[
 V=\E[Y(Y^2-\rho Y)^2]=\eta_5-2\rho\mu_4+\rho^3.
\]
Since
$\delta=\E[Y\operatorname{sgn}(X)(Y^2-\rho Y)]$,
weighted Cauchy gives $\delta^2\le2mV$. Hence
\begin{equation}\label{ms:fourier:rescore:eq:weighted-constant-loss}
 \frac\rho2\mu_4-\frac14\eta_5
 \le\frac{\rho^3}{4}-\frac{\delta^2}{8m}.
\end{equation}
The maximum of these three separately valid losses can be subtracted.
Their sum cannot. A zero atom contributes zero to the weighted measure.

\paragraph{An actual central gap without assuming $D<2$}
The already proved gap decomposition \eqref{ms:defect-decomposition} shows that
$D<(\sigma-1)\rho$ implies an actual central gap. Use its same
nonnegative terms, writing $T=T_++T_-$ and
$Q=(\sqrt{\sigma+1}A-\sqrt{\sigma-1}B)^2$.
Hence $T\le D/(\sigma-1)$ and
\[
 (A+B)^2\ge A^2+B^2\ge\frac{\rho-D/(\sigma-1)}m.
\]
Moreover $A\le a_-/m$, $B\le a_+/m$ and $J\le\rho+2m$.
The last bound gives the compulsory defect
\[
 D\ge \sigma\rho-\gamma-6m.
\]
The inequalities $T-J_o\ge2BR_++2AR_-$ give the useful additional
endpoint constraint
\[
 D\ge(\sigma-1)b_++6a_+B-(\sigma+5)mB^2.
\]
If $D<(\sigma-1)b_+$, this downward quadratic forces $B$ to be at
least its positive root. The reflected constraint has $a_-,b_-,A$
and the coefficients $\sigma+1,\sigma+7$. In a closed moment box,
positive coefficients are rounded down, negative coefficients up, and
the lower square-root endpoint is used in the resulting root bound.

The actual opposite-sign event has probability $2pq$.
The first two moment constraints give lower bounds for both $p$ and $q$.
For independent copies one has exactly
\[
 \E[(X-X')^2\mid XX'<0]
 =\frac{1+\beta(1-2p)+4m^2}{2pq}
 =\frac{1+4m^2+\beta}{2p}
  +\frac{1+4m^2-\beta}{2q}.
\]
For an upper bound, use the upper $m$ endpoint and both endpoints of
$p$ and $\beta$. The expression is convex in $p$ and affine in $\beta$.
This improves the former $2/\Pp(XX'<0)$ bound inside the complete high-frequency
estimate. Both copies still have the same actual marginal.

The gap can also be licensed before the sufficient inequality
$D<(\sigma-1)\rho$ is available. If an already proved bound $D\le d$
and a signed moment cell give
\[
 P_0=\max\{(\sigma_l-1)\rho_l,\sigma_l\rho_l-\gamma_u\}-d>0,
 \qquad \sigma_l\le\sqrt{10},
\]
then zero cannot belong to support: in that case
$D\ge\sigma\rho-\gamma\ge P_0+d$.
For the resulting actual gap, $J\le\rho+2mAB$ gives
$6mAB\ge P_0$. Thus a closed panel $p\in[l,u]\subset(0,1)$ has
\begin{equation}\label{ms:fourier:rescore:eq:general-endpoint-floors}
 A\ge\frac{P_0}{6(1-l)},\qquad B\ge\frac{P_0}{6u},\qquad
 (A+B)^2\ge\frac{2P_0}{3m_u}.
\end{equation}
Indeed $A\le m/(1-p)$ and $B\le m/p$, while
$m^2\le p(1-p)$, so $mAB\le(1-p)A$ and $mAB\le pB$.
The last inequality uses $4AB\le(A+B)^2$.
If earlier steps have acquired $A\ge a>0$, $B\ge b>0$ and $m\ge m_l$,
one may replace the lower price $P_0$ by
\begin{equation}\label{ms:fourier:rescore:eq:retained-endpoint-product}
 P=\max\{P_0,6m_l a b\}\le6mAB.
\end{equation}
This retains an already proved physical gap; it does not license a gap
from an unsuccessful initialization. In fact the new term alone gives
$2(6m_l a b)/(3m_u)\le4ab\le(a+b)^2$.

\paragraph{Preserving central and outer integral costs}
Let
\[
 P_1(A,B)=B^3/3+AB^2/2-A^3/6,\quad
 P_2(A,B)=B^4/4+AB^3/3-A^4/12,\quad G=\rho P_1-P_2.
\]
Then $J_{1,c}=mP_1$, $J_{2,c}=mP_2$. Since
$\partial_BG=B(A+B)(\rho-B)$,
\[
 G(A,B;\rho)\le G(A,\rho;\rho)
 =\rho^4/12+\rho^3A/6-\rho A^3/6+A^4/12.
\]
The last expression increases with $A$ because its derivative is
$(\rho-A)^2(\rho+2A)/6$, and is convex in $\rho$ because its second
$\rho$ derivative is $\rho^2+\rho A$. Its closed-rectangle maximum is
therefore bounded using the upper $A$ endpoint and both $\rho$ endpoints.

Positive outer gaps satisfy
$(t-a)(\rho-t)\le(\rho-B)_+^2/4$. On negative outer gaps the positive
part is bounded by $(R-\rho)^2/4$ after extracting $|t|$.
Integration gives
\[
 3(\rho J_1-J_2)_o
 \le\frac38\{(\rho-B)_+^2T_++(R-\rho)^2T_-\}.
\]
The available weighted outer budget is $D-mQ$, not the entire $D$.
The cover intersects this budget with the individual bounds
$T_+\le b_{+,u}-m_lB_l^2$ and $T_-\le b_{-,u}-m_lA_l^2$.

For the signed phase term $U=J_1+(\mu_4-3)/6$, retain
\[
 -RT_-/2\le J_{1,o}\le RT_+/2.
\]
The central minimum of $P_1$ occurs at $B_l$ and an $A$ endpoint;
its maximum occurs at $B_u$ and $A$ clipped to the interval at $A=B_u$.
A nonnegative polynomial bound is multiplied by $m_u$ for an upper
bound and $m_l$ for a lower bound; a negative one uses the opposite
endpoint. For example, the binary endpoints $A=3/2$, $B=2/3$ make this
maximum negative; the sign switch is necessary even when $D<2$.
The fourth-moment sign intervals also sharpen
$K_1=\E[X|X|^3]-3\beta-3\rho\gamma/2$.
The complete $K_1,U,\phi$ rectangles are used in the signed expression
$|CK_1+\phi U|$, without replacing it by two separate absolute values.


\paragraph{A complete coarse bound before a central gap is licensed.}
Always $-\mu_4/6\le J_1\le\mu_4/2$ and
$-3J_2-\eta_5/4\le0$. On a gap of length $h=b-a$, these follow from
\[
 \int_a^b(t-a)t\,dt+\tfrac12\int_a^bt^2dt=hb^2/2,\quad
 \tfrac32\int_a^bt^2dt-\int_a^b(t-a)t\,dt=h(2a+b)^2/6,
\]
and, on a negative gap $(-v,-u)$, from
$\tfrac13\int_u^v x^3dx-\int_u^v(v-x)x^2dx=u^3(v-u)/3$.
For $1\le r\le\rho$ and actual support radius $R\le6$, set
\[
 \Gamma(R,r)\ge\max_{r\le v\le R}
 \left\{-\frac{v^2}4+\frac{(R+r)v}3-\frac{Rr}2+
                      \frac{r^3(2R-r)}{12v^2}\right\}_+.
\]
Then
\[
 3(\rho J_1-J_2)\le\frac32\Gamma(R,r)(\rho-\gamma)
                              +\frac38\rho^2J.
\]
On negative gaps extract $x=|t|$ and discard the nonpositive part
$x<\rho$. Replacing $\rho$ by $r$ and the outer endpoint by $R$
gives $\int_r^v(R-x)x(x-r)\,dx$. Dividing this integral by $v^2$
gives exactly the displayed expression. Signwise Tonelli now charges
$\Gamma\E X_-^3=\Gamma(\rho-\gamma)/2$.
On the positive portion of each gap, $t(\rho-t)\le\rho^2/4$.
Multiply by $t-a\ge0$ and integrate; since
$\int_a^b(t-a)\,dt=(b-a)^2/2$, the resulting contribution is
at most $\rho^2J/8$. This includes a gap crossing zero.
Partition the $v$-range into 64 closed intervals: on each, maximize the
concave quadratic at its clipped vertex $2(R+r)/3$ and the positive
reciprocal term at the left endpoint. This proves its interval bound.
The alternative bound $\Gamma\le5/2$ follows at $r=1,R=6$ by
$3v^4-28v^3+66v^2-11>0$ on $[1,6]$: the endpoint values are $30,205$,
and the interior critical values are $(85\mp35\sqrt5)/2>0$.
These monotone enclosing parameters cover all $r\ge1,R\le6$.
\paragraph{Moment and gap coefficients for a whole cell}
\label{ms:fourier:rescore:mom-conditional-boxes}

The following coefficient rules are collected here to make the finite
calculation transparent. Every rule is an inequality for an actual law;
there is no extremal-law assumption inside a cell.
Write $v=1-\theta$. For $p\in[l,u]\subset(0,1)$ set
\begin{gather*}
 s_0=\min\{l(1-l),u(1-u)\},\quad
 v_* =\max\{\sqrt{u/(1-u)},\sqrt{(1-l)/l}\},\\
 r_B=(1-2s_0)/\sqrt{s_0},\quad
 g_B=\max_{t=l,u}|1-2t|/\sqrt{t(1-t)},\quad b_B=\max_{t=l,u}|1-2t|,\\
 L_3=R+2v_*,\quad L_4=R^2+2Rv_*+3v_*^2,\\
 L_5=R^3+2R^2v_*+3Rv_*^2+4v_*^3,\quad
 e_B=r_B\frac{1-4s_0+2s_0^2}{s_0(1-2s_0)}.
\end{gather*}
The polynomials $T_3=x+2a$, $T_4=x^2+2ax+3a^2$ and
$T_5=x^3+2ax^2+3a^2x+4a^3$ are nonnegative. Their maxima at
$x=R$, $a=v_*$, in Lemma~\ref{ms:fourier:rescore:re-pricing}, give, if $\theta\le\Theta<1$,
\begin{align}
 \rho&\le r_B+(L_3-r_B)\Theta,&
 |\gamma|&\le g_B+(L_3-g_B)\Theta,\notag\\
 |\beta|&\le b_B+(1-b_B)\Theta,&
 \mu_4&\le(1-\Theta)^2(1/s_0-3)+L_4\Theta,\notag\\
 \eta_5&\le(1-\Theta)^2e_B+L_5\Theta.&&\label{ms:fourier:rescore:mom-boxes}
\end{align}
Replacing $\theta$ by $\Theta$ uses the explicitly checked inequalities
$L_3>r_B,g_B$, $L_4>2(1/s_0-3)$ and $L_5>2e_B$.
The exact binary factors are $v^{3/2},v,v^2,v^{5/2}$; the upper bounds
use $v^{3/2}\le v$ and $v^{5/2}\le v^2$. Thus cells crossing $1/2$ and
zero conditional variances are included. The same cubic expansion, in the
later domain $p\in[0.37,0.47]$, $A\ge0.75$, $B\ge1.05$, gives
\begin{equation}\label{ms:fourier:rescore:mom-rho-lower}
 \rho\ge v^{3/2}\rho_B(p)+3\min(A,B)\theta
       \ge\rho_B(p)+0.45\theta,\qquad
 \rho_B(p)=\frac{1-2pq}{\sqrt{pq}}.
\end{equation}
Indeed $\rho_B<1.11$ there and $v^{3/2}\ge1-3\theta/2$.
Also $\beta\ge1-2p-2(1-p)\theta$.

For the signed coefficient
$K_1=\E[X|X|^3]-3\beta-3\rho\gamma/2$, define
\[
 K(p,v)=(1-2p)\{(3v^3/2-v^2)/(pq)-3v^3+v^2+3v\}.
\]
Direct substitution gives $K_1(V)=-K(p,v)$, with the coefficient $3$
unchanged. For $v\ge2/3$, $K$ decreases in $p$: both coefficients in
braces are nonnegative and $(1-2p)/(pq)$ decreases. Its derivative in $v$
has the sign of $1-2p$, since the derivative of the braces is at least
$9v^2-6v+3>0$. Hence both variance endpoints are retained. The error is
\begin{equation}\label{ms:fourier:rescore:mom-K-error}
 |K_1(X)-K_1(V)|\le
 \Theta\{L_4+3+\tfrac32L_3(R_\rho+g_B)\}
\end{equation}
when $\rho\le R_\rho$. This follows from the conditional expansions and
$\rho\gamma-\rho_V\gamma_V=
\rho(\gamma-\gamma_V)+\gamma_V(\rho-\rho_V)$.

Here are the remaining signed charges. They also explain the two
alternative bounds retained in the initial, coarse moment cover:
\begin{align}
 -\mu_4/6&\le J_1\le\mu_4/2,&
 -3J_2-\eta_5/4&\le0,                                  \label{ms:fourier:rescore:mom-gap-signs}\\
 \rho J_1-J_2&\le\tfrac54(\rho-\gamma)+\tfrac18\rho^2(\rho+1),
                                                              \label{ms:fourier:rescore:mom-oldgap}\\
 -75D/104&\le J_{1,o}\le3D/2.                           \label{ms:fourier:rescore:mom-outer}
\end{align}
The first two lines follow from the coarse gap identities proved above.
For~\eqref{ms:fourier:rescore:mom-outer}, use $J_{1,o}\ge-3T_-$,
$J_{1,o}\le3J_o\le3T$ and~\eqref{ms:defect-decomposition}.

The central integrals are $J_{j,c}=mP_j(A,B)$, with
\[
 P_1=B^3/3+AB^2/2-A^3/6,\qquad
 P_2=B^4/4+AB^3/3-A^4/12.
\]
They increase in $B$; their $A$ derivative changes from positive to negative
at $A=B$. A rectangle minimum is therefore at $B_l$ and an $A$ endpoint.
Multiply a negative minimum by an upper bound for $m$, and a nonnegative
minimum by a lower bound. For the upper central charge, maximizing at
$(\bar a,\bar b)$ if $p\le1/2$ gives
$J_{1,c}\le v^2(2-3p)/(6pq)$, a decreasing function of $p$ on that interval.
If $p\ge1/2$, maximize first at $A=B$ to get
$J_{1,c}\le(2/3)v^2q^2/p\le1/3$. Thus every cell may use
\begin{equation}\label{ms:fourier:rescore:mom-central-upper}
 J_{1,c}\le\max\{1/3,(2-3l)/(6l(1-l))\}.
\end{equation}

For $|X|\le R$, $R\ge1$, define
\[
 k_4(R)=\frac{(R+1)^2}{R+1/2},\qquad
 k_5(R)=\frac{R^3+2R^2+3R+3/2}{R+1/2}.
\]
Then
\begin{equation}\label{ms:fourier:rescore:mom-universal-moments}
 \mu_4\le\min\{R\rho,1+k_4(R)(\rho-1)\},\quad
 \eta_5\le\min\{R\mu_4,1+k_5(R)(\rho-1)\},\quad
 \mu_6\le R\eta_5.
\end{equation}
Writing $u=|X|$ in~\eqref{ms:fourier:rescore:re-pricing}, the first two constants
are the endpoint values of the ratios
$u+3/2+1/[4(u+1/2)]$ and
$u^2+3u/2+9/4+3/[8(u+1/2)]$, whose derivatives are nonnegative.
The other bounds are pointwise. The following coarse cancellation estimates
continue to use the valid radius-six majorants:
The last identity in~\eqref{ms:fourier:rescore:re-pricing} gives
\begin{equation}\label{ms:fourier:rescore:mom-K-cancellation}
 |K_1+3\rho\beta|\le(96-18\rho)(\rho-1)/13.
\end{equation}
Indeed $u(u+2-3r/2)/(u+1/2)$ is convex on $[0,6]$, its maximum is
$(96-18r)/13$, and its negative part is at most $31/40$, smaller than that
maximum for $1\le r\le37/20$.
The other bounds needed below are
\begin{align}
 |K_1|&\le6\rho-3\rho^2/2+3|\beta|,\notag\\
 K_1&\le M+3|\beta|-3\rho g_-/2,&
 -K_1&\le3\rho+3|\beta|+(3\rho/2-3)g_-,\notag\\
 -K_q&\le(\rho/2-3/8)M-\rho/2+1/8+3|\beta|g/2,&
 -K_q&\le8\rho^3/27+(3|\beta|/2-1/2)\rho,               \label{ms:fourier:rescore:mom-K-bounds}
\end{align}
where $\mu_4\le M$, $|\gamma|\le g$, $\gamma\ge g_-$, and
$K_q=\eta_5/4-\rho\mu_4/2-3\beta\gamma/2+\rho/2$.
For the first two lines use
$|\E[X|X|^3]-3\gamma|\le3\rho$, and note the negative coefficients
of $\gamma$. For the last line use
$\eta_5\ge\mu_4^{3/2}\ge3\mu_4/2-1/2$, or maximize
$\rho y^2/2-y^3/4$ over $y\ge0$.


\paragraph{Compatible weaker resource constants.}
One may also use the weaker binary charge
$b_I=(616/25)(1-\Theta)^2\operatorname{dist}(p_*,I)^2$.
It follows from $2(6+2\sigma)>616/25$ and
$(1-\theta)^{3/2}\ge(1-\theta)^2$.
With $K=\min(K_+,K_-)$ and $v_* =\max(v_+,v_-)$, the merged charge
\eqref{ms:fourier:rescore:re-residual} gives
\[
 T_++T_-\le T_*=(d-b_I)\max\{v_*/[3(A_l+B_l)],(R+2v_*)/K\},
 \qquad d_o=\min\{d,(\sigma+1)T_*\}.
\]
Thus $W_o\le d_o$ and, for $R\le6$,
\[
 -75d_o/104\le J_{1,o}\le3d_o/2,\qquad
 3(\rho J_{1,o}-J_{2,o})\le\frac38
 \max\{\rho^2/(\sigma-1),(R-\rho)^2/(\sigma+1)\}d_o.
\]
These are precisely the merged-cost versions of the sign-specific
resource rule; they do not introduce another available budget.

\paragraph{Conditional lower moments improve two signed charges}

Let $s_u=\max_{p\in[l,u]}pq$ and $\theta\le\Theta$.
The unsigned cases of~\eqref{ms:fourier:rescore:re-conditional} give
\begin{align}
 \mu_4&\ge M_l:=\max\{1,(1-\Theta)^2(1/s_u-3)\},\label{ms:fourier:rescore:entry201-m4lower}\\
 \eta_5&\ge E_l:=
 (1-\Theta)^{5/2}\frac{1-4s_u+2s_u^2}{s_u^{3/2}}.\label{ms:fourier:rescore:entry201-e5lower}
\end{align}
Both binary moment functions decrease in $s\in(0,1/4]$.
For the fifth this follows because the derivative has the sign of
$s^2+2s-3/2<0$; the fourth is immediate.
The fourth-moment lower bound sharpens the lower endpoint of
$J_1+(\mu_4-3)/6$; its explicit use appears in the cell map below.

The constant signed charge is
\[
 -K_q=\frac{\rho\mu_4}{2}-\frac{\eta_5}{4}
       +\frac32\beta\gamma-\frac{\rho}{2}.
\]
One coherent estimate uses $\eta_5\ge3\mu_4/2-1/2$ and
$\mu_4\le M$. A second estimate uses \eqref{ms:fourier:rescore:entry201-e5lower}.
Their difference is independent of $\rho$. Consequently the common
bound improves by subtracting
\begin{equation}\label{ms:fourier:rescore:entry201-eta-improvement}
 \frac14\max\{0,E_l-3M/2+1/2\}.
\end{equation}
This is a constant shift on the cell and does not spoil convexity.
It does not assert the generally invalid implication
$\eta_5\ge3M/2-1/2$ from an upper bound on $\mu_4$.

\paragraph{Signed conditional Taylor cancellation}

Let $K_1=\E[X|X|^3]-3\beta-\tfrac32\rho\gamma$.
On a sign class put $x=|X|$, $a=|V|$, $W=X-V$ and
$\varepsilon=1$ on the positive class, $-1$ on the negative class.
In the actual central-gap domain there is no atom at zero.
With $w_3=x+2a$, $w_4=x^2+2ax+3a^2$ and
$\gamma_V=\E V^3$, exact conditional expansion gives
\begin{equation}\label{ms:fourier:rescore:entry201-Kexact}
 K_1(X)-K_1(V)
 =\E\!\left[W^2
 \left\{\varepsilon(w_4-3-\tfrac32\rho w_3)
                   -\tfrac32\gamma_Vw_3\right\}\right].
\end{equation}
Apply~\eqref{ms:fourier:rescore:re-conditional} to $j=2,3,4$, with $h=1$
or $h=\varepsilon$, and use
$\rho\gamma-\rho_V\gamma_V
=\rho(\gamma-\gamma_V)+\gamma_V(\rho-\rho_V)$.

The positive coefficient in \eqref{ms:fourier:rescore:entry201-Kexact} is
$F(x,a;t)=x^2+2ax+3a^2-3-\tfrac32t(x+2a)$
with $t=\rho+\gamma_V$; the negative coefficient is $-F$
with $t=\rho-\gamma_V$. Their ranges are bounded on complete
rectangles of actual conditional means and support magnitudes.
The positive-support magnitude is at most $R$.
The same formula permits separate actual bounds on the two signed support magnitudes.
On the late domain $p<1/2$, bounds for $\gamma_V$ are
\[
 \sqrt{(1-\Theta)^3(1-2u)^2/[u(1-u)]}
 \le\gamma_V\le(1-2l)/\sqrt{l(1-l)}.
\]
The positive mean lies in
$[\sqrt{(1-\Theta)(1-u)/u},\sqrt{(1-l)/l}]$, and the negative
magnitude mean lies in
$[\sqrt{(1-\Theta)l/(1-l)},\sqrt{u/(1-u)}]$.
A valid common lower bound for the actual absolute third moment is
\begin{equation}\label{ms:fourier:rescore:relaxed-rho}
 r_*:=\max\{1,r_--(3r_+/2-3a_l)_+\Theta\}\le\rho,
 \quad r_-\le\rho_B(u),\quad r_+\ge\rho_B(l),
 \quad a_l\le\min(A,B).
\end{equation}
Indeed conditional cubic expansion gives
$\rho\ge(1-\theta)^{3/2}\rho_B(p)+3a_l\theta$.
Convexity gives $(1-\theta)^{3/2}\ge1-3\theta/2$, and $\rho_B$
decreases on $(0,1/2)$, so
$\rho\ge r_--(3r_+/2-3a_l)_+\Theta$; also $\rho\ge1$.
The stronger bound $\rho\ge\rho_B(u)$ is recovered if
$3a_l\ge3r_+/2$. Its reserve condition is sufficient, not necessary.
Use the same $r_*$ in the paired Fourier rectangle, signed Taylor
coefficients, phase curvature and exported common-law quadratics.
Recomputing a larger lower endpoint in one of those locations would
invalidate their association.

Since $x+2a>0$, $F$ decreases with $t$.
For fixed $t$ its Hessian is
$\left(\begin{smallmatrix}2&2\\2&6\end{smallmatrix}\right)>0$.
Its rectangle maximum is at a corner. Its minimum is among the four
corners, the minimizer on each edge, and the stationary point
$(x,a)=(3t/8,3t/8)$ if that lies in the rectangle.
On fixed-$x$ edges use $a=t/2-x/3$, clipped to the edge; on
fixed-$a$ edges use $x=3t/4-a$, also clipped.
This is a complete finite extremum list.
Multiplication by $0\le\theta\le\Theta$ uses
$\Theta\min(0,L)$ and $\Theta\max(0,U)$ for the lower and upper
coefficient bounds. It intersects this new interval with the retained
absolute-error interval; neither need be discarded.

\paragraph{Generic event and complete central rectangle}
\begin{lemma}[Generic small-defect event]\label{ms:fourier:rescore:entry201-event}\label{ms:fourier:rescore:mom-third}
Assume $|X|\le R\le6$ and $\rho\le37/20$. Let $0\le d<2$ and $D\le d$. Define
\[
\begin{gathered}
 r_*=\min\{37/20,(3+d)/(\sigma-1)\},\qquad W=\sqrt{37d/10},\\
 c=\sqrt{\sigma+1},\qquad e=\sqrt{\sigma-1},\qquad C_0=c+e,
\end{gathered}
\]
\[
 h_*=\frac{C_0}{6}\sqrt{12+W^2}+\frac{W}{3C_0},\qquad
 L_*=\max\left\{h_*,\frac{(2c-e)h_*+3W}{C_0}\right\}.
\]
Then $\rho\le r_*$, $A+B\le h_*$ and $|2B-A|\le L_*$.
If $\alpha=L_*/6\le3$, then
\[
 |J_1|\le\alpha(\rho+1)+(3-\alpha)d/2,\quad
 \mu_4\le h_*\rho+Rd/2,\quad
 \eta_5\le h_*^2\rho+R^2d/2.
\]
For independent copies, the actual event $E=\{XX'<0\}$ obeys
\[
 \Pp(E)\ge a:=1/(2r_*^2),\quad
 |X-X'|\ge d_0:=\sqrt{2-d}\text{ on }E,\quad
 \E[(X-X')^2\mid E]\le2/a.
\]
\end{lemma}
\begin{proof}
The inequality $J\le\rho+1$ gives $D\ge(\sigma-1)\rho-3$.
Equation \eqref{ms:defect-decomposition} gives $w^2\le2\rho D\le37d/10$ for
$w=cA-eB$. Since $ce=3$,
\[
 A+B=\frac{C_0}{6}\sqrt{w^2+12AB}-\frac{w}{3C_0},
 \qquad 2B-A=\frac{(2c-e)(A+B)-3w}{C_0}.
\]
Use $AB\le1$, $|w|\le W$; also $2B-A\ge-(A+B)$.
Now $J_{1,c}=J_c(2B-A)/6$, $|J_{1,o}|\le3J_o$ and
$0\le J_o\le T\le d/2$ imply
$|J_1|\le\alpha J+(3-\alpha)J_o$.

For either inner endpoint $b_*$ and $b_*\le x\le R$,
\[
 x^4-h_*x^3\le R x(x^2-b_*^2),\qquad
 x^5-h_*^2x^3\le R^2x(x^2-b_*^2).
\]
Indeed $h_*\ge b_*$, and after factoring $x-b_*$ the remaining
comparisons follow from $x\le R$. Integrate and use $T\le d/2$.
Finally $m^2\le pq$ and $m\ge1/(2\rho)$ give $2pq\ge a$;
the conditional squared distance is bounded by its unconditional
expectation $2$, divided by the event probability.
\end{proof}

The central part of $\rho J_1-J_2$ is exactly
$mG(A,B;\rho)$, where $G=\rho P_1-P_2$.
Let a proved probability cell give
$A\in[A_l,A_u]$, $B\in[B_l,B_u]$, $m\in[m_l,m_u]$.
Here
\[
 A_u=\sqrt{u/(1-u)},\quad B_u=\sqrt{(1-l)/l},\quad
 m_l=\sqrt{(1-\Theta)s_0},\quad
 m_u=\sqrt{\max_{p\in[l,u]}pq}.
\]

\paragraph{The closed-box Bernstein identity.}\label{ms:fourier:rescore:re-bernstein}
The Bernstein identity in Section~\ref{ms:certificates} applies after an
affine change from each closed coordinate interval to $[0,1]$.
Its tensor-product form includes every face of the resulting closed box.
We use that identity with coordinate degrees four below.

\begin{lemma}[Complete central rectangle]\label{ms:fourier:rescore:entry201-central}
Write $G(A_l+(A_u-A_l)x,B_l+(B_u-B_l)y;\rho)
=\sum t_{kh}x^ky^h$. Define the 25 rational coefficients
\[
 b_{ij}(\rho)=\sum_{\substack{k\le i\\h\le j}}
 t_{kh}\frac{\binom{i}{k}}{\binom4k}
          \frac{\binom{j}{h}}{\binom4h},\qquad 0\le i,j\le4,
\]
and
\[
 \mathcal G(\rho)=\max_{i,j}
       \{m_l b_{ij}(\rho),\,m_u b_{ij}(\rho)\}.
\]
Then $\rho J_{1,c}-J_{2,c}\le\mathcal G(\rho)$.
Moreover $\mathcal G$ is convex in $\rho$.
\end{lemma}
\begin{proof}
Apply the Bernstein identity~\ref{ms:fourier:rescore:re-bernstein} separately to the
two coordinates after the affine substitution. The Bernstein basis is nonnegative and sums to one on the closed square,
so $G$ is a convex combination of its 25 coefficients.
For each coefficient retain both $m$ endpoints, or equivalently use
$m_u$ when it is nonnegative and $m_l$ otherwise. This gives a uniform
bound on the entire three-dimensional rectangle, including boundaries.
Each endpoint-multiplied coefficient is affine in $\rho$; their maximum
is convex. No interior stationary point is omitted.
\end{proof}

\paragraph{Refining a rectangular parameter enclosure}
\label{ms:fourier:rescore:entry-feedback}\label{ms:fourier:rescore:fourier-uniform}

The raw remainder in \eqref{ms:fourier:rescore:re-raw} is
\begin{equation}\label{ms:fourier:rescore:entry-remainder}
 \mathcal R=C\{3\beta J/2+3\rho J_1-3J_2-K_q\}
 +H(CK_1+\varphi U)-H^2(C\rho+\gamma\varphi/2)+E_n,
\end{equation}
where $U=J_1+(\mu_4-3)/6$, $\varphi=\phi(H/\sqrt n)$.
For any one coherent low-frequency triple and complete high-frequency value $K_H$ at floor $N$,
\begin{equation}\label{ms:fourier:rescore:entry-fourier-budget}
 |E_n|\le K_F+s_F|H|,\quad
 K_F=c_\pi(L_0+K_H+L_2/\sqrt N),\quad s_F=c_\pi L_H,\quad c_\pi\ge2/\pi.
\end{equation}
Both coefficient choices, $16/25$ before the signed phase and $212/333$
afterward, are valid since $\pi>333/106$.

\begin{lemma}[Continuous cells and feedback]\label{ms:fourier:rescore:entry-cell-theorem}\label{ms:fourier:rescore:entry-corners}
Suppose bounds for one actual cell give
\[
 r_0\le\rho\le r_1,\quad0<c_0\le C\le c_1,\quad
 |\beta|\le b,\quad |\gamma|\le g,\quad \mu_4\le M,\quad D\le d,
\]
$K_1\in[k_l,k_u]$, $U\in[u_l,u_u]$, and
$0\le\varphi_-\le\varphi\le2/5$. Freeze the shared outer charge $d_o$,
lower fifth moment $E_l$, upper $s_u=\max_{p\in[l,u]}p(1-p)$, and
central envelope $\mathcal G$ from the preceding bounds.
Define
\begin{align*}
 A(r)&=\tfrac32 b(r+2\sqrt{s_u})+3\mathcal G(r)
 +\tfrac38\max\{r^2/(\sigma-1),(R-r)^2/(\sigma+1)\}d_o\\
 &\quad +(r/2-3/8)M-r/2+1/8+\tfrac32bg
 -\tfrac14\max\{0,E_l-3M/2+1/2\},\\
 B(c)&=\max_{k\in\{k_l,k_u\},\,u\in\{u_l,u_u\},\,
                      \psi\in\{\varphi_-,2/5\}}|ck+\psi u|,\\
 Q(c,r)&=cr+\tfrac12\varphi_-(3.16r-3-d)
       -\tfrac12(2/5-\varphi_-)(d-0.16)_+.
\end{align*}
If all four corner values of $Q$ are positive, then, for every real $H$,
\begin{align}
 \mathcal R&\le\max_{(c,r)\in\{c_0,c_1\}\times\{r_0,r_1\}}
       \{K_F+cA(r)+(B(c)+s_F)|H|-Q(c,r)H^2\}\notag\\
 &\le\max_{c,r\text{ corners}}
       \{K_F+cA(r)+(B(c)+s_F)^2/(4Q(c,r))\}.           \label{ms:fourier:rescore:entry-four-corners}
\end{align}
If $T$ is a strict upper bound for the final maximum in
\eqref{ms:fourier:rescore:entry-four-corners}, the weight-removal step in
Lemma~\ref{ms:fourier:rescore:re-weight} gives
\begin{equation}\label{ms:fourier:rescore:entry-feedback-values}
 D<\frac{15000T}{997N},\quad
 \varphi>\frac{997}{2500}\left(1-\frac{5T}{2N}\right),\quad
 C<\CE+V/N\quad\text{if also }\mathcal R/\rho<V.
\end{equation}
Only positive density floors are used. These statements remain valid at
$H=0$, on closed cell edges, and for zero conditional variance.
\end{lemma}
\begin{proof}
The constant coefficient follows from the shared outer budget,
Lemma~\ref{ms:fourier:rescore:entry201-central}, and~\eqref{ms:fourier:rescore:mom-K-bounds}, using
\[
 J\le\rho+2ABm\le\rho+2(1-\theta)^{3/2}\sqrt{pq}
       \le\rho+2\sqrt{s_u}.
\]
This follows from $A\le\bar a$, $B\le\bar b$, $m^2=(1-\theta)pq$;
it is used only in the charge $3|\beta|J/2$. Since $\gamma\ge3.16\rho-3-d$, its negative
part is at most $(d-0.16)_+$. Using the upper density only for that negative
part gives $C\rho+\gamma\varphi/2\ge Q$. Absolute value is separately
convex in each coordinate of $ck+\psi u$, proving the formula for $B$.
For fixed $r$ and $|H|$, the resulting upper bound is convex in $c$.
Take those endpoints \emph{before} completing the square. At either
endpoint $c>0$, it is convex in $r$, since $\mathcal G$ and the maximum of the two outer quadratics are
convex, every other term of $A$ is affine, and $Q$ is affine in $r$. This proves the first line of~\eqref{ms:fourier:rescore:entry-four-corners};
$fh-qh^2\le f^2/(4q)$ proves the second. Bilinearity of $Q$ proves
its positivity throughout the rectangle from its corners.

\end{proof}

\paragraph{Signed phase curvature in the resource schedule}
\label{ms:fourier:rescore:re-signed-phase-repair}
The auxiliary Fourier orientation does not change the signed third moment
in the actual transfer identity. The paired phase estimate must therefore
retain its sign. On a fixed resource cell put
\[
 L=\frac{997}{2500}\left(1-\frac{5r}{2N}\right),\qquad
 g_0(\rho)=3.16\rho-3-d,\qquad U=0.4,
\]
where $r,d$ are the preceding raw and defect upper bounds, and the
bounds must satisfy $0\le L\le U$. The actual density lies in $[L,U]$
and the actual signed moment satisfies $\gamma\ge g_0(\rho)$. Consequently
\begin{equation}\label{ms:fourier:rescore:re-sign-aware-price}
 C\rho+\tfrac12\gamma\phi
 \ge q_{\rm s}(C,\rho):=
 C\rho+\tfrac12\min\{L g_0(\rho),U g_0(\rho)\}.
\end{equation}
This follows by first multiplying the lower bound on $\gamma$ by the
nonnegative actual density, and then minimizing over its two endpoints.
The simpler expression $C\rho+Lg_0(\rho)/2$ is licensed if $g_0\ge0$,
or if nonnegativity of the actual $\gamma$ has been separately proved.
It is not licensed merely by reflecting the auxiliary Fourier integral.

The corrected price preserves the four-corner reduction. For fixed $C$,
$q_{\rm s}$ is concave in $\rho$, being affine plus the minimum of two
affine functions. For fixed $\rho$ it is affine in $C$. In the phase
envelope $A+f|H|-q_{\rm s}H^2$, the remaining terms are separately
convex in $C,\rho$, by the moment and coefficient bounds already proved.
Thus the entire envelope is separately convex, and two successive
one-dimensional endpoint bounds place its maximum at the four corners.
The same argument applies after adding the affine conserved-resource
curvature $e_\rho$ and subtracting a barrier price $K\rho$.
If the four corrected divisors are positive, concavity makes them
positive throughout the rectangle, and square completion gives the
valid raw and resource bounds
\[
 \max_{C,\rho\text{ corners}}\left(A+\frac{f^2}{4q_{\rm s}}\right),
 \qquad
 \max_{C,\rho\text{ corners}}
       \left(A+\frac{f^2}{4(q_{\rm s}+e_\rho)}-K\rho\right).
\]
Any threshold-root export likewise uses the corrected corner curvature.

When the actual signed lower bound \eqref{ms:fourier:rescore:re-gamma} is nonnegative,
the simpler nonnegative-skewness price is licensed. Otherwise retain
\eqref{ms:fourier:rescore:re-sign-aware-price}. Auxiliary reflection alone never supplies
this license. Every initial or left finite pass must retain the sign-aware price
unless its certificate supplies a separate license for a simpler
specialization.

There is also a convenient cell-constant correction. Since $\rho\ge1$,
$(-g_0(\rho))_+\le(d-0.16)_+$, so the weaker price
\[
 C\rho+\tfrac12L g_0(\rho)
       -\tfrac12(U-L)(d-0.16)_+
\]
is valid. This weaker cell-constant form and the exact sign-aware form are both
permitted. Raw remainders and barrier-subtracted resource budgets remain
distinct. A threshold-root export uses the same corrected curvature as
the coefficient envelope from which it is derived.

\paragraph{The unsigned first-stage cover is the same enclosure}
\label{ms:fourier:rescore:re-unsigned}
The first finite barrier uses a cover of $\rho$ alone before introducing
the signed cover. Here are its exact coefficient choices, so that the
first barrier uses the same analytic enclosure. Fix a closed
cell $\rho\in[l,u]\subset[1,1.85]$, actual radius $1\le R\le6$, and
write $k_j(R)$ for the radial ratios of Lemma~\ref{ms:fourier:rescore:re-pricing}. Set
\begin{gather*}
 M=\min\{11.1,Ru,1+k_4(R)(u-1)\},\quad
 g=\min\{1.85,u,\sqrt{M-1}\},\\
 E_5=\min\{RM,1+k_5(R)(u-1)\},\qquad
 b=\sqrt{[u\sqrt{u^2+8}-u^2-2]/2}.
\end{gather*}
The last expression solves $\rho\ge(1+\beta^2)/\sqrt{1-\beta^2}$,
proved by $\rho m\ge a_+^2+a_-^2$ and $m^2\le a_+a_-$.
If a previous pass licenses $D\le d<2$, the generic-event lemma also
permits the ceilings $h_*u+Rd/2$ and $h_*^2u+R^2d/2$ for $M,E_5$;
all dependent quantities are recomputed after taking these minima.
Use the triangle low-frequency triple and the minimum of complete permitted high-frequency
values, with the same moment cap. Put
\[
 K_F=\tfrac{16}{25}(L_0+K_H+L_2/\sqrt N),\qquad s_F=\tfrac{16}{25}L_H.
\]
The radial identity bounds $|K_1|$ by
\[
 k_{\rm rad}=3ub+
     \frac{R(R+2-3l/2)}{R+1/2}(u-1),
\]
provided the displayed radial coefficient is at least $31/40$.
Indeed its positive maximum is at $R$ and its negative part is at most
$31/40$ on $[0,R]$ for $1\le\rho\le1.85$; convexity of
$x(x+2-3\rho/2)/(x+1/2)$ proves this directly.
Two additional valid bounds are
\[
 k_{\rm old}=6u-3u^2/2+3b,\qquad
 k_R=Ru/2+\max_{t=l,u}|R/2-3t/2|g+3b.
\]
For $k_R$, bound the positive and negative fourth moments separately by
$R(\rho\pm\gamma)/2$. Before defect feedback use
$k=\min(k_{\rm rad},k_{\rm old},k_R)$.
After feedback one may replace $k_{\rm old}$ by $\max(k_+,k_-)$, where
\begin{gather*}
 g_l=3.16l-3-d,\quad g_u=3.16u-3-d,\\
 k_+=M+3b-\tfrac32\min_{r=l,u;,v=g_l,g_u}rv,\\
 k_-=3u+3b+\max_{r=l,u;,v=g_l,g_u}(3r/2-3)v.
\end{gather*}
These substitutions bound the lower affine graph for the actual $\gamma$,
not an independently chosen skewness.

Put $\Delta=2u$ initially and $\Delta=(1-3.16)l+3+d$ after feedback.
The constant coefficient of the raw identity is bounded by
\begin{align*}
 A_* =\min\bigl\{&\tfrac32b(2u+1)+\tfrac32\Gamma(R,l)\Delta
       +\tfrac38u^2(u+1)+\tfrac8{27}u^3-l/2,\\
      &\tfrac32b(2u+1)+2uM-l/2\bigr\}.
\end{align*}
This is the coarse gap bound above together with the two bounds for
$-K_q$. The density coefficient obeys
\[
 v_H=\min\{0.4j+0.4\max(2,M-3)/6,
                     0.4\max(1/2,2M/3-1/2)\},
\]
where $j=3(u+1)$ initially; the generic-gap choice is
$j=\alpha(u+1)+(3-\alpha)d/2$, with its checked $\alpha<3$.
For a score endpoint enclosure $C\in[c_-,c_+]$, the resulting raw
bound is
\begin{equation}\label{ms:fourier:rescore:re-unsignedbound}
 K_F+c_+A_*+\frac{(c_+k+v_H+s_F)^2}{4q_*},
 \qquad q_*>0,\quad k\ge0,\quad A_*>0,
\end{equation}
where initially $q_*=c_-l-g/5$. After feedback use the larger of that
quantity and
\[
 c_-l+\varphi_-(3.16l-3-d)/2
                  -(2/5-\varphi_-)(d-0.16)_+/2.
\]
The score endpoints are the lower score floor and the upper allowed
score cap; if the latter depends on $\rho$, use its values at $u,l$
for $c_-,c_+$, respectively. Convexity in the actual score precedes the
completion of the square, and these outward endpoint intervals cover it.
For the phase resource add \eqref{ms:fourier:rescore:re-extra} before maximizing. A barrier
$e\ge K/n$ subtracts $K l$ from the nonnegative resource. The first raw
target must pass the license of Lemma~\ref{ms:fourier:rescore:re-weight} first.

\paragraph{Optional signed improvement inside an unsigned cell.}
If $g_0=\max(-g,\sigma_l l-3-d)>g$, exclude the cell. When $g_0\ge0$,
the sign-moment constraint forces
$(1-\beta)^3\le(1+\beta)(u-g_0)^2$.
This inequality is monotone on $(-1,1)$, so exact bisection yields a
lower $b_0\le\beta$; if it fails even at $b$, the cell is empty.
For $x\in[g_0,g]$, set
\begin{align*}
 f_-(x)&=(l+x)^2/[2(1+b)]-R(u-x)/2-3b-3ux/2,\\
 f_+(x)&=R(u+x)/2-(l-x)_+^2/[2(1-b_0)]-3b_0-3lx/2.
\end{align*}
The signwise fourth-moment Cauchy bounds prove
$\min f_-\le K_1\le\max f_+$. The first minimum is at its clipped
vertex $3u(1+b)/2-l-R(1+b)/2$. The second maximum is among its endpoints,
the clipped kink $l$, and the clipped vertex
$l+(R/2-3l/2)(1-b_0)$. Intersect with $[-k,k]$ and replace the linear
term in \eqref{ms:fourier:rescore:re-unsignedbound} by the smaller full rectangle bound
\[
 \max_{c=c_-,c_+;,k=k_l,k_u;,\psi=\varphi_-,0.4;,
             v=-1/2,\,2M/3-1/2}|ck+\psi v|+s_F.
\]
All sixteen corners are retained. If $b_0>0$, the opposite-sign event
also has probability at least $(1+b_0^2)^2/(2u^2)$, from
$m\ge(1+\beta^2)/(2\rho)$ and $m^2\le pq$. It may improve the complete
high-frequency value. These are intersections with necessary conditions; the cover
still examines every $\rho$-cell.

\paragraph{Moving from a small defect to probability coordinates.}
For $0<d\le0.7$, start with
\[
 A_l=\frac{1.47\sqrt{2-d}-\sqrt{3.7d}}{3.52},\qquad
 B_l=\frac{2.04\sqrt{2-d}-\sqrt{3.7d}}{3.52}.
\]
They follow from $A+B\ge\sqrt{2-d}$ and
$|\sqrt{\sigma+1}A-\sqrt{\sigma-1}B|\le\sqrt{3.7d}$;
their positivity is checked, rather than assumed at every $d$.
The pointwise inequality
$x(x+b)\ge(29/5)b(x-b)$ follows because the difference is
$(x-12b/5)^2+b^2/25>0$. For $x>b>0$, applying it to the two residual
third moments, gives
$D\ge(29/5)\min((\sigma+1)A,(\sigma-1)B)\theta$.
Start $\Theta=1$ and perform the declared finite number of updates:
\begin{align*}
 \Theta&\leftarrow\min\{\Theta,d/[(29/5)
                 \min((104/25)A_l,(54/25)B_l)]\},\\
 \epsilon&=\sqrt{d/[(616/25)(1-\Theta)^2]},\qquad
 [l,u]=[(4-\sigma_u)/2-\epsilon,(4-\sigma_l)/2+\epsilon],\\
 A_l&\leftarrow\max\{A_l,\sqrt{(1-\Theta)l/(1-l)}
                           -d/[2(1-u)(179/25)A_l^2]\},\\
 B_l&\leftarrow\max\{B_l,\sqrt{(1-\Theta)(1-u)/u}
                           -d/[2l(129/25)B_l^2]\}.
\end{align*}
The last two bounds use
$D\ge2(\sigma+2)B^2M_++2(\sigma+4)A^2M_-$ and
$\Delta_A=M_-/q$, $\Delta_B=M_+/p$.
Every division, interval endpoint and variance domain is gated.
Subdivide this entire probability interval, then use the common
cell update and coefficient atlas. Thus the unsigned first barrier,
the signed initialization and the final probability refinements are
three data choices for the same finite-enclosure theorem.


\begin{lemma}[Asymmetric moment and probability restrictions]\label{ms:fourier:rescore:u-asymmetric}
For a sign subprobability measure supported on $[0,R]$, denote its
moments of orders $0,1,2,3,4,5$ by $P,m,a,b,c,e$.  Then
\begin{align}\label{ms:fourier:rescore:u-weighted-cauchy}
 (Rm-a)^2&\le(RP-m)(Ra-b),\nonumber\\
 (Ra-b)^2&\le(Rm-a)(Rb-c),\\
 (Rb-c)^2&\le(Ra-b)(Rc-e).\nonumber
\end{align}
Consequently, whenever the respective denominators are positive,
\begin{equation}\label{ms:fourier:rescore:u-asymmetric-moments}
 P\ge\frac mR+\frac{(Rm-a)^2}{R(Ra-b)},\quad
 c\le Rb-\frac{(Ra-b)^2}{Rm-a},\quad
 e\le Rc-\frac{(Rb-c)^2}{Ra-b}.
\end{equation}
These apply separately with the actual radii $R_-$ and $R_+$.
Furthermore, on a signed cell with $m\ge m_-$ and
$\beta\in[\beta_-,\beta_+]$,
\begin{equation}\label{ms:fourier:rescore:u-p-cauchy}
 \frac{2m_-^2}{1+\beta_+}\le p\le
 1-\frac{2m_-^2}{1-\beta_-}.
\end{equation}
\end{lemma}
\begin{proof}
Apply Cauchy--Schwarz to the pairs $(1,x)$, $(x^{1/2},x^{3/2})$
and $(x,x^2)$ under the nonnegative measure $(R-x)\mu(dx)$;
this gives the three
moment log-convexity inequalities in \eqref{ms:fourier:rescore:u-weighted-cauchy}.
Rearrangement gives \eqref{ms:fourier:rescore:u-asymmetric-moments}.
Zero denominators are handled in \eqref{ms:fourier:rescore:u-weighted-cauchy} before any
division: a positive numerator is incompatible, while omitting a
zero-over-zero nonnegative correction is a valid weaker bound.
The two sign masses sum to at most one, and to one after zero is
excluded from support.  Their necessary probability inequalities are
therefore valid in both cases.  Finally
$m^2\le p\E X_+^2=p(1+\beta)/2$ and
$m^2\le(1-p)(1-\beta)/2$ prove \eqref{ms:fourier:rescore:u-p-cauchy}.
\end{proof}

All occurrences of the signed third moment remain signed.  In particular,
with a prior raw bound $r$ and defect bound $d_D$, set
$L_\phi=0.3988(1-5r/(2N))$ and $g_0(\rho)=3.16\rho-3-d_D$.
The safe lower phase price is
\begin{equation}\label{ms:fourier:rescore:u-signed-price}
 C\rho+\tfrac12\min\{L_\phi g_0(\rho),\ 0.4g_0(\rho)\}.
\end{equation}
This is precisely \eqref{ms:fourier:rescore:re-sign-aware-price}:
a negative lower bound for $\gamma$ must be multiplied by the upper
density.  Reflection used inside an absolute Fourier bound does not
license changing the sign of the actual contact law.

\begin{lemma}[Positive endpoint floors from the defect]\label{ms:fourier:rescore:u-product}
If $D\le d_D<\sqrt{10}-1$, the central gap $(-A,B)$ is genuine and
\[
 A\ge\frac{\sqrt{10}-1-d_D}{6(1-p)},\qquad
 B\ge\frac{\sqrt{10}-1-d_D}{6p}.
\]
On a closed panel $p\in[l,u]\subset(0,1)$, replace the numerators
by $\sigma_- -1-d_D$ with $\sigma_-^2<10$, and the denominators
by $6(1-l)$ and $6u$, respectively, whenever the numerator is positive.
\end{lemma}
\begin{proof}
The gap identity and its signwise Tonelli estimate in
\eqref{ms:defect-decomposition} give
\[
 J\le m(A+B)^2+T_++T_-=\rho+2mAB,
 \qquad D\ge(\sqrt{10}-1)\rho-6mAB.
\]
If zero were in the support, the same proof gives
$D\ge(\sqrt{10}-1)\rho$, a contradiction.  Thus the displayed gap
exists.  Since $\rho\ge1$, $m^2\le p(1-p)$,
$A\le m/(1-p)$ and $B\le m/p$, we have
$mAB\le(1-p)A$ and $mAB\le pB$.
Substitution proves both floors.  Their panel versions weaken positive
denominators in the safe direction and apply on both probability halves.
They can be intersected with all prior endpoint floors.
\end{proof}

\paragraph{Three consequences of a signed defect bound}\label{ms:fourier:rescore:u-geometry}

Once the gap is acquired, use the actual conditional representation
\begin{equation}\label{ms:fourier:rescore:u-coordinates}
 X=\ell(B_0+U-p),\quad\Pp(B_0=1)=p,\quad
 \E(U\mid B_0)=0,\quad
 \ell^2=\frac{1-\theta}{p(1-p)}.
\end{equation}
Write $q=1-p$, $s_j=\E(U^2\mid B_0=j)$,
$\lambda=n(qs_0+ps_1)$,
$\theta_-=q\ell^2s_0$, $\theta_+=p\ell^2s_1$, and
$\theta=\theta_-+\theta_+$.  The actual endpoints satisfy
\[
 U_0\in[-A_0,a_0],\qquad U_1\in[-a_1,A_1],\qquad
 d=a_0+a_1,\quad A=\ell(p-a_0),\quad B=\ell(q-a_1).
\]
Thus $a_0<p$ and $a_1<q$; neither inward endpoint is discarded.
Here $D$ is the moment defect and $d$ is the inward displacement.

Let $A_*,B_*>0$ be acquired lower gap bounds, $S=A_*+B_*$,
$k_+=3(\sqrt{10}+1)B_*$ and
$k_-=3(\sqrt{10}+3)A_*$.
The separate-charge identity \eqref{ms:simultaneous-gap-formula} gives
\begin{equation}\label{ms:fourier:rescore:u-one-resource}
 D\ge b+6(1-\theta)S d+k_+\theta_++k_-\theta_-.
\end{equation}
For example, a panelwise lower binary charge is
\[
 b=\frac{(6+2\sqrt{10})(1-\Theta)^{3/2}}{\sqrt{v_{\max}}}
       \operatorname{dist}\left(\frac{4-\sqrt{10}}2,[l,u]\right)^2,
\]
rounded downward.  To check the change of coordinates in
\eqref{ms:fourier:rescore:u-one-resource}, the two conditional-mean gaps in
\eqref{ms:simultaneous-gap-formula} sum to $\ell d$ and
$m\ell=1-\theta$; its additional product term is nonnegative.

\begin{lemma}[Affine, moment, and variance consequences]\label{ms:fourier:rescore:u-three-uses}
Suppose $p\in[l,u]\subset(0,1)$, $\theta\le\Theta<1$ and
$D\le D_*$.  Put
\[
 v_{\min}=\min\{l(1-l),u(1-u)\},\qquad
 v_{\max}=\max_{x\in[l,u]}x(1-x),\qquad T=D_*-b.
\]
The maximum includes $1/4$ if the panel crosses $1/2$.
Then the following are consequences for the same law.

\textup{(i)} With $k_*\le\min(k_+,k_-)$ positive,
\begin{equation}\label{ms:fourier:rescore:u-affine-interface}
 D\ge b+g d+h\lambda,\qquad
 g=6(1-\Theta)S,\qquad h=\frac{k_*(1-\Theta)}{n v_{\max}}.
\end{equation}
If $T<0$ the panel is empty; if $T=0$ then $d=\lambda=0$.
For $T>0$ one has the simultaneous inequality
\begin{equation}\label{ms:fourier:rescore:u-joint}
 \frac d{d_0}+\frac\lambda\Lambda\le1,
 \qquad d_0=T/g,\quad\Lambda=T/h.
\end{equation}
Reducing the three lower coefficients in \eqref{ms:fourier:rescore:u-affine-interface},
or enlarging either intercept in \eqref{ms:fourier:rescore:u-joint}, is valid.
The affine coefficients cannot be inferred from intercept bounds alone.

\textup{(ii)} If $A_0,A_1>0$, the stronger signed prices are
\begin{equation}\label{ms:fourier:rescore:u-moment-price}
 \left(k_++\frac{6Sq}{A_1}\right)\theta_+
 +\left(k_-+\frac{6Sp}{A_0}\right)\theta_-\le T.
\end{equation}
For $j=3,4,5$, let $b_j(p)=(pq^j+qp^j)/(pq)^{j/2}\le B_j$ on
the panel, and choose upper Taylor coefficients
\[
 L_+=L_j(R_+,\sqrt{(1-l)/l}),\quad
 L_-=L_j(R_-,\sqrt{u/(1-u)}),
\]
where
\[
 L_3(R,x)=R+2x,\quad L_4(R,x)=R^2+2Rx+3x^2,
 \quad L_5(R,x)=R^3+2R^2x+3Rx^2+4x^3.
\]
Set $a_j=[1-(1-\Theta)^{j/2}]/\Theta$ if $\Theta>0$, and
$a_j=j/2$ otherwise.  If $\widehat k_\pm>0$ are lower bounds
for the two prices in \eqref{ms:fourier:rescore:u-moment-price}, then
\begin{equation}\label{ms:fourier:rescore:u-moment-bound}
 \E|X|^j\le B_j+T\max\left\{0,
 \frac{L_+-a_jB_j}{\widehat k_+},
 \frac{L_--a_jB_j}{\widehat k_-}\right\}.
\end{equation}

\textup{(iii)} Actual signed support bounds $[-R_-,R_+]$ give
\begin{equation}\label{ms:fourier:rescore:u-amplitudes}
 \ell\ge\ell_0:=\sqrt{(1-\Theta)/v_{\max}},\quad
 A_0\le R_-/\ell_0-l,\quad A_1\le R_+/\ell_0-(1-u).
\end{equation}
If both outward caps are at most $A^*>0$ and \eqref{ms:fourier:rescore:u-joint} has
positive intercepts, then
\begin{equation}\label{ms:fourier:rescore:u-harmonic}
 \lambda\le L_*:=\left\{\frac1\Lambda+
          \frac1{nA^*\max(u,1-l)d_0}\right\}^{-1},\qquad
 \theta\le\frac{L_*}{n v_{\min}+L_*}.
\end{equation}
In addition,
\begin{equation}\label{ms:fourier:rescore:u-mean}
 \sqrt{v_{\min}(1-\Theta)}\le m\le\sqrt{v_{\max}}.
\end{equation}
\end{lemma}
\begin{proof}
For (i), use $\theta=\lambda(1-\theta)/(npq)$ in
\eqref{ms:fourier:rescore:u-one-resource} to obtain the displayed affine
coefficients. Apply Theorem~\ref{ms:joint-resource} with its index floor
chosen as the actual $n$, and with $d_D=D_*$, $v_*=v_{\max}$,
$S_*=S$ and $K=k_*$. The input $nD\le nD_*$ follows from $D\le D_*$;
this proves the stated joint inequality and both zero cases.
For (ii), take expectations of
$(U_0+A_0)(a_0-U_0)\ge0$ and
$(U_1+a_1)(A_1-U_1)\ge0$.  Conditional centering gives
\[
 s_0\le A_0a_0,\quad s_1\le A_1a_1,\quad
 (1-\theta)s_0=p\theta_-,\quad(1-\theta)s_1=q\theta_+.
\]
Substituting $d\ge s_0/A_0+s_1/A_1$ in
\eqref{ms:fourier:rescore:u-one-resource} proves \eqref{ms:fourier:rescore:u-moment-price}.
This spends its displacement term for the moment estimate; that term
must not be added again to \eqref{ms:fourier:rescore:u-moment-price}.
The original inequality \eqref{ms:fourier:rescore:u-one-resource} remains separately true.

The positive-measure Taylor identity
\eqref{ms:fourier:rescore:re-conditional} gives
\[
 \E|X|^j\le(1-\theta)^{j/2}b_j(p)
                          +L_+\theta_++L_-\theta_-.
\]
First replace $b_j(p)$ by $B_j$ in its nonnegative coefficient.
Then convexity gives $(1-\theta)^{j/2}\le1-a_j\theta$ on
$[0,\Theta]$.  Maximize the remaining linear functional over the
nonnegative signed variances subject to \eqref{ms:fourier:rescore:u-moment-price},
using \eqref{ms:fourier:rescore:re-simplex}.  This proves
\eqref{ms:fourier:rescore:u-moment-bound}, even if one Taylor cost after subtraction is
negative.  A zero outward amplitude makes that centered conditional
class deterministic and is treated directly before division.

For (iii), subtract the conditional centers $-p\ell,q\ell$ from
actual support to obtain \eqref{ms:fourier:rescore:u-amplitudes}.  The variance bounds
just proved imply
\[
 \lambda\le nA^*(qa_0+pa_1)
          \le nA^*\max(u,1-l)d.
\]
Insert this lower bound for $d$ into \eqref{ms:fourier:rescore:u-joint}; then use
$\theta=\lambda/(npq+\lambda)$, increasing in $\lambda$ and
decreasing in $pq$, to get \eqref{ms:fourier:rescore:u-harmonic}.
Finally conditional variance gives $m^2=pq(1-\theta)$ and hence
\eqref{ms:fourier:rescore:u-mean}.  These proofs also cover $p\ge1/2$ and panels crossing
$1/2$; no replacement of $v_{\min}$ by $l(1-l)$ is made there.
\end{proof}

For the upper Bernoulli moments in (ii), writing $v=pq$ gives
$B_3(v)=(1-2v)/\sqrt v$, $B_4(v)=1/v-3$, and
$B_5(v)=(1-4v+2v^2)/v^{3/2}$.  Each decreases on $(0,1/4]$, so its
value at $v_{\min}$ is a valid bound on either probability half.
Lower prices in \eqref{ms:fourier:rescore:u-moment-price} use the lower numerators
$q\ge1-u$, $p\ge l$ and the upper amplitudes in
\eqref{ms:fourier:rescore:u-amplitudes}.  When a full two-sided noise amplitude $A^*$ is
exported, the inward bound $d_0\le A^*$ is also checked.
Previously acquired fourth and fifth moment bounds may be retained by
intersection, including $\E|X|^5\le
\max(R_-,R_+)\E X^4$.  The updated variance cap in
\eqref{ms:fourier:rescore:u-harmonic} can be rounded upward and intersected with its
previous value; it does not invalidate any earlier inequality.

\paragraph{Keeping the score and phase coherent}\label{ms:fourier:rescore:u-phase}

The raw envelope in \eqref{ms:fourier:rescore:re-quadratics} and its phase
improvement \eqref{ms:fourier:rescore:re-extra} give
\begin{equation}\label{ms:fourier:rescore:u-conserved}
 \kappa nD\le A(C,\rho)+F_{\operatorname{sgn}H}(C,\rho)|H|
       -Q(C,\rho)H^2-n(C-C_{\mathrm E})\rho,
 \qquad\kappa=997/15000.
\end{equation}
The curvature $Q$ already includes the extra phase price
\[
 \frac12\left(1-\frac{5r}{2N}\right)
                   \left(0.4095\rho-\frac{d_D}{15}\right).
\]
It is not added a second time.  The signed price
\eqref{ms:fourier:rescore:u-signed-price} is retained before any square completion.

A positive barrier uses the charge $K\rho$.  The selected
zero-barrier probability refinements also retain the actual score and
the fixed rational upper bound $C_{\mathrm E}<0.40974$:
\begin{equation}\label{ms:fourier:rescore:u-actual-charge}
 n(C-C_{\mathrm E})\rho\ge\chi(C)\rho,
 \qquad\chi(C)=n(C-0.40974)_+.
\end{equation}
Split the score interval at $0.40974$ when that point is inside.
The rectangle proof above is separately convex in the
actual $C$ and $\rho$.  On either closed score piece,
$\chi(C)\rho$ is affine in either variable when the other is fixed;
its subtraction preserves that convexity.  Hence both moment endpoints
at each score endpoint give a valid coherent corner envelope.
The selected score endpoints $0.4095,0.40974,c_1$ give six
corners.  Keeping only the two extreme score endpoints would miss the
kink and is not this argument. Raw remainders and the first unweighted passes continue to use their
uncharged envelope.  Actual-score charging here is a refinement of the
probability kernel and its exported phase quadratics.

\begin{lemma}[The signed phase interface]\label{ms:fourier:rescore:u-phase-interface}
Suppose the acquired charged corner bounds are
$q_i^{\pm}(h)=A_i+F_i^{\pm}h-Q_i h^2$, and $0\le h\le H_\pm$.
Then
\begin{equation}\label{ms:fourier:rescore:u-phase-export}
 D\le\frac{15000}{997n}\max_i q_i^{\operatorname{sgn}H}(|H|).
\end{equation}
A fixed target envelope is valid if, for each source corner and each
sign, one target corner dominates it on the entire consumed interval
$[0,H_\pm]$.  For a difference $a+bh+ch^2$, it is enough to check
its values at the two endpoints and also $a-b^2/(4c)$ when $c>0$
and $0<-b/(2c)<H_\pm$.
\end{lemma}
\begin{proof}
Equation~\eqref{ms:fourier:rescore:u-phase-export} is \eqref{ms:fourier:rescore:u-conserved} after
corner domination and division by the positive $\kappa n$.
The minimum of a quadratic on a compact interval occurs at an endpoint
or an interior stationary point that is a minimum.  The stated tests
therefore dominate every source corner pointwise; taking maxima
preserves the inequality.  Negative linear coefficients keep their
signs, and the normalization $15000/(997n)$ is unchanged.
\end{proof}

In particular, if $D\le D_*$ and \eqref{ms:fourier:rescore:u-affine-interface} holds,
define $D_H$ as the smaller of $D_*$ and the right side of
\eqref{ms:fourier:rescore:u-phase-export}.  The actual law satisfies
\begin{equation}\label{ms:fourier:rescore:u-phase-joint}
 g d+h\lambda\le D_H-b.
\end{equation}
A negative right side excludes that phase.  If $T=D_*-b>0$,
\eqref{ms:fourier:rescore:u-phase-joint} is equivalently
$d/d_0+\lambda/\Lambda\le(D_H-b)/T$.
The entire joint residual changes with phase, not just one intercept.
At a count threshold $k+u$, the same actual phase is
$H=(k+u-np)/\sqrt{pq+\lambda/n}$.  This is the interface to the
count proof; it prevents independent maximization of count phase and
moment budget.

\paragraph{Radius feedback and the final local interface}

In the actual coordinates \eqref{ms:fourier:rescore:u-coordinates}, suppose a
separate confinement argument has already proved $|U|\le a$ on a
probability panel $p\in[l,u]\subset(0,1)$. Since
$\ell^2=(1-\theta)/(pq)\le1/(pq)$, it follows that
\begin{equation}\label{ms:fourier:rescore:re-radiusfeedback}
 |X|\le
 \frac{\max\{u+a,1-l+a\}}{\sqrt{\min\{l(1-l),u(1-u)\}}}.
\end{equation}
Indeed $|B_0+U-p|\le\max\{p+a,1-p+a\}$, and the concavity of
$p(1-p)$ places its minimum on the closed panel at an endpoint.
The implication includes zero noise. The new radius is used only
after confinement has established its hypothesis, not in the argument
that establishes confinement.

For completeness, the signed lower bound \eqref{ms:fourier:rescore:re-gamma}
requires only $p\in[l,u]\subset(0,1/2)$, $\theta\le\Theta<1$, and
the already established radius $R$. To see the sign directly, write
$X=V+W$, $\E[W\mid B_0]=0$. The bound $X\ge-R$ gives
$W^3\ge(-R-V)W^2$, so
\[
 \gamma\ge\E V^3+2\E(VW^2)-R\theta.
\]
The smaller conditional value of $V$ is
$-\sqrt{(1-\theta)p/q}\ge-\sqrt{u/(1-u)}$. The binary cubic
$(1-\theta)^{3/2}(1-2p)/\sqrt{p(1-p)}$ decreases in both $p<1/2$
and $\theta$. These observations prove precisely
\eqref{ms:fourier:rescore:re-gamma}; its nonnegativity, whenever needed, is a
separate scalar inequality checked on the actual panel.

Finally suppose a phase cell has a coherent upper envelope
$\mathcal B(H)\le T_H$ in \eqref{ms:fourier:rescore:re-phaseprice}, and a proved
moment inequality for the same law is
\[
 D\ge\alpha d_Y+b\theta+P,\qquad
 \alpha>0,\quad b\ge0,\quad P\ge0.
\]
If $n\ge N$, $\ell\ge\ell_-\ge0$ and $\lambda\ge L_0\ge0$, then
\begin{equation}\label{ms:fourier:rescore:re-localG}
 nd_Y\le G:=
 \frac{T_H/\kappa-NP-b\ell_-^2L_0}{\alpha}.
\end{equation}
This follows from $n\theta=\ell^2\lambda$ and $nP\ge NP$.
Thus $G<0$ excludes the entire cell. If the upper envelope is strict,
the exported inequality is strict and $G\le0$ excludes the cell.
For a concave quadratic $A+Fh-Qh^2$, $Q>0$, its maximum on a closed
phase interval is its value at the clipped vertex $F/(2Q)$. Taking
the maximum over coherent corners, followed by the minimum with an
independent scalar bound, produces $T_H$. This operation reuses the
single resource \eqref{ms:fourier:rescore:re-phaseprice}; it creates no additional
allowance for the subsequent count argument.

\paragraph{Finite sequences of enclosure refinements}\label{ms:fourier:rescore:u-propagation}

\begin{theorem}[Preservation of simultaneous bounds]\label{ms:fourier:rescore:u-contract}
Start with a complete closed probability cover and the actual support
and scalar hypotheses \eqref{ms:fourier:rescore:u-profile}.  Each surviving panel carries
simultaneous moment, phase, support and resource inequalities for one
hypothetical maximizing law in that panel.  Perform any declared finite
sequence of the following operations:
\begin{enumerate}
\item a complete signed initial or probability resource cover, using
      the enclosure theorem and all its domain gates;
\item a signed support-ray exclusion using only previously valid
      same-law moment bounds, followed by moment estimates on that
      acquired support;
\item the endpoint, moment, variance and phase consequences proved
      above, or intersection with a previously acquired bound;
\item outward rounding; or pooling all active panels by taking maxima
      of upper bounds and minima of lower bounds, then using their
      closed interval hull for a new complete initial cover.
\end{enumerate}
Keep an external set of already excluded panels.  The result remains a
complete cover of every possible maximizing law, with valid simultaneous
bounds on every surviving panel.  A final scalar, affine and quadratic
interface to a fixed count contract is valid if all the inequalities
above are checked in the required directions on the same row.
\end{theorem}
\begin{proof}
Theorem~\ref{ms:fourier:rescore:re-enclosure} proves the
within-cover induction, including complete incidences, empty cells,
positive denominators and square completion.  The support-ray implications and Lemma~\ref{ms:fourier:rescore:u-phase-interface}
prove the additional implication steps.
No operation changes the marginal law or uses a proposed radius before
that radius has been proved.  Intersections and outward rounding preserve
true inequalities for that law.

If no active panel remains, the band is excluded and no minimum or
maximum of an empty family is taken.  Otherwise the actual surviving
law belongs to at least one active panel.  Every pooled lower bound is no larger than that panel's lower
bound, and every pooled upper bound is no smaller.  This applies
separately to the two support radii and all signed moment endpoints.
The interval hull enlarges the surviving union, so it is a safe domain
for the next implication with these pooled hypotheses.  It does not
reinstate previously excluded laws: their exclusions remain external.
Shared boundary points belong to their adjacent closed panels.
Induction over the finite schedule now proves the assertion.  A bound
on a maximum of coherent phase corners is transferred only by
Lemma~\ref{ms:fourier:rescore:u-phase-interface}, and a direct affine resource is
transferred by decreasing its three actual lower coefficients.
\end{proof}

For a band $N\le n\le L$, a nonnegative conserved budget $T$ gives
$D\le T/(\kappa N)$; if $T<0$, \eqref{ms:fourier:rescore:u-conserved} already
contradicts $D\ge0$ before any replacement of $n$ by $N$.
Weighted resources $nD\le N d_D$ are propagated in their weighted
form using Theorem~\ref{ms:fourier:rescore:joint-resource}.
An unweighted inequality $D\le d_D$ is never multiplied by $N$ and
reported as a bound for $nD$.
Index monotonicity is used only where the retained formula has its
explicit gate; support is checked separately at every listed integer.
The low-frequency transfer needs $n>2$ and the high-frequency comparator
$n\ge5$, both satisfied here. There is no residual $n\ge36$
assumption in the earlier-index moment and Fourier formulas.

\subsection{Confinement and conditional-count witnesses}
\label{ms:fourier:local}

\begin{remark}[Preserved index scope]\label{ms:fourier:scope47}
The supplied confinement calculation below has floor $N=47$. Its rational
normalizers, the $n-2$ concentration bound and all numerical applications
retain that scope. The subsequent local-comparison identities are stated
with explicit parameters, but this does not establish their numerical
hypotheses at floor $42$. Appendix~\ref{ms:uniform-bridge} gives the
conditional parameter changes for that floor. The complete selected
resource, confinement, normalization and local-cover predicates and their
composition for every $n\ge42$ remain obligations. The former index-$46$ barrier assembly
and the historical tail-completion assertion are not used here.
\end{remark}

\subsubsection{One confinement lemma with two concentration choices}\label{ms:fourier:cut49:conf:re:confinement}
A proved amplitude is used as a support update, never as a premise of
its own proof. The following derivation applies to every acquired row,
with the same formulas for the central and left probability ranges.
\paragraph{Statement and hypotheses}
Let $n\ge N=47$, $Y=B+U$, $\Pp(B=1)=p<1/2$ and $\E(U\mid B)=0$.
Put $q=1-p$, $s=\E U^2$, $\lambda=ns$ and $X=\ell(Y-p)$, where
$\E X^2=1$. Thus $\ell^2=(1-\theta)/(pq)$ and $\theta=\ell^2s$.
The conditional label-zero support is at most $a\ge0$ and the label-one
support is at least $b\le1$, where $a,b$ are support points and
$b-a=1-d$. In particular $0\le a\le d$, $1-d\le b\le1$.
Assume the initial strip $\operatorname{supp}Y\subset(-0.55,1.55)$.
Write $\rho=\E|X|^3$, $\beta=\E[X|X|]$ and $\epsilon=\E|U|$.

Select an acquired closed probability panel $[l,h]$.
Its fields give upper bounds $\Theta,d_0,\rho_0,\beta_+,T,\Lambda$,
a lower bound $\beta_-$, and separate bounds $H_+,H_-$ for the two signs of
$H$. Assume
\[
 \theta\le\Theta,\quad d\le d_0,\quad \rho\le\rho_0,
 \quad \beta_-\le\beta\le\beta_+,\quad 0\le\lambda\le\Lambda,
 \quad 0.4<C\le C_0=0.4253.
\]
For the positive sign $0\le H\le H_+$; for the negative sign
$-H_-\le H\le0$. Let $x=np+H/\ell$ and
$\nu=\mathcal L(x-\sum_{i<n}Y_i)$. Assume a contact majorant $Q$ for
$g(y)=\nu([y,\infty))$, with equality at each support point of $Y$,
and derivative
\begin{equation}\label{ms:fourier:cut49:conf:derivative}
 \sqrt n(-Q'(y))=\ell\phi(H/\sqrt n)
 +\frac\ell n\{3C\beta+(H\phi(H/\sqrt n)+3C\rho)t-3Ct|t|\},
 \quad t=\ell(y-p).
\end{equation}
Both actual deleted convolutions satisfy, separately,
\begin{equation}\label{ms:fourier:cut49:conf:concentration}
 Q_j(w)\le\frac{\ell\phi(0)w+2C_0\rho}{\sqrt j},
 \qquad j=n-1,n-2,
\end{equation}
for every interval of width $w\ge0$.

Here is the precise common resource assumption. Put
$\kappa=997/15000<\phi(0)/6$, $J=\int_{[a,b)}(y-a)\nu(dy)$,
and note that the central gap contributes $m\ell J=(1-\theta)J$ to $L$.
Thus $L\ge(1-\theta)J$; retain $e\ge0$ and
\begin{equation}\label{ms:fourier:cut49:conf:budget}
 n e\rho+n\phi(0)D/6+n^{3/2}L<\mathcal B(H),\qquad
 \mathcal B(H)=\min\{T,\max_j(A_j+F_{j,\pm}|H|-Q_jH^2)\}.
\end{equation}
The four quadratics are the table's quadratic rows, using
$Q_{\rm budget}$, not $Q_{\rm raw}$. The Gaussian phase contribution
has already been included in $Q_{\rm budget}$.

For each panel let $E,F_\ell,m_0$ be respectively strict lower, upper,
and lower rational bounds for
\[
 \sqrt{(1-\Theta)/(h(1-h))},\qquad
 1/\sqrt{l(1-l)},\qquad \sqrt{(1-\Theta)l(1-l)}.
\]
The bounds use the grid $10^{-6}$, with an extra unit outward. Define
\begin{align*}
 c_+&=2(5.16)E^2(1-h-d_0)^2,&
 c_-&=2(7.16)E^2(l-d_0)^2,\\
 b_*&=\min\{3(4.16)E(1-h-d_0),3(6.16)E(l-d_0)\},\\
 P&=m_0E^2\pos{3.5\operatorname{dist}([l,h],[0.41886,0.41887])-2.041d_0}^2.
\end{align*}
The simultaneous moment inequality has the following usable form:
\begin{equation}\label{ms:fourier:cut49:conf:moment}
 D\ge c_+M_++c_-M_-+b_*\theta+P,\quad
 d=(M_+/p+M_-/q)/\ell,\quad
 \epsilon\le2(M_++M_-)/\ell,\quad M_\pm\ge0.
\end{equation}
To derive it, the endpoint expansion of \eqref{ms:defect-decomposition} gives
$D\ge m(\sqrt{\sigma+1}A-\sqrt{\sigma-1}B)^2
+2(\sigma+2)B^2M_++2(\sigma+4)A^2M_-
+3(\sigma+1)B\theta_++3(\sigma+3)A\theta_-$.
Insert $A\ge E(l-d_0)$, $B\ge E(1-h-d_0)$ and the rational lower
coefficients displayed above. For the square, write
$A=\ell(p-a)$, $B=\ell(1-p-(1-b))$, with
$a+(1-b)=d\le d_0$. Its unshifted coefficient is
$(\sqrt{\sigma+1}+\sqrt{\sigma-1})(p-p_*)$;
$\sqrt{\sigma+1}+\sqrt{\sigma-1}>3.5$ and
$\sqrt{\sigma+1}<2.041$ prove the stated lower square $P$.
For completeness, on the positive class put $Z=X-B\ge0$.
Then $T_+=2B^2M_++3B\E[Z^2;X>0]+\E[Z^3;X>0]$ and
$R_+=BM_++\E[Z^2;X>0]$. Consequently
$(\sigma-1)T_++6BR_+\ge2(\sigma+2)B^2M_+
+3(\sigma+1)B\theta_+$, since $\E[Z^2;X>0]\ge\theta_+$.
The reflected expansion gives the negative-class bound.
The centered absolute deviation on each class is at most twice its
mean distance from the inner endpoint, proving the bound on $\epsilon$.
All quantities refer to the same marginal.

\paragraph{Concentration from an antichain}
\begin{lemma}[Strict-gap concentration]
Suppose $Y=U_0+B(1+U_1-U_0)$, where $B$ is Bernoulli$(p)$ independent
of the pair $(U_0,U_1)$, and $1+U_1-U_0\ge c>0$ almost surely.
For independent copies $Y_1,\ldots,Y_j$ and any interval $I$ of diameter
strictly less than $c$,
\[
 \Pp\left(\sum_{i=1}^jY_i\in I\right)
 \le \max_{0\le k\le j}\binom jk p^k(1-p)^{j-k}.
\]
The same bound holds after adding any independent random variable.
\end{lemma}
\begin{proof}
Apply the conditional random-chain argument in
Appendix~\ref{ms:uniform-bridge}. That argument uses only independent
Bernoulli labels and positive increments at least $c$, so it applies to
the present coupled pairs at every $j$. Its interval diameter is strictly
less than $c$, including when either endpoint is closed. Conditioning on
any added independent variable gives the last assertion.
\end{proof}
For the actual marginal we may couple the two conditional noises
independently and then sample an independent label. Since $U_0\le a$ and
$U_1\ge b-1$, the step is at least $b-a=1-d$.
Strict diameter is necessary: a Bernoulli law on $\{0,c\}$ assigns mass
one to $[0,c]$, exceeding its largest atom when $0<p<1$.

\paragraph{A rational cap for every remaining index.}
Put $J_0=N-2=45$ and define
\[
 r_L=l-(1-l)/J_0,\qquad r_R=h+h/J_0,\qquad
 v_* =\min\{r_L(1-r_L),r_R(1-r_R)\}.
\]
Require $0<r_L\le r_R<1$. Adjacent binomial-mass ratios
show that any mode $k$ at index $j\ge J_0$ satisfies
$(j+1)p-1\le k\le(j+1)p$, hence $r=k/j\in[r_L,r_R]$.
In particular $1\le k\le j-1$. Finite Stirling bounds give
\[
 \sqrt j\binom jk p^k(1-p)^{j-k}
 \le\frac{\exp(1/(12j)-jD(r\Vert p))}{\sqrt{2\pi r(1-r)}}
 <\frac{0.4}{\sqrt{v_*}(1-1/(12J_0))}.
\]
Here $D(r\Vert p)\ge0$ follows from Jensen's inequality for the logarithm;
$\exp x\le(1-x)^{-1}$ for $0\le x<1$ follows term by term from the
power series. The finite Stirling inequalities used are
$\sqrt{2\pi m}(m/e)^m<m!<\sqrt{2\pi m}(m/e)^m\exp(1/(12m))$.
Choose the strict $10^{-6}$ grid upper square root $M$ of
$0.16/[v_*(1-1/(12J_0))^2]$. Thus, whenever $w<1-d$,
\begin{equation}\label{ms:fourier:cut49:conf:antichainconcentration}
 Q_j(w)<M/\sqrt j,\qquad j\ge n-2\ge J_0.
\end{equation}
This covers the $n-1$ full summands of $\nu$ and the $n-2$ full summands
remaining after one mark or deletion. An independent remaining conditional
noise preserves the cap. We do not replace the latter index by $n-1$.

Either the score-derived cap or the strict-gap cap may be selected on
an interval where its hypotheses hold. The first confinement uses the
score cap; later central and left schedules use the strict-gap cap.
For the left panels replace the conditional cubic coefficient 24 below
by 28; in general any $c_\rho$ satisfying
$c_\rho^2[l(1-l)]^3>[3(1-l)+0.55]^2$ is valid.
\paragraph{Signed deletion, including atoms}
Let $m=n-1$, $Z=x-\sum_{i=1}^mY_i$, $\mu=x-mp$ and let $K=\sum B_i$.
For a fixed removed label put
$T_b(I)=\Pp(x-\sum_{i<m}Y_i-b-U_b\in I)$.
Exchangeability and $mp-K=Z-\mu+\sum U_i$ give the exact identities
\[
 \nu(I)-T_1(I)=\frac{\E[(Z-\mu+\sum U_i)1_I(Z)]}{mp},\qquad
 T_0(I)-\nu(I)=\frac{\E[(Z-\mu+\sum U_i)1_I(Z)]}{mq}.
\]
There is no absolute value around the sum of noises.
For a full removed pair $(Y_i,U_i)$, condition on the pair. The remaining
sum gives a function $f(Y_i)\in[0,Q_{m-1}(w)]$. Since $\E U_i=0$,
\[
 \E[U_i f(Y_i)]\le Q_{m-1}(w)\E(U_i)_+
 =\tfrac12\epsilon Q_{m-1}(w).
\]
This also works when one conditional law is deterministic.

Let independent conditional noises satisfy $\Delta=U_1-U_0\ge-d$.
For $\tau>0$ put $I^+=I+1+[-d,\tau]$. On $\Delta\le\tau$, the
point for $T_1$ in $I$ maps to the point for $T_0$ in $I^+$. On the
complement, conditioning on both noises gives error at most
$Q_{m-1}(w)\Pp(\Delta>\tau)$. Moreover
\begin{equation}\label{ms:fourier:cut49:conf:bad}
 \Pp(\Delta>\tau)\le
 \min\left\{\frac{s}{s+p\tau^2},\frac d{d+\tau}\right\}.
\end{equation}
Indeed $\E\Delta=0$, $\operatorname{Var}\Delta\le s/p$, and
$\Delta+d\ge0$ has mean $d$. Cantelli and Markov prove the two bounds.
The exceptional event is strict, so $\Delta=\tau$ belongs to the good event.

If $I,I^+$ are disjoint, combine the two signed noise terms before
bounding them. For the same remaining-sum law their conditional
probabilities $f_I,f_J$ satisfy
\[
 0\le f_I\le Q_I,\quad0\le f_J\le Q_J,\quad f_I+f_J\le Q_{\rm hull}.
\]
Since $p<q$,
\[
 \frac{f_I}{p}+\frac{f_J}{q}
 \le\min\left\{\frac{Q_I}p+\frac{Q_J}q,
 \frac{Q_{\rm hull}}q+\left(\frac1p-\frac1q\right)Q_I\right\}.
\]
Multiplication by $\E U_+=\epsilon/2$ proves the combined bound.
Closed, open and half-open intervals are all covered by containment in
the corresponding closed interval. No density or atom separation is used.

\paragraph{Slope bounds and the two exterior intervals}
Split $|H|$ into 16 closed equal cells $[h_l,h_h]$ and $\lambda$ into
24 closed equal cells $[s_l,s_h]$. On each box put
\[
\begin{aligned}
 \rho_* &=\min\{\rho_0,\rho_B(l)+24s_h/N\},\\
 b_- &=\max\{\beta_-,1-2h-2(1-h)\Theta\},\qquad
 b_+ =\min\{\beta_+,1-2l+\Theta\}.
\end{aligned}
\]
Here $\rho_B(p)=(1-2p(1-p))/\sqrt{p(1-p)}$. Write $s_i=\E U_i^2$. Conditional centering and the two sign supports give
\[
 \rho=\ell^3\{pq(1-2pq)+3pq(s_0+s_1)
                  +p\E U_1^3-q\E U_0^3\}.
\]
The initial strip and $d_0<0.55$ imply $|U_i|<0.55$.
Thus $\rho\le\rho_B+\ell^3(3q+0.55)s\le\rho_B+24s$; the last
coefficient inequality is checked as
$24^2[l(1-l)]^3>[3(1-l)+0.55]^2$.
Similarly $\beta=(1-\theta)(1-2p)+\ell^2(ps_1-qs_0)$, and
$|\ell^2(ps_1-qs_0)|\le\theta$. These identities prove the stated
outward bounds on $\rho$ and $\beta$.
Set $f=0.4F_\ell$, $\varphi_-=0.39886(1-h_h^2/(2N))>0$.
The exact upper bound for $H\phi(H/\sqrt n)$ used below is
\[
 v_+=\begin{cases}0.4h_h&H\ge0,\\-\varphi_-h_l&H\le0.\end{cases}
\]
For all $t\le0$, completing the square in \eqref{ms:fourier:cut49:conf:derivative} gives
\[
 \mathcal M_-=F_\ell\pos{(v_++3C_0\rho_*)^2/(12(0.4))-3(0.4)b_-},
 \qquad L_-=E\varphi_- -\mathcal M_-/N.
\]
A positive bound $B_*>\sqrt{(1-l)/l}$ gives valid central-gap lower
and positive-half-line upper slope bounds
\begin{align*}
 \mathcal M_+&=F_\ell\pos{0.4h_hB_*+3C_0B_*(B_*-1)-1.2b_-},\\
 L_g&=E\varphi_--\max(\mathcal M_-,\mathcal M_+)/N,\\
 U_+&=f+\frac{F_\ell}{N}
 \left\{3C_0b_++\frac{(0.4h_h+3C_0\rho_*)^2}{12(0.4)}\right\}.
\end{align*}
They follow by completing a square on the positive half-line and using
$t(t-1)\le B_*(B_*-1)$ for $1\le t\le B_*$; on $0\le t\le1$ that
term has the favorable sign. Negative values are covered by $\mathcal M_-$.
All divisions and positivity requirements are explicit gates.

\paragraph{The sharper lower-ray slope.}
Take $\tau_-=5r_0/12$, $\eta=10^{-5}$ and
$I=[a-r,a-d-\tau_--\eta]$, where $r\ge r_0$.
Since $a\le d$, the entire interval has
$t\le t_*=-E(l+\tau_-+\eta)<0$.
Write $B(t)=b_-+\rho_*t+t^2$ and
$B_*^{\min}=B(\min(t_*,-\rho_*/2))$. For the actual $C$ and $t\le t_*$,
\[
 3C(\beta+\rho t+t^2)+H\phi t
 \ge1.2B(t)+v_+t+3(C_0-0.4)\min(0,B_*^{\min}).
\]
The inequality uses $t<0$ and separates the possible negative part of
$B(t)$ before varying $C$. Minimize the right side at the clipped vertex
$\min(t_*,-(v_++1.2\rho_*)/2.4)$ and call its value $q_*$. Then
\[
 L=E\varphi_--F_\ell\max(0,-q_*)/N
\]
is a lower slope on the whole interval, for every $n\ge N$.
Positive corrections are discarded. Therefore contact gives
\[
 \sqrt n\nu(I)\ge L(r-d-\tau_--\eta).
\]
This one slope replaces the previous separate bounds on the long and
short parts of the interval.

\paragraph{Lower interval geometry.}
Its translate is $I^+=[a-r+1-d,b-\eta]$, of width $r-\eta$.
The initial strip implies $I^+\subset[a+g,b)$,
$g=1-0.55-2d_0>0$, hence $\nu(I^+)\le J/g$.
The intervals are disjoint throughout $r\le0.55+d_0$, since
$\inf I^+-\sup I=1-r+\tau_-+\eta>0$.
At $r=r_0$, concentration and directed distance bounds are
\begin{align*}
 Q_I&=f(r_0-\tau_--\eta)+2C_0\rho_*,&
 Q_J&=f(r_0-\eta)+2C_0\rho_*,\\
 Q_{\rm hull}&=f(1+r_0-\eta)+2C_0\rho_*,\\
 D_I&=\pos{v_H-l-\tau_--\eta},&
 D_J&=\pos{v_H+1-l-\eta},
\end{align*}
where $v_H=-h_l/F_\ell$ for $H\ge0$ and $v_H=h_h/E$ for $H\le0$.
These signs follow from $\mu=p+H/\ell$.

\paragraph{Upper interval geometry.}
For a support point $y\ge1+r_0$, take $\tau_+=7r_0/15$,
$I=[a,a+\tau_+]$ and $I^+=[b,a+1+2\tau_+]$.
The strict gates $d_0+2\tau_+<r_0$ and $1-d_0-\tau_+>0$ give
$I^+\subset[b,y)$ and disjointness. Contact, including the atom at $b$,
and the leakage definition give
\[
 \sqrt n\nu(I)\ge L_g(1-d)-\sqrt nJ/\tau_+,\qquad
 \sqrt n\nu(I^+)\le U_+(0.55+d).
\]
Use
\begin{align*}
 Q_I&=f\tau_++2C_0\rho_*,&Q_J&=f(2\tau_++d_0)+2C_0\rho_*,\\
 Q_{\rm hull}&=f(1+2\tau_+)+2C_0\rho_*,\\
 D_I&=\pos{d_0+\tau_++h_h/E-l},&
 D_J&=\pos{1+d_0+2\tau_++h_h/E-l}.
\end{align*}
The upper side deliberately keeps the absolute phase bound.

\paragraph{The strict-gap alternative.}
If $1-0.55-2d_0>0$ and the displayed lower intervals have width below
$1-d$ throughout $r\in[r_0,0.55+d_0]$, take $Q_I=Q_J=M$.
For the upper intervals the corresponding gate is
$1-2d_0-2\tau_+>0$; then again $Q_I=Q_J=M$.
Use the same directed distances. Their concentration caps are now
independent of $r$, and use
$a_e=F_dM\{1/l+1/(1-h)\}/2$ in the next formulas.
The lower-ray gate reduces to $L>0$: its contact side increases with
$r$, while the entire error is constant on the acquired finite strip.
For the score-derived cap retain the explicit affine slope gate below.
No assertion outside that strip is needed.

\paragraph{One budget, two error bounds}
Let $F_d$ be the strict grid upper bound for
$\sqrt{\max\{(47/46)^3,47/45\}}$. It simultaneously covers the two
concentration factors in \eqref{ms:fourier:cut49:conf:concentration} after multiplication by
$\sqrt n$, uniformly for $n\ge47$. For either pair of intervals set
\begin{align*}
 a_e&=\frac{F_d}{2}\min\left\{\frac{Q_I}l+\frac{Q_J}{1-h},
 \frac{Q_{\rm hull}}{1-h}+\left(\frac1l-\frac1{1-l}\right)Q_I\right\},\\
 A&=F_d\left\{Q_ID_I/l+Q_JD_J/(1-h)\right\},\qquad B=F_dQ_I.
\end{align*}
The probability denominators are outward because $1/p-1/(1-p)$ is
decreasing on $(0,1/2)$.
Use $a_d=L$, $a_J=1/g$ on the lower side, and
$a_d=L_g+U_+$, $a_J=1/\tau_+$ on the upper side.
For the phase cell, maximize each concave budget quadratic at its
vertex clipped to $[h_l,h_h]$, take their maximum, then the minimum
with $T$; call this upper bound $T_*$. Put
\[
 R=T_*/\kappa-b_*E^2s_l,
 \quad\alpha=\min(c_+lE,c_-(1-h)E),\quad G=(R-NP)/\alpha.
\]
If $G\le0$, the strict common budget excludes the entire closed box.
Otherwise $nd<G$ and
\[
 B_\tau=\min\{s_h/(s_h+Nl\tau^2),G/(N\tau+G)\}
\]
bounds \eqref{ms:fourier:cut49:conf:bad}. With
\[
 z=\max\left\{\frac{a_d/l+2a_e}{Ec_+},
 \frac{a_d/(1-h)+2a_e}{Ec_-},\frac{\kappa a_J}{1-\Theta}\right\},
\]
one valid total error bound is
\begin{equation}\label{ms:fourier:cut49:conf:oldbound}
 \mathcal E_0=(zR+A)/N-zP+B B_\tau.
\end{equation}
This follows by multiplying the single budget by $z/\kappa$.
The inverse-index coefficient $zR+A$ is checked positive.

A stronger bound spends the same resource on the pair event. Let
$R_0=R-NP>0$ and define the three pairs
\[
 (v_+,w_+)=\left(\frac{a_d/l+2a_e}{Ec_+}\frac{R_0}{N},
                    \frac{R_0}{Ec_+l}\right),
\]
\[
 (v_-,w_-)=\left(\frac{a_d/(1-h)+2a_e}{Ec_-}\frac{R_0}{N},
                    \frac{R_0}{Ec_-(1-h)}\right),\qquad
 (v_J,w_J)=\left(\frac{\kappa a_J}{1-\Theta}\frac{R_0}{N},0\right).
\]
They are upper cost/displacement coordinates at the vertices of the
simplex with resources $nc_+M_+$, $nc_-M_-$ and $n^{3/2}L/\kappa$.
For any $t\ge0$,
\[
 \frac{D}{N\tau+D}\le\frac{N\tau}{(N\tau+t)^2}D+
 \frac{t^2}{(N\tau+t)^2}\qquad(D\ge0).
\]
The difference is $N\tau(D-t)^2/[(N\tau+D)(N\tau+t)^2]$.
Consequently a further valid error bound is
\[
 \mathcal E_t=\frac AN+
 \max_{i\in\{+,-,J\}}\left(v_i+B\frac{N\tau}{(N\tau+t)^2}w_i\right)
 +B\frac{t^2}{(N\tau+t)^2}.
\]
Take the minimum of \eqref{ms:fourier:cut49:conf:oldbound} and these eleven
bounds with $t=jG/10$, $j=0,\ldots,10$. Each is individually valid;
no numerical minimization or derivative of their minimum is assumed.

The lower contradiction is the strict gate
$Lr_0-L(\tau_-+\eta)>\mathcal E$.
To extend it to every $r\ge r_0$, bound increments directly in the original
affine-in-width signed deletion bound for the same law. The budget gives
$n\epsilon<E_*=2T/(\kappa E\min(c_+,c_-))$. The sufficient slope gate is
\[
 L>\frac{F_df}{N}\left\{D_I/l+D_J/(1-h)
 +\frac{E_*}{2}\left(1/l+1/(1-h)\right)\right\}+F_dfB_\tau.
\]
Using the uncoupled increment is safe even when the hull bound was used
at $r_0$. The upper contradiction is $L_g-0.55U_+>\mathcal E$.
Thus no support point lies below $a-r_0$ or above $1+r_0$.
Since $a\ge0$, $b\ge1-d_0$ and $d_0<r_0$, the label-zero and label-one
supports both satisfy $|U|<r_0$.

\paragraph{The selected schedules.}
The ordinary, central-antichain, and left-antichain certificates substitute
their recorded $r_0$ into these same gates. Noise-dependent amplitudes
are the same theorem with a positive lower noise endpoint in
$R=T_*/\kappa-b_*E^2s_l$. In the central rows the additional amplitudes
are $0.35,0.30,0.275$ on $[0.05,0.08],[0.08,0.10],[0.10,\Lambda]$;
on the left they are $0.375,0.325,0.30$. Intersect each interval with its
own row's $\Lambda$ and each amplitude with that row's global amplitude.
The required acceptance checks both signs of $H$, all closed phase/noise
cells, clipped endpoints, and strict budget exclusions. No proof of
whole-count, nearest-count, or binomial score comparison is used by this confinement lemma.

\subsubsection{The conditional-count comparison}\label{ms:fourier:cut49:shortloc:section}

A law concentrated near two points has a natural integer count and a small
centered error. We compare it with its own Bernoulli count. The same exact
identity handles every noise regime. There are two useful estimates of its
right side: discard all but the nearest outgoing count, or retain all counts
through a random cutoff. The estimates share their normalization, moments,
and tail bounds; only the final scalar witness changes.

\paragraph{The identity and sufficient inequalities}

Let $0<p\le1/2$, $q=1-p$, $v=pq$, and let
\[
 Y=B+U,\qquad \Pp(B=1)=p,\qquad \E(U\mid B)=0.
\]
The two actual conditional laws $U_0,U_1$ obey
\begin{equation}\label{ms:fourier:cut49:shortloc:support}
 |U_j|\le A,\quad U_0\le a_0,\quad U_1\ge-a_1,\quad
 a_0,a_1\ge0,\quad d=a_0+a_1\le d_0<\min(p,q),\quad d_0\le A.
\end{equation}
Write $s_j=\E U_j^2$, $s=qs_0+ps_1$, $\lambda=ns$,
$\ell=(v+s)^{-1/2}$, $\theta=s/(v+s)$ and
$\rho=\E|\ell(Y-p)|^3$. At the actual threshold choose
\[
 x=np+H/\ell=k+u,\qquad k\in\{1,\ldots,n-1\},\quad |u|\le\tfrac12,
 \qquad \delta=k-np=H/\ell-u.
\]
Either nearest integer is allowed at a tie. Put
\[
 \pi_j=\binom njp^jq^{n-j},\quad F_B(j)=\Pp(\operatorname{Bin}(n,p)\le j),
 \qquad \rho_B=(1-2v)/\sqrt v,
\]
with $\pi_j=0$ outside $[0,n]$. Let $W_j$ be an independent sum of $j$
copies of $U_1$ and $n-j$ copies of $U_0$. The signed scores are
\[
 C=\frac{\sqrt n}{\rho}\{\Pp(\textstyle\sum Y_i\le x)-\Phi(H/\sqrt n)\},
 \qquad C_B=\frac{\sqrt n}{\rho_B}\{F_B(k)-\Phi(\delta/\sqrt{nv})\}.
\]
Throughout this section $C\ge c_*>0$; the selected application uses $c_*=0.4$.
Define
\begin{align}
 f(u)&=\frac{\Phi((\delta+u)/\sqrt{n(v+s)})-
                    \Phi(\delta/\sqrt{n(v+s)})}{\pi_k},\label{ms:fourier:cut49:shortloc:f}\\
 Q&=\Pp(\textstyle\sum Y_i\le k+u)-F_B(k),\nonumber\\
 \mathcal R&=\sqrt n\{\Phi(\delta/\sqrt{nv})-
                 \Phi(\delta/\sqrt{n(v+s)})\}+C(\rho_B-\rho).\nonumber
\end{align}
Equation~\eqref{ms:master-identity}, with $V_j=W_j$ and the same
interior count $k$, gives the exact identity
\begin{equation}\label{ms:fourier:cut49:shortloc:master}
 \rho_B(C-C_B)=\sqrt n\pi_k\{Q/\pi_k-f(u)\}+\mathcal R.
\end{equation}
In particular, the phase $u$, the sign of $H$, and both conditional laws
refer to the same marginal throughout.

\begin{theorem}[Parameterized local comparison]\label{ms:fourier:cut49:shortloc:theorem}
Fix $n\ge N$, $\lambda>0$ and a closed parameter cell on which
\eqref{ms:fourier:cut49:shortloc:support} holds. Suppose $\sqrt n\pi_k\ge m>0$.
Either of the following witnesses proves $\rho_B(C-C_B)<-\varepsilon\lambda$.
\begin{enumerate}
\item \emph{Nearest-count witness.} Valid bounds $O,T,R_0,\kappa$ satisfy
\[
 \Pp(W_k>u)+f(u)\ge O,\quad
 \sum_{i\ge1}\Pp(W_{k+i}\le u-i)\le T,\quad
 \pi_{k+i}/\pi_k\le R_0,\quad \mathcal R\le\kappa\lambda/n.
\]
Only legal counts occur in the sum; adding positive upper bounds for
other counts is allowed. If $D=O-R_0T>0$ and
\begin{equation}\label{ms:fourier:cut49:shortloc:directtest}
             mD-\kappa\lambda/n>\varepsilon\lambda,
\end{equation}
the conclusion follows.
\item \emph{Whole-count witness.} Bounds $B(\lambda)<0$ and $I(\lambda)\ge0$
satisfy
\[
 \rho_B(C-C_B)\le\sqrt n\pi_k\lambda B(\lambda)
                                      +\lambda I(\lambda)/n.
\]
It is enough that
\begin{equation}\label{ms:fourier:cut49:shortloc:whole-test}
                  mB(\lambda)+I(\lambda)/N<-\varepsilon.
\end{equation}
\end{enumerate}
The witnesses may differ from cell to cell. A finite closed cover and a
separate treatment of $\lambda=0$ therefore suffice for a complete local
comparison.
\end{theorem}
\begin{proof}
The signed count identity \eqref{ms:count-difference}, with $V_j=W_j$, is
\begin{equation}\label{ms:fourier:cut49:shortloc:count}
 Q=-\sum_{i=0}^{k}\pi_{k-i}\Pp(W_{k-i}>u+i)
       +\sum_{i=1}^{n-k}\pi_{k+i}\Pp(W_{k+i}\le u-i).
\end{equation}
Apply Proposition~\ref{ms:local-witnesses} with $m_0=m$.
The legal counts, closed incoming events and strict outgoing events are
identical. In the second alternative the present hypothesis $I\ge0$
makes its $\max(0,I)/N$ exactly $I/N$, yielding
\eqref{ms:fourier:cut49:shortloc:whole-test}.
\end{proof}

\paragraph{Binomial normalization and the variance correction}

Here is an explicit finite recipe for the constants required above.
Let $p\in[p_-,p_+]\subset(0,1/2)$, $n\ge N$, $\lambda\le\Lambda$,
$\theta\le\Theta$, and $|\delta|\le\Delta$, where $\Delta\ge1$ also includes
every binomial mode. Choose $0<r_-<p_- -\Delta/N$ and
$p_++\Delta/N<r_+<1/2$. Let $v_*>0$ be a lower bound for $p(1-p)$.
Choose positive $E,L,R$ bounding, on $[r_-,r_+]$,
\[
 |D'''(r\Vert p)|=\frac{1-2r}{r^2(1-r)^2},\qquad
 |(\log(r(1-r)))'|=\frac{1-2r}{r(1-r)},\qquad
 \frac1{12r(1-r)},
\]
respectively, where $D(r\Vert p)=r\log(r/p)+(1-r)\log((1-r)/q)$.
Then, with $Z(t)=Lt/2+R+Et^3/(6N)$,
\begin{equation}\label{ms:fourier:cut49:shortloc:mass}
 \sqrt n\pi_k=\frac{\phi(0)}{\sqrt v}
       e^{-\delta^2/(2nv)+\zeta},\qquad n|\zeta|<Z(|\delta|).
\end{equation}
For completeness, the finite Stirling formula used here is
\[
 \log(j!)=(j+\tfrac12)\log j-j+\tfrac12\log(2\pi)+r_j,
 \qquad 0<r_j<1/(12j).
\]
If $a_j=\log(j!)-(j+1/2)\log j+j$ and $z=(2j+1)^{-1}$, the integrated
geometric series gives
$0<a_j-a_{j+1}=\sum_{i\ge1}z^{2i}/(2i+1)<1/[12j(j+1)]$.
Summing gives the remainder relative to $\lim a_j$. The integration-by-parts
recursion for $\int_0^{\pi/2}\sin^j t\,dt$ and the inequalities between three
successive integrals give Wallis's product; its finite product equals
$e^{4a_j-2a_{2j}}j/[2(2j+1)]$ and tends to $\pi/2$. Thus
$\lim a_j=\tfrac12\log(2\pi)$. Applying this formula to $n!,k!,(n-k)!$
and Taylor-expanding $D$ and the logarithmic prefactor proves
\eqref{ms:fourier:cut49:shortloc:mass}; the factorial error is at most $1/[12nr(1-r)]$.

Consequently a valid mass lower bound is any positive
\begin{equation}\label{ms:fourier:cut49:shortloc:m}
 m<\frac{0.39886}{\sqrt{p_+(1-p_+)}}
          \left[1-\frac{\Delta^2/(2v_*)+Z(\Delta)}N\right].
\end{equation}
The same formula at a mode proves $\sqrt n\max_j\pi_j<1$ whenever
$v_*[1-Z(\Delta)/N]^2>0.4^2$ and $Z(\Delta)<N$.
The rational bounds $0.39886<\phi(0)<0.4$ follow, for example, by the
alternating arctangent series in Machin's identity.

Differentiating \eqref{ms:fourier:cut49:shortloc:f} gives
\begin{equation}\label{ms:fourier:cut49:shortloc:logder}
 \log f'=-\tfrac12\log(1+s/v)+\delta^2/(2nv)
                   -(\delta+u)^2/(2n(v+s))-\zeta.
\end{equation}
Thus $n|\log f'|\le B_0\lambda+P(D,a)$ on $|\delta|\le D\le\Delta$,
$|u|\le a\le1/2$, where
\[
 B_0>\frac1{2v_*}+\frac{\Delta^2}{2Nv_*^2},\qquad
 P(D,a)=\frac{2Da+a^2}{2v_*}+Z(D).
\]
Choose $0<b<1$ above $[B_0\Lambda+P(\Delta,1/2)]/N$ and put
$F=(1-b)^{-1}$. The elementary estimates $e^z\le(1-z)^{-1}$ and
$1-e^{-z}\le z$ imply
\begin{equation}\label{ms:fourier:cut49:shortloc:derivative}
 |f'-1|\le(B_1\lambda+P_1)/n,\qquad B_1=FB_0,\quad P_1=FP(D,a).
\end{equation}
This supplies both a derivative interval $l_f\le f'\le h_f$ and the
finer phase-error bound $|u-f(u)|\le a(B_1\lambda+P_1)/n$.

The exact count variance is
\begin{equation}\label{ms:fourier:cut49:shortloc:variance}
 t_j=\E W_j^2=\lambda+(j-np)(s_1-s_0),\qquad
 |s_1-s_0|\le s/p\le c_v s,\quad c_v\ge1/p_-.
\end{equation}
It gives $\omega_-\lambda\le t_k\le\omega_+\lambda$ with
$\omega_\pm=1\pm c_vD/N$. All uses below require the checked condition
$\omega_->0$. Also
$t_{k+i}\le iV_*\lambda$ for $V_*=1+c_v(\Delta+1)/N$.
The binomial recursion gives, for every legal $i\ge1$,
\begin{equation}\label{ms:fourier:cut49:shortloc:ratio}
 \frac{\pi_{k+i}}{\pi_k}\le R_0,
 \qquad R_0>1+\frac{p_++\Delta-1}
 {Nv_*[1-(\Delta-1)/(Np_-)]},
\end{equation}
provided the denominator is positive and $p_++\Delta<2$.
Indeed the first ratio has this bound and all subsequent rising factors
are at most one. No count outside $[0,n]$ is used in this recursion.

Finally \eqref{ms:fourier:cut49:shortloc:support} fixes the signs inside the absolute cubic.
Expansion and centering give
\[
 \E|Y-p|^3=v(1-2v)+3pq(s_0+s_1)
                   +p\E U_1^3-q\E U_0^3
 \ge v(1-2v)+(3p-d)s.
\]
The inequality $(1+z)^{3/2}\le1+3z/2+3z^2/8$ for $z\ge0$ yields
\begin{equation}\label{ms:fourier:cut49:shortloc:cubic}
 (v+s)^{3/2}(\rho-\rho_B)\ge s
     \left(6p-3p^2-\tfrac32-d-\frac{3(1-2v)s}{8v}\right).
\end{equation}
Let $Q_*>0$ be a lower bound for the bracket and
$m_c\le c_*/[p_+(1-p_+)+\Lambda/N]^{3/2}$.
The normal variance correction is nonpositive for $\delta\le0$;
for $\delta>0$ its derivative in $s$ is at most $0.4\delta/(2v^{3/2})$.
Therefore
\begin{equation}\label{ms:fourier:cut49:shortloc:repair}
 \mathcal R\le(P_*D_+-m_cQ_*)\lambda/n,
 \qquad P_*\ge\frac{0.4}{2v_*^{3/2}},\quad D_+\ge\delta_+.
\end{equation}
One may take $D_+=D$, but retaining the positive part is often decisive.
All bounds concern $v+s$, which can exceed $1/4$.

\paragraph{Two elementary tools for probabilities}

\begin{lemma}[A shifted quartic]\label{ms:fourier:cut49:shortloc:quartic}
Suppose $W$ is a sum of independent centered variables bounded in absolute
value by $A$, and $t=\E W^2$. Let $K_0\ge(9+6\sqrt3)/4$ and $a>0$.
For $r\ge0$ and $\sigma\in\{-1,1\}$,
\begin{equation}\label{ms:fourier:cut49:shortloc:quarticformula}
 \Pp(W+\sigma r>0)\ge
 \frac{(3a^2-A^2)t-3t^2-4Art-6r^2t-r^4
                      +3a^2r^2+2\sigma a^3r}{K_0a^4}.
\end{equation}
\end{lemma}
\begin{proof}
The polynomial $h_a(z)=(-z^4+3a^2z^2+2a^3z)/(K_0a^4)$ is below
$\mathbf1_{\{z>0\}}$: it is nonpositive for $z\le0$, and its maximum for
$z>0$ is at $z/a=(1+\sqrt3)/2$ and is at most one. Independence gives
$|\E W^3|\le At$ and $\E W^4\le A^2t+3t^2$.
Expand $\E h_a(W+\sigma r)$ and apply these two inequalities.
The value at the threshold is zero, so atoms require no limiting argument.
\end{proof}

\begin{lemma}[One common tail bound]\label{ms:fourier:cut49:shortloc:tails}
If independent centered $X_i\le A$ have total variance at most $V$, then,
for $b>0$, their sum has upper tail bounded by the minimum of one and
\begin{equation}\label{ms:fourier:cut49:shortloc:tailmin}
 \begin{gathered}
 \frac{V}{V+b^2},\qquad
 \mathcal B(A,V,b)=e^{b/A}\left(\frac{V}{V+Ab}\right)^{b/A+V/A^2},\qquad
 \frac{\E(P_\mu-j)_+^3}{(b/A+\mu-j)^3},\\
 \mu=V/A^2,\qquad 0\le j<b/A+\mu.
 \end{gathered}
\end{equation}
where $P_\mu$ is Poisson. The same holds for a lower tail when $|X_i|\le A$.
All bounds apply to closed events, hence also to strict events.
\end{lemma}
\begin{proof}
Cantelli's bound follows by applying Markov to $(\sum X_i+V/b)^2$.
Taylor's integral remainder gives
$\E e^{zX_i}\le\exp\{\E X_i^2(e^{zA}-1-zA)/A^2\}$
for $z\ge0$. Multiply and use $z=A^{-1}\log(1+Ab/V)$ to get $\mathcal B$.
The zero-variance case is immediate.

For the third bound, smooth $f(x)=(x-t)_+^3$ by a nonnegative smooth
compactly supported density. Then $f'',f''',f''''\ge0$. For
$\mathcal Af(x)=f(x+A)-f(x)-Af'(x)$, the Taylor quotient
\[
 \frac{f(x+y)-f(x)-yf'(x)}{y^2}
       =\int_0^1(1-r)f''(x+ry)\,dr
\]
is nondecreasing in $y$. A centered $X\le A$ of variance $v$ consequently
satisfies $\E f(x+X)\le f(x)+(v/A^2)\mathcal Af(x)$.
For $\mathcal P_zf(x)=\E f(x+A(P_z-z))$,
$\partial_z\mathcal P_zf=\mathcal P_z\mathcal Af$ and
\[
 \partial_z^2\mathcal P_zf=\mathcal P_z\mathcal A^2f\ge0,\qquad
 \mathcal A^2f(x)=\int_0^A\!\int_0^A(A-r)(A-w)f''''(x+r+w)\,dr\,dw.
\]
Thus $\mathcal P_{v/A^2}f\ge f+(v/A^2)\mathcal Af$. Replace the independent
variables successively; translations and expectations preserve the derivative
signs and Poisson rates add. Increasing the total variance to $V$ is permitted
by Jensen after adding independent centered Poisson noise. Removing the
smoothing is justified by bounded actual upper supports and finite Poisson
third moments. Markov at $t=A(j-\mu)$ proves the displayed bound.
\end{proof}

These bounds have finite rational evaluations. For example,
\begin{align}\label{ms:fourier:cut49:shortloc:poissonfinite}
 \E(P_\mu-j)_+^3={}&\mu^3+3(1-j)\mu^2+(1-3j+3j^2)\mu-j^3\\
 &+e^{-\mu}\sum_{h=0}^{\lfloor j\rfloor}\frac{\mu^h(j-h)^3}{h!}.\nonumber
\end{align}
Replacing $e^{-\mu}$ by $(\sum_{h=0}^{24}\mu^h/h!)^{-1}$ enlarges the
nonnegative correction. Alternatively sum the positive Poisson series
through 32; the remaining tail is at most its first term divided by
$1-\mu(34-j)^3/[34(33-j)^3]$, provided $j<33$ and that ratio is below one.
For Bennett, put $x=Ab/V$ and
$L=[(b/A+\mu)\log(1+x)-b/A]$. The positive series
$\log(1+x)=2\sum_{h\ge0}z^{2h+1}/(2h+1)$, $z=x/(2+x)$, gives lower
bounds for $L$; any positive lower exponential series gives an upper bound
for $e^{-L}$. Rounding a positive lower $L$ down before exponentiation is
also valid. These are choices of rational witnesses, not optimization claims.

\paragraph{The nearest-count witness, including vanishing noise}

For $u=-r\le0$, add $-h_fr$ to \eqref{ms:fourier:cut49:shortloc:quarticformula} with
$\sigma=1$; for $u=r\ge0$, add $l_fr$ with $\sigma=-1$.
Both bound $\Pp(W_k>u)+f(u)$ from below. For a cell
$t\in[t_-,t_+]$, $r\in[r_-,r_+]$, this expression is concave in $t$
(second derivative $-6/(K_0a^4)$), so only $t_-,t_+$ are needed.
At each endpoint, bound a monomial with positive coefficient at $r_-$
and one with negative coefficient at $r_+$. Any rational $a>0$ is a valid
witness. Other valid choices are
\begin{equation}\label{ms:fourier:cut49:shortloc:ownother}
 \frac{r_-^2}{t_++r_-^2}-h_fr_+\quad(u\le0),\qquad l_fr_-\quad(u\ge0).
\end{equation}
They follow from Cantelli and nonnegativity, respectively. Their maximum
or one selected witness is valid without an optimality proof.

For incoming counts use $t_{k+i}\le iV_*\lambda$. On $u\in[-r_+,-r_-]$
the first distance is $1+r_-$ and every later distance is at least $i$.
On $u\in[r_-,r_+]$ they are $1-r_+$ and $i(1-r_+/2)$ for $i\ge2$.
Let $b,b_o$ be these first and later distances. On $\lambda\in[L,U]$,
let $T_1$ be any bound from Lemma~\ref{ms:fourier:cut49:shortloc:tails} at $(A,V_*U,b)$.
Since $\mathcal B(A,iV,ib)=\mathcal B(A,V,b)^i$, all later counts total at most
$B^2/(1-B)$ for $B=\mathcal B(A,V_*U,b_o)<1$. Thus
\begin{equation}\label{ms:fourier:cut49:shortloc:midtest}
 D=O-R_0\{T_1+B^2/(1-B)\}>0,\qquad
 mD-\kappa U/N>\varepsilon U
\end{equation}
proves the entire cell when $\kappa\ge0$. This sign condition is required on each cell. The argument applies to
both central and left probability panels.

A cell touching $\lambda=0$ requires estimates proportional to $\lambda$;
a positive lower endpoint may not be substituted. The following rescaling
provides them. Put $r=zt$, $t_*=\omega_+U$ and $d_a=K_0a^4$.
The shifted quartic plus $br$, divided by $t$, becomes
\begin{equation}\label{ms:fourier:cut49:shortloc:scaledown}
 F_{\sigma,b}(z,t)=\frac{3a^2-A^2-3t}{d_a}
 +\left(\frac{2\sigma}{K_0a}+b-\frac{4At}{d_a}\right)z
 +\frac{3a^2t-6t^2}{d_a}z^2-\frac{t^3z^4}{d_a}.
\end{equation}
It extends to $t=0$ and is concave in $t$. On $0\le z\le z_0$, replacing
$z^4$ by $z_0^2z^2$ gives a lower quadratic at $t=t_*$; at $t=0$ it is
linear. Endpoints and a clipped quadratic vertex therefore give an exact
lower bound $J(\sigma,b,z_0)$ for the whole rectangle.

For negative phase choose $z_-,c_->0$ and verify
\begin{gather}\label{ms:fourier:cut49:shortloc:tinynegative}
 J(1,-h_f,z_-)>c_-,\quad
 \frac{z_-^2}{1+z_-^2t_*}-h_fz_->c_-,\quad
 \frac{2z_-}{(1+z_-^2t_*)^2}>h_f,\\
 \frac1{1+4t_*}-\frac{h_f}{2}>c_-t_*,\qquad z_-t_*<1/2.\nonumber
\end{gather}
The quartic handles $r\le z_-t$. Beyond that point use
$r^2/(t+r^2)-h_fr$. Its derivative at $r=z_-t$ is positive by the third
gate; its second derivative has just one sign change, so any later critical
point is a maximum. Its minimum is at $z_-t$ or $1/2$, covered by the other
gates. Hence the own count exceeds $c_-t$.
For positive phase choose $z_+,c_+>0$ with
\begin{equation}\label{ms:fourier:cut49:shortloc:tinypositive}
 J(-1,l_f,z_+)>c_+,\qquad l_fz_+>c_+,\qquad z_+t_*<1/2.
\end{equation}
The quartic handles $r\le z_+t$ and the normal term handles the rest.
The own count exceeds $c_+t$ and also exceeds $l_fr$.

For incoming tails proportional to $\lambda$, fix a Poisson anchor $j$.
A positive series term of index $h$ divided by
$\lambda(b/A+\mu-j)^3$ has logarithmic derivative, multiplied by $\lambda$,
\[
 h-1-\mu-\frac{3\mu}{b/A+\mu-j},\qquad \mu=V_*\lambda/A^2.
\]
If $b/A-j>0$ and this is nonnegative at $\lambda=U$ for the first positive
index, it is nonnegative for every term and all $0<\lambda\le U$.
Therefore the Poisson bound divided by $\lambda$ is largest at $U$.
Likewise $B^2/[\lambda(1-B)]$ is increasing whenever
$\lambda B'/B\ge1/2$. Differentiating the Bennett formula gives
\[
 \lambda B'/B=b/A-\mu\log(1+b/(A\mu));
\]
The function $g(\mu)=b/A-\mu\log(1+b/(A\mu))$ decreases because
$\log(1+x)>x/(1+x)$ for $x>0$. Moreover the logarithmic derivative of
$B^2/[\lambda(1-B)]$, multiplied by $\lambda$, is
$2g-1+gB/(1-B)\ge0$ when $g\ge1/2$. Thus checking the gate at $U$
proves the whole interval. A rational upper bound for the logarithm
verifies that gate. On cells using the fixed positive bound $l_fr_-$ instead,
unscaled endpoint tails suffice: after a positive bracket is checked, divide
it by $U$ to obtain a lower bound per $\lambda$. These arguments, followed by
\eqref{ms:fourier:cut49:shortloc:directtest}, prove the nearest-count comparison uniformly down to arbitrary positive
noise, including a deterministic conditional component.

\paragraph{The whole-count witness}

Assume $c=1-d_0>1/2$. Couple all $2n$ conditional variables independently and
set $W_j=\sum_{i\le j}U_{1,i}+\sum_{i>j}U_{0,i}$, extending it constantly
outside $[0,n]$. Every increment of $j+W_j$ is at least $c$. Define
\[
 K=\max\{j\in\mathbb Z:j+W_j\le k+u\},\qquad R=K-k.
\]
The coupling of Section~\ref{ms:counts}, with $V_j=W_j$, gives
\begin{equation}\label{ms:fourier:cut49:shortloc:cutoff}
                  \Pp(\textstyle\sum Y_i\le k+u)=\E F_B(k+R).
\end{equation}
Let $\chi=(n+1)p-k<2$ and $m_B=\max_j\pi_j$. For every integer $r$,
\begin{equation}\label{ms:fourier:cut49:shortloc:curvature}
 F_B(k+r)-F_B(k)-r\pi_k\le\frac{m_B}{(n+1)v}
 \left[(\chi-1)_+r_++\chi_+\binom{r_-}{2}+\binom{r_-}{3}\right].
\end{equation}
This is Lemma~\ref{ms:crossing-curvature}, with the same zero extension,
$\chi<2$ and interior count. Its proof requires only an increment lower
bound $c>1/2$. Here $d\le d_0$ gives increments at least $1-d_0=c$,
and $c\le1$ follows from $d_0\ge0$.

Write $t=\E W_k^2$. Crossing monotonicity and Cauchy--Schwarz give
\begin{equation}\label{ms:fourier:cut49:shortloc:crossingmoments}
 \E R_+\le\frac{t}{(c-u)^2},\quad
 \E\binom{R_-}{2}\le\frac{t}{[c+\min(u,c/3)]^2},\quad
 \E\binom{R_-}{3}\le\frac{t\sqrt{A^2+3t}}{(2c-1/2)^3}.
\end{equation}
These are the crossing-moment bounds of Lemma~\ref{ms:crossing-curvature}
with $A_U=A$, $V_k=W_k$, $|u|\le1/2$ and the increment lower bound
$c=1-d_0$. The lemma includes the strict outgoing crossing and the
fourth-moment estimate from independence and conditional centering.

\paragraph{The marked crossing.}
Let $D_i=U_{1,i}-U_{0,i}$, $\vartheta=\E(D_i)_+=\E(D_i)_-$,
$0<\eta\le1$, and $R_f=\E\{u+\eta-W_k\}$. The two deleted laws are
\[
 A_0\overset{d}=kU_1+(n-k-1)U_0,\qquad
 A_1\overset{d}=(k-1)U_1+(n-k)U_0.
\]
Then
\begin{align}
 \E R&\le u+\eta-R_f+J,\label{ms:fourier:cut49:shortloc:mark}\\
 J&=\frac{\vartheta}{\eta}\sum_{i\ge1}
 \{\Pp(A_0\le u+d_0-ci)+\Pp(A_1>u+\eta+i-A)\}.\nonumber
\end{align}
Here are the pathwise details. Put $z=u-W_k$ and mark
\[
 Z_f=\sum_{i=1}^{n-k}(D_{k+i})_-\mathbf1_{\{W_k\le u-ci\}},\quad
 Z_b=\sum_{i=1}^{k}(D_{k-i+1})_+\mathbf1_{\{W_k>u+\eta+i\}}.
\]
If $R>z+\eta$, the crossing activates the marks through $R$: for $R>0$
use $z\ge cR$, while for $R<0$ use $z<R-\eta$. In either case
$R\le z+Z_f+Z_b$. Hence if $h=R-\lfloor z+\eta\rfloor>0$,
$Z_f+Z_b>h-1+\eta\ge\eta h$. Thus
$R\le\lfloor z+\eta\rfloor+(Z_f+Z_b)/\eta$ on every path.
On an active forward mark $U_0>U_1\ge-d_0$; deleting its pair implies
$A_0\le u+d_0-ci$. On a backward mark delete $U_1\le A$ to obtain the
strict event for $A_1$. The remaining sum is independent of the entire
pair, so the expectations factor. Extending to all $i\ge1$ adds only
nonnegative bounds. Finally $|R|\le1+(|W_k|+|u|)/c$ proves integrability
and permits expectations. No coupling has replaced either deleted marginal.

\paragraph{A finite bound for every mark.}
Fix $\lambda\in[L,U]$, $L>0$, and $t\le\omega_+\lambda$; put
$T_L=\omega_+L$, $T_U=\omega_+U$. On a phase subcell $u\in[v_l,v_h]$,
for $i=1,2$ let $b_1=ci-d_0-v_h>0$ and $b_2=i-A+v_l>0$.
Lemma~\ref{ms:fourier:cut49:shortloc:tails} gives a bound $\omega_+\lambda K_i^{\rm ind}$,
where
\[
 K_i^{\rm ind}=\sum_{j=1}^2\min\left\{\frac1{b_j^2+T_L},
                          \frac{\overline P(A,T_U,b_j)}{T_L}\right\}.
\]
The denominator $T_L$ is a lower bound for the chosen upper envelope
$\omega_+\lambda$, not for any actual variance.
Alternatively couple $A_0=Z+U_1$, $A_1=Z+U_0$ with
$Z\overset d=(k-1)U_1+(n-k-1)U_0$. The two events project into
$Z\le-\alpha$ and $Z>\beta$, where
$\alpha=ci-2d_0-v_h>0$, $\beta=i-A-d_0+v_l>0$.
For any $z\in(-\alpha,\beta)$ their sum is at most
$(\omega_+\lambda+z^2)/\min(z+\alpha,\beta-z)^2$.
Consequently a second kernel is
\[
 K_i^{\rm com}=\frac{1+z^2/T_L}{\min(z+\alpha,\beta-z)^2},\qquad
 z=\max\{-T_L/\beta,\min((\beta-\alpha)/2,T_L/\alpha)\}.
\]
This center is inside the interval; differentiating the two quadratic
ratios on either side of the midpoint proves its optimality, although
any interior rational center suffices. Take the smaller kernel on each
phase subcell and the maximum over their closed cover. Denote it $K_i$.
For $i\ge3$ use fourth moments and the convex midpoint integral for
$(cx-b)^{-4}$. On $u\in[u_l,u_h]$ this gives
\begin{equation}\label{ms:fourier:cut49:shortloc:K}
 K=K_1+K_2+\frac{A^2+3T_U}{3c}
  \left\{\frac1{((5/2)c-A-u_h)^3}
                 +\frac1{((5/2)c-A+u_l)^3}\right\},\qquad
 J\le\frac{\vartheta}{\eta}\omega_+\lambda K.
\end{equation}
All denominators must be positive. Indeed, the tail fourth moment is at
most $A^2t+3t^2$, and integration from $5/2$ to infinity covers every
remaining integer, including the harmless added illegal counts.

If the same-law resource supplies $nd<G$ and
$G^2<c_vNL$, then
\begin{equation}\label{ms:fourier:cut49:shortloc:theta}
 \vartheta\le\frac{c_vG\lambda}{c_vn\lambda+G^2}.
\end{equation}
To prove this, for centered $X\ge-b$, $\E X^2\le V$ and $b^2\le V$,
the quadratic $b(X-V/b)^2/(b+V/b)^2$ majorizes $X_-$ on $[-b,\infty)$.
Its expectation is at most $bV/(V+b^2)$. Apply this to $D_i\ge-d$,
using $\operatorname{Var}(D_i)=s_0+s_1\le c_v\lambda/n$, and monotonicity
in $b\le\sqrt V$. The case $d=0$ is immediate. With $\eta=z\lambda$,
$z>0$, $zU\le1$, \eqref{ms:fourier:cut49:shortloc:K} becomes
\begin{equation}\label{ms:fourier:cut49:shortloc:Jfinal}
 \frac J\lambda\le\frac{\omega_+Kc_vG}{z(c_vN\lambda+G^2)}.
\end{equation}

\paragraph{The fractional-part gain.}
The following implication assumes the scalar majorization and its rational
constant enclosures stated below. Their finite arithmetic acceptance is a
separate condition.
If each centered summand has support span at most $b<1/2$ and the sum has
variance $t$, then, for any real $w$ and
$0<r_*\le2\sin^2(\pi b)/b^2$,
\begin{equation}\label{ms:fourier:cut49:shortloc:harmonic}
             \E\{w-W\}\ge\frac5{23}(1-e^{-r_*t}).
\end{equation}
Indeed $1-\cos(2\pi x)\le(23/5)x$ for $0\le x\le1$.
A finite proof of this sole scalar inequality is recorded below.
For an independent copy $Z'$ of one summand, concavity of sine on
$[0,\pi b]$ gives
\[
 |\E e^{2\pi iZ}|^2
 =1-2\E\sin^2\pi(Z-Z')
 \le e^{-4\sin^2(\pi b)\operatorname{Var}(Z)/b^2}.
\]
Multiply square roots over the independent summands and use periodicity.
This argument includes zero variances and integer collisions.
Here the conditional spans are at most $A+d_0$, so this is the $b$ to use.

For a reproducible scalar proof, $x\ge10/23$ follows from $1-\cos\le2$.
On $[0,10/23]$ use the degree-ten alternating upper polynomial for
$1-\cos(2\pi x)$, with directed coefficients:
\[
 Q(x)=\sum_{j=1}^5\frac{(-1)^{j+1}(2\pi_jx)^{2j}}{(2j)!},\qquad
 \pi_j=\begin{cases}3.141592653589794,&j\text{ odd},\\
                       3.141592653589793,&j\text{ even}.\end{cases}
\]
Put
$P(x)=23/5-Q(x)/x$. On the eight intervals with left endpoint $a_j=5j/92$
and width $w=5/92$, its degree-nine Bernstein coefficients are
\[
 b_{jk}=\sum_{i=0}^k\frac{\binom ki}{\binom9i}
           \sum_{h=i}^9[x^h]P(x)\binom hi a_j^{h-i}w^i,
 \qquad 0\le j<8,\quad0\le k\le9.
\]
Positivity of these 80 rational coefficients is a finite certificate
condition; its acceptance is not asserted in this appendix. Expansion in powers of the unit-interval coordinate
proves the formula; Bernstein weights are nonnegative and sum to one.
The alternating remainder and the stated rational bounds on $\pi$ follow
from the same elementary series above. This completely specifies the
finite scalar obligation. A rational rate follows by replacing
$\sin(\pi b)$ with $x-x^3/6+x^5/120-x^7/5040$, $x=3.14159b$, on the
checked alternating-series domain.

Set $r=r_*\omega_-$, require $rU<3$, and define
\begin{equation}\label{ms:fourier:cut49:shortloc:gain}
 g=\frac5{23}\left(\frac{r^3L^2}{6}-\frac{r^4L^3}{24}
                     +\frac{r^5L^4}{120}-\frac{r^6L^5}{720}\right).
\end{equation}
The sixth-order Taylor lower bound for $1-e^{-r\lambda}$ gives
$R_f/\lambda\ge5r/23-5r^2\lambda/46+g$. The discarded higher-order
part is increasing as a function of $\lambda$ for $r\lambda<3$:
its derivative has the sign of $48-18x+(24/5)x^2-x^3>0$ on $[0,3]$.
For instance this polynomial has positive degree-three Bernstein
coefficients $48,30,132/5,51/5$ there. Thus freezing it at $L$ is valid.

\paragraph{One final expression.}
On a signed phase cell find an interval $[x_l,x_h]$ containing $H/\ell$
by its four endpoint quotients, using positive bounds for $\ell$.
Let
\[
 D=\max(|x_l-u_h|,|x_h-u_l|),\quad
 D_+=\max(0,x_h-u_l),\quad a=\max(|u_l|,|u_h|),\quad b_\chi=p_+-x_l.
\]
The same formula allows coarser upper bounds $D,D_+,b_\chi$, as used by
the central certificate. Check $D\le\Delta$ and $b_\chi+u_h<2$.
With $\omega_\pm=1\pm c_vD/N$, the curvature coefficient is
\begin{equation}\label{ms:fourier:cut49:shortloc:V}
 V_0=\frac{\omega_+}{v_*}\left[
 \sup_{u_l\le u\le u_h}\left\{
 \frac{(b_\chi+u-1)_+}{(c-u)^2}
 +\frac{(b_\chi+u)_+}{[c+\min(u,c/3)]^2}\right\}
 +\frac{\sqrt{A^2+3T_U}}{(2c-1/2)^3}\right].
\end{equation}
On a rational phase subcell, both numerators are bounded at the right
endpoint; the first denominator is smallest there, and the second is
smallest at the left endpoint. This proves the supremum bound, including
both kinks. Equations \eqref{ms:fourier:cut49:shortloc:cutoff}--\eqref{ms:fourier:cut49:shortloc:crossingmoments},
$\sqrt n m_B\le1$, and \eqref{ms:fourier:cut49:shortloc:master} now give the whole-count
witness with
\begin{align}\label{ms:fourier:cut49:shortloc:BI}
 B(\lambda)&=z+\frac{\omega_+Kc_vG}{z(c_vN\lambda+G^2)}
                      -\frac{5r}{23}+\frac{5r^2\lambda}{46}-g,\\
 I(\lambda)&=V_0+aB_1+\frac{aP_1}{\lambda}+P_*D_+-m_cQ_*.
 \nonumber
\end{align}
We used $|u-f(u)|\le a(B_1\lambda+P_1)/n$ and $\sqrt n\pi_k\le1$.
Every constant is fixed before $\lambda$ varies. Both $B$ and
$mB+I/N$ are convex on $\lambda>0$; $I$ is decreasing. Thus their two
endpoint gates $B<0$, $I>0$, $mB+I/N<-0.001$ prove the entire closed noise
cell. Thus this argument supplies the analytic implication for each of the three
whole-count estimates, conditional on its own complete cover and endpoint gates.


\paragraph{The zero-noise endpoint and the remaining composition}
\label{ms:fourier:zero-noise}
If $\lambda=0$, then $p,q>0$ and conditional centering imply
$U_0=U_1=0$ almost surely. The actual law is its Bernoulli comparison law.
The zero-noise argument in Section~\ref{ms:counts} places an attained
positive closed-CDF maximum at a binomial atom, so $u=0$ and $C=C_B$.
The external all-probability binomial theorem of
\cite[Theorem~1]{Schulz2016}, when applied with its stated hypotheses,
then supplies the strict Bernoulli bound. For positive noise, the local
conclusion requires every parameter cell to satisfy one of the displayed
strict tests. No selection of historical table rows, completeness of a
finite cover, or fresh acceptance of those tests is asserted here.

The mathematical composition therefore still requires: the actual contact
witness; already acquired common-law support, moments and signed resources;
both deleted-law concentration caps; valid normalization and positivity
of every divisor; a complete cover of all signs, phases and noise values;
and acceptance of the stated finite gates on that cover. No step replaces
an $n-2$ convolution by an $n-1$ convolution or infers a uniform statement
from a fixed-law asymptotic limit.
